\documentclass[11pt,a4paper]{article}
\usepackage[margin=25mm,headheight=14pt]{geometry}
\usepackage{fontspec}
\setmainfont{Latin Modern Roman}
\setsansfont{Latin Modern Sans}
\setmonofont{Latin Modern Mono}[Scale=0.86]
\usepackage{amsmath,amssymb}
\usepackage{unicode-math}
\setmathfont{Latin Modern Math}
\usepackage{microtype}
\usepackage{graphicx}
\usepackage{longtable,booktabs,array}
\newcounter{none}
\usepackage{calc}
\usepackage{xcolor}
\definecolor{linkblue}{HTML}{244B70}
\usepackage{fvextra}
\DefineVerbatimEnvironment{verbatim}{Verbatim}{fontsize=\small,breaklines=true,breakanywhere=true,breaksymbolleft={},breaksymbolright={},xleftmargin=5pt}
\usepackage{enumitem}
\setlist{itemsep=2pt,parsep=1pt,topsep=4pt}
\usepackage{fancyhdr}
\pagestyle{fancy}
\fancyhf{}
\fancyhead[L]{\small Cubic SU(2) lattice Yang-Mills theory}
\fancyhead[R]{\small Consolidated reader}
\fancyfoot[C]{\thepage}
\renewcommand{\headrulewidth}{0.3pt}
\usepackage{hyperref}
\hypersetup{colorlinks=true,linkcolor=linkblue,urlcolor=linkblue,pdftitle={Fourth-order
vacuum sources and sixth-order energy in cubic SU(2) lattice Yang-Mills
theory},pdfauthor={KokunoYumeto},pdfsubject={Fourth-order source and sixth-order energy with complete supplied proofs},pdfcreator={Pandoc and XeLaTeX}}
\usepackage{bookmark}
\urlstyle{same}
\setlength{\emergencystretch}{4em}
\hfuzz=0.2pt
\setlength{\parindent}{0pt}
\setlength{\parskip}{5pt plus 1pt minus 1pt}
\setcounter{secnumdepth}{2}
\setcounter{tocdepth}{2}
\providecommand{\tightlist}{\setlength{\itemsep}{0pt}\setlength{\parskip}{0pt}}
\providecommand{\pandocbounded}[1]{#1}
\let\originalsection\section
\renewcommand{\section}{\clearpage\originalsection}
\let\originaldisplaystart\[
\let\originaldisplayend\]
\renewcommand{\[}{\begingroup\small\originaldisplaystart}
\renewcommand{\]}{\originaldisplayend\endgroup}
\begin{document}
\begin{titlepage}
\vspace*{20mm}
{\sffamily\Large MATHEMATICAL READER\par}
\vspace{12mm}
{\LARGE\bfseries Fourth-order vacuum sources and sixth-order energy in
cubic SU(2) lattice Yang-Mills theory\par}
\vspace{9mm}
{\large A consolidated reader for two exact finite-lattice
calculations\par}
\vspace{15mm}
{\large KokunoYumeto\par}
\vspace{3mm}
{17 September 2026\par}
\vfill
\begin{minipage}{0.94\textwidth}
An introduction to two exact calculations in the same original finite cubic lattice coordinates, followed by all six supplied proof texts. The fourth source and sixth energy are common to both packets. Their further analytic and computational conclusions are distinguished explicitly.
\end{minipage}
\vspace{12mm}
\end{titlepage}
\pagenumbering{roman}
\tableofcontents
\clearpage
\pagenumbering{arabic}
\section{The mathematical problem}\label{the-mathematical-problem}

This reader brings together two calculations of the fourth coefficient
of the logarithmic vacuum and the sixth coefficient of the ground energy
for the same finite cubic SU(2) lattice Hamiltonian. Both calculations
retain the original links, Haar measure, coupling constants, and open
boundary. Their common coefficients are expressed in different
coordinate systems. One calculation also supplies bounds on the complete
interacting vacuum and physical spectrum, uniform in the size of the
finite box. The other supplies explicit signed derivatives of the vacuum
coefficients and an independently proved, volume-dependent analytic
remainder for the energy.

The purpose of this introduction is to identify the mathematical
objects, state the shared results, and explain precisely what each
calculation adds. The six proof texts follow in full as Appendices A-F.
The introduction is an expository guide to those proofs, not an
additional independent analytical verification of them. The associated
exact tables and executable evidence are companion material; this volume
does not print thousands of polynomial entries. No claim of global
historical priority or of a solution of the four-dimensional continuum
Yang-Mills problem is made.

\subsection{Original lattice and
Hamiltonian}\label{original-lattice-and-hamiltonian}

Fix an integer \(L\geq2\). The vertices are \(\{-L,\ldots,L\}^3\); the
links are the positively oriented nearest-neighbor edges contained in
that box. Let \(E_L\) and \(P_L\) denote the sets of links and
elementary plaquettes, and put \(M=|P_L|\). A configuration assigns
\(U_e\in\mathrm{SU}(2)\) to each positive link. Integration is with
respect to product Haar probability on \(\mathrm{SU}(2)^{E_L}\). The
physical Hilbert space consists of the functions invariant under every
vertex gauge transformation, including the transformations at boundary
vertices.

For \(p=(n;i,j)\), \(i<j\), the ordered plaquette holonomy and its
fundamental trace are

\[
\begin{aligned}
\Omega_p&=U_i(n)U_j(n+e_i)U_i(n+e_j)^{-1}U_j(n)^{-1},\\
W_p&=\operatorname{tr}(\Omega_p),\qquad S=\sum_{p\in P_L}W_p.
\end{aligned}
\]

With \(T_\alpha=-i\sigma_\alpha/2\) and the original left derivatives
\(X_{e,\alpha}\), the kinetic operator and Hamiltonian are

\[
\begin{aligned}
K&=-\sum_{e,\alpha}X_{e,\alpha}^{2},\\
H&=\kappa\bigl[K+\xi(2M-S)\bigr],\qquad
\kappa=\frac{2g^2}{a},\qquad \xi=\frac1{4g^4}.
\end{aligned}
\]

Here \(a>0\) is the lattice spacing and \(g>0\) the original coupling.
In particular, a small value of \(\xi\) corresponds to a large value of
\(g\). The scalar \(2\kappa M\xi\) is part of the original physical
energy throughout the calculations.

\subsection{The logarithmic vacuum and the coefficient
equation}\label{the-logarithmic-vacuum-and-the-coefficient-equation}

Let \(P_H f=\int f\,dU\) and \(Q_H=I-P_H\), and define

\[
\Gamma(f,h)=\sum_{e,\alpha}(X_{e,\alpha}f)(X_{e,\alpha}h),
\qquad B(f,h)=K^{-1}Q_H\Gamma(f,h).
\]

The inverse is taken on the nonconstant physical component. At finite
coefficient order the papers implement it using the actual Casimir
spectrum of the polynomial source. They do not impose a finite-spin
approximation to the full interacting Hamiltonian.

The source is the Haar-centered logarithm \(v=Q_H\log\psi\) of the
positive vacuum. Its equation and formal expansion are

\[
v=\xi v_1+B(v,v),\qquad v_1=\frac S3,\qquad
v=\sum_{n\geq1}\xi^n v_n,\qquad
v_n=\sum_{i=1}^{n-1}B(v_i,v_{n-i})\quad(n\geq2).
\]

The source is a function. Its coefficient representatives carry the
original multisets of plaquette labels as well; these labels record
where each term came from even when an assembled function or channel
vanishes. The physical normalization and energy are recovered by

\[
\begin{aligned}
c_L&=-\frac12\log\int e^{2v}\,dU,\qquad \psi_L=e^{v+c_L},\\
E_{0,L}&=\kappa\left[2M\xi-P_H\Gamma(v,v)\right].
\end{aligned}
\]

The distinction between this normalized vacuum and auxiliary polynomial
representatives is essential. The second calculation also uses the
formal eigenline representative \(u=1+\sum_{\nu\ne0}u_\nu\lambda^\nu\),
with \(P_Hu_\nu=0\), for \(K-\sum_p\lambda_pW_p\). Its relation to the
same logarithmic source is

\[
v=Q_H\log u,\qquad u=\frac{e^v}{P_He^v}.
\]

The physical diagonal is \(\lambda_p=\xi\) for every plaquette; taking
it after coefficient calculations retains every directional response and
multiplicity.

\subsection{Source identity and reading
map}\label{source-identity-and-reading-map}

The terms \emph{trace packet} and \emph{quaternion packet} below
identify two supplied collections dated 17 September 2026. They do not
denote different physical theories. The trace packet evaluates original
trace words and builds analytic bounds from their coefficient data. The
quaternion packet evaluates polynomials in chord variables obtained by
an explicit spanning-tree coordinate map, with the original derivatives
and Haar measure transported through that map.

{\def\LTcaptype{none} % do not increment counter
\begin{longtable}[]{@{}
  >{\raggedright\arraybackslash}p{(\linewidth - 2\tabcolsep) * \real{0.5000}}
  >{\raggedright\arraybackslash}p{(\linewidth - 2\tabcolsep) * \real{0.5000}}@{}}
\toprule\noalign{}
\begin{minipage}[b]{\linewidth}\raggedright
Appendix and supplied proof text
\end{minipage} & \begin{minipage}[b]{\linewidth}\raggedright
Packet and principal role
\end{minipage} \\
\midrule\noalign{}
\endhead
\bottomrule\noalign{}
\endlastfoot
A.
{\ttfamily Q\allowbreak{}U\allowbreak{}A\allowbreak{}R\allowbreak{}T\allowbreak{}I\allowbreak{}C\allowbreak{}\_\allowbreak{}S\allowbreak{}O\allowbreak{}U\allowbreak{}R\allowbreak{}C\allowbreak{}E\allowbreak{}.\allowbreak{}m\allowbreak{}d\allowbreak{}}
& Trace: complete fourth source and coefficient bounds \\
B.
{\ttfamily P\allowbreak{}H\allowbreak{}Y\allowbreak{}S\allowbreak{}I\allowbreak{}C\allowbreak{}A\allowbreak{}L\allowbreak{}\_\allowbreak{}R\allowbreak{}E\allowbreak{}T\allowbreak{}U\allowbreak{}R\allowbreak{}N\allowbreak{}.\allowbreak{}m\allowbreak{}d\allowbreak{}}
& Trace: all-order vacuum correction and physical spectral estimate \\
C.
{\ttfamily S\allowbreak{}I\allowbreak{}X\allowbreak{}T\allowbreak{}H\allowbreak{}\_\allowbreak{}E\allowbreak{}N\allowbreak{}E\allowbreak{}R\allowbreak{}G\allowbreak{}Y\allowbreak{}.\allowbreak{}m\allowbreak{}d\allowbreak{}}
& Trace: sixth energy, cube term, remainder, and plaquette observable \\
D.
{\ttfamily F\allowbreak{}O\allowbreak{}U\allowbreak{}R\allowbreak{}T\allowbreak{}H\allowbreak{}\_\allowbreak{}O\allowbreak{}R\allowbreak{}D\allowbreak{}E\allowbreak{}R\allowbreak{}\_\allowbreak{}S\allowbreak{}O\allowbreak{}U\allowbreak{}R\allowbreak{}C\allowbreak{}E\allowbreak{}.\allowbreak{}m\allowbreak{}d\allowbreak{}}
& Quaternion: complete fourth source in exact tree and quaternion
coordinates \\
E.
{\ttfamily S\allowbreak{}I\allowbreak{}X\allowbreak{}T\allowbreak{}H\allowbreak{}\_\allowbreak{}O\allowbreak{}R\allowbreak{}D\allowbreak{}E\allowbreak{}R\allowbreak{}\_\allowbreak{}E\allowbreak{}N\allowbreak{}E\allowbreak{}R\allowbreak{}G\allowbreak{}Y\allowbreak{}.\allowbreak{}m\allowbreak{}d\allowbreak{}}
& Quaternion: sixth energy and a separate finite-volume remainder \\
F.
{\ttfamily S\allowbreak{}I\allowbreak{}G\allowbreak{}N\allowbreak{}E\allowbreak{}D\allowbreak{}\_\allowbreak{}T\allowbreak{}A\allowbreak{}N\allowbreak{}G\allowbreak{}E\allowbreak{}N\allowbreak{}T\allowbreak{}\_\allowbreak{}C\allowbreak{}E\allowbreak{}R\allowbreak{}T\allowbreak{}I\allowbreak{}F\allowbreak{}I\allowbreak{}C\allowbreak{}A\allowbreak{}T\allowbreak{}E\allowbreak{}.\allowbreak{}m\allowbreak{}d\allowbreak{}}
& Quaternion: explicit signed directional coefficients and tangent
identities \\
\end{longtable}
}

Equation labels in the supplied texts remain local to their original
files. Thus ``Appendix A, Q14'' and ``Appendix D, Q15'' refer to two
presentations of the same single-face component; an unqualified Q or E
label inside an appendix belongs to that source or its named companion.
In the sixth-energy formulas the trace packet uses \((P,T,C,B)\) for
paths, common-edge triples, corners, and cubes, whereas the quaternion
packet uses \((P_3,T_e,T_v,C)\). The introduction uses the latter
notation to avoid that collision. Both use \(J\) for unordered adjacent
plaquette pairs.

The pinned predecessor is the cubic calculation at Git revision
{\ttfamily 9\allowbreak{}1\allowbreak{}4\allowbreak{}3\allowbreak{}4\allowbreak{}b\allowbreak{}6\allowbreak{}9\allowbreak{}6\allowbreak{}2\allowbreak{}0\allowbreak{}6\allowbreak{}2\allowbreak{}b\allowbreak{}d\allowbreak{}8\allowbreak{}0\allowbreak{}4\allowbreak{}3\allowbreak{}9\allowbreak{}d\allowbreak{}4\allowbreak{}c\allowbreak{}b\allowbreak{}2\allowbreak{}c\allowbreak{}a\allowbreak{}e\allowbreak{}9\allowbreak{}d\allowbreak{}2\allowbreak{}4\allowbreak{}7\allowbreak{}9\allowbreak{}2\allowbreak{}6\allowbreak{}4\allowbreak{}d\allowbreak{}d\allowbreak{}e\allowbreak{}}
in the
\href{https://github.com/KokunoYumeto/yang-mills-interacting-workbench}{Yang-Mills
interacting workbench}. The source texts identify which inherited
definitions and bounds they use. Repository status is time-dependent;
historical status statements in the appendices describe their source
packets.

The appendix proof texts and the algorithms accompanying the edition
preserve their supplied source content. Raw conversational input is
omitted from the public edition for privacy. References in the original
texts to an upload, a session, or a ZIP describe the original
provenance; they do not imply that conversational input or every
companion file is embedded in this PDF. Formatting and navigation have
been added, while the supplied formulas and proofs have been retained.
In particular, plain-text mathematical displays are printed in their
original notation so that transcription introduces no new mathematical
choices.

\section{The shared fourth source}\label{the-shared-fourth-source}

\subsection{What is being completed}\label{what-is-being-completed}

At the quadratic reference \(q_2=\xi v_1+\xi^2v_2\), the residual is

\[
\xi v_1+B(q_2,q_2)-q_2
=\xi^3v_3+\xi^4 B(v_2,v_2).
\]

The fourth coefficient of the solution is, however,

\[
\boxed{v_4=2B(v_1,v_3)+B(v_2,v_2).}
\]

Both packets calculate this complete signed sum. The mixed term
\(2B(v_1,v_3)\) is indispensable; computing the degree-four part of the
quadratic-reference residual alone would not calculate \(v_4\).

\subsection{Completeness and two
representations}\label{completeness-and-two-representations}

A nonzero logarithmic coefficient has an edge-connected plaquette
multiset. Beginning with a plaquette containing a fixed anchor link, and
adding either a repeated plaquette or an edge-adjacent plaquette, gives
respectively

\[4,\quad46,\quad612,\quad8621\]

anchored multisets at orders one through four. At order four, signed
permutations of the three axes and translations reduce the catalogue to
78 classes:

{\def\LTcaptype{none} % do not increment counter
\begin{longtable}[]{@{}lrrrrrr@{}}
\toprule\noalign{}
Partition & \(4\) & \(3+1\) & \(2+2\) & \(2+1+1\) & \(1+1+1+1\) &
Total \\
\midrule\noalign{}
\endhead
\bottomrule\noalign{}
\endlastfoot
Classes & 1 & 2 & 2 & 19 & 54 & 78 \\
\end{longtable}
}

Both packets retain explicit maps from these representatives to all
8,621 anchored multisets. Reconstruction in a finite box includes each
contained connected multiset once, with the original inverse-link
convention and multiplicities.

The trace catalogue has 743 nonzero rational trace-monomial terms. The
quaternion catalogue has 4,044 nonzero rational monomials after the
sphere relations are used. These counts concern different presentations
of the same coefficient functions, and are not a discrepancy in \(v_4\).
A trace monomial can expand into many quaternion monomials. In the
quaternion calculation, the variables are the actual chords of the
cluster graph: the four-face lateral tube of a cube has five chord
variables, which the catalogue retains.

The equality of the intended coefficients is mathematical: both
constructions solve the same Haar-centered coefficient equation in the
original variables, and the positive Casimir makes its centered solution
unique. This identifies their results once the operator maps and
residual identities established in the respective proofs are applied. A
separate exact reconciliation for this edition also expands all 743
trace terms into the quaternion coordinates and verifies equality of the
complete polynomials in all 78 classes, together with agreement of all
8,621 signed-coordinate transports. The edition's accompanying
{\ttfamily V\allowbreak{}A\allowbreak{}L\allowbreak{}I\allowbreak{}D\allowbreak{}A\allowbreak{}T\allowbreak{}I\allowbreak{}O\allowbreak{}N\allowbreak{}.\allowbreak{}m\allowbreak{}d\allowbreak{}},
in
{\ttfamily c\allowbreak{}o\allowbreak{}n\allowbreak{}s\allowbreak{}o\allowbreak{}l\allowbreak{}i\allowbreak{}d\allowbreak{}a\allowbreak{}t\allowbreak{}i\allowbreak{}o\allowbreak{}n\allowbreak{}/\allowbreak{}2\allowbreak{}0\allowbreak{}2\allowbreak{}6\allowbreak{}0\allowbreak{}9\allowbreak{}1\allowbreak{}7\allowbreak{}/\allowbreak{}},
records that reconciliation and its evidence. The two different term
counts alone neither prove nor refute that identification.

For one plaquette, the shared closed expression is

\[
\boxed{
v_4^{(p^4)}=\frac{17}{10368}\chi_1(\Omega_p)
-\frac7{51840}\chi_2(\Omega_p),
}
\]

where \(\chi_1=W_p^2-1\) and \(\chi_2=W_p^4-3W_p^2+1\). This is one
component of the full catalogue, not a substitute for its mixed-face
coefficients.

\subsection{Signed tangent data in the quaternion
packet}\label{signed-tangent-data-in-the-quaternion-packet}

For a plaquette \(p\) and a coupling multiindex \(\rho\), the
coefficient of the directional derivative of the logarithmic vacuum is

\[Z_{p,\rho}=(\rho_p+1)v_{\rho+e_p}.\]

Appendix F computes these signed coefficients through total degree three
and verifies the differentiated source equation coefficientwise. Its
evidence comprises 78 independently reconstructed logarithmic sources,
282 explicitly expanded degree-three response polynomials, and 2,156
residual identities. Of those identities, 282 are the degree-zero
equations \(KZ_{p,0}=W_p\) and 1,874 are homogeneous positive-degree
equations. Thus every residual is zero after subtracting its prescribed
forcing. This clarifies the appendix's compressed phrase ``2,156 exactly
zero polynomial tangent residuals.'' The integer prefactor is the
occurrence multiplicity of the differentiated plaquette. These are
finite formal-parameter identities; the packet does not infer from them
a bound on the full infinite linearized inverse.

\section{The shared sixth energy
coefficient}\label{the-shared-sixth-energy-coefficient}

\subsection{A common coefficient and complementary
derivations}\label{a-common-coefficient-and-complementary-derivations}

Write the ground eigenvalue of \(K-\xi S\) near zero as

\[e(\xi)=e_2\xi^2+e_4\xi^4+e_6\xi^6+\cdots,
\qquad E_{0,L}=\kappa\bigl(2M\xi+e(\xi)\bigr).\]

A link-center involution makes the nontrivial energy branch even. The
first two coefficients are

\[e_2=-\frac M3,\qquad e_4=\frac{5M}{216}-\frac{2J}{1053}.\]

The trace packet obtains \(e_6\) from the logarithmic-source equation,
eliminating \(v_5\) with its exact equation and retaining sources
through \(v_4\). It checks the result against a separate eigenvector
recurrence. The quaternion packet uses Haar-centered eigenline
coefficients \(u_1,u_2,u_3\) and proves

\[
e_6=-\langle u_3,Ku_3\rangle_H
-e_4\|u_1\|_H^2-e_2\|u_2\|_H^2.
\]

The two state-normalization terms in this identity are part of the
energy calculation. Neither method drops the additive physical scalar or
replaces the physical norm by a Euclidean norm of the displayed
coefficients.

\subsection{Geometric contributions, including the
cube}\label{geometric-contributions-including-the-cube}

Let \(P_3\) count three-plaquette paths, \(T_e\) triples sharing a
common edge, \(T_v\) triples at a cube corner, and \(C\) elementary
cubes. The common result is

\[
\boxed{\begin{aligned}
e_6={}&-\frac{289}{77760}M
+\frac{22285}{23654592}J
-\frac{4909}{118272960}P_3\\
&+\frac{244}{4312035}T_e
-\frac{212}{542997}T_v
-\frac{83}{1944}C.
\end{aligned}}
\]

In the trace packet the coefficient for the labelled multiset \(p^4q^2\)
is \(22285/47309184\). The roles \(p\) and \(q\) may be exchanged.
Summing both roles for each unordered adjacent pair gives the
coefficient \(22285/23654592\) of \(J\) above. This factor of two is a
difference in counting convention, not a disagreement between the
packets.

The last term is a three-dimensional closed-surface contribution.
Link-center parity shows that the set of odd-multiplicity plaquettes in
a nonzero scalar coefficient must have even boundary. At total order six
the only nonempty such set is the boundary of an elementary cube. With
the six faces oriented outward, Haar contraction gives

\[\left\langle\prod_{p\in\partial c}W_p\right\rangle_H=\frac{2^8}{2^{12}}=\frac1{16}.\]

For an ordering \(\pi\) of the six faces let \(A_k(\pi)\) be its first
\(k\) faces, and let \(|\partial A_k|\) count the boundary links. The
cube coefficient is

\[
-\frac1{16}\sum_{\pi\in S_6}\prod_{k=1}^{5}
\frac4{3|\partial A_k(\pi)|}
=-\frac1{16}\frac{166}{243}
=-\frac{83}{1944}.
\]

All 720 orders are included. The intermediate denominators are the
Casimirs of the retained boundary states. Appendix E also derives the
same value by pairing complementary three-face subsets.

\subsection{Open-box formula}\label{open-box-formula}

For \(m=2L\geq4\), the original boundary is included in the exact counts

\[
\begin{gathered}
M=3m^2(m+1),\qquad J=6m(3m^2-1),\qquad C=m^3,\\
T_e=12m^2(m-1),\qquad T_v=8m^3,\\
P_3=138m^3-126m^2-24m+12.
\end{gathered}
\]

They give

\[
\boxed{e_6=-\frac{211396463m^3+30959193m^2+21845782m+2336684}{4691494080}.}
\]

Consequently

\[
\frac{E_{0,L}}\kappa
=2M\xi-\frac M3\xi^2
+\left(\frac{5M}{216}-\frac{2J}{1053}\right)\xi^4
+e_6\xi^6+\text{higher even orders}.
\]

The coefficient-density limit is
\(\lim_{m\to\infty}e_6/M=-211396463/14074482240\). This statement
concerns a coefficient. A limit of the full series requires a uniform
remainder or a separate convergence argument, as supplied for a
particular observable in the trace packet.

\section{What is controlled beyond finite
order}\label{what-is-controlled-beyond-finite-order}

\subsection{The trace packet: full vacuum and physical
spectrum}\label{the-trace-packet-full-vacuum-and-physical-spectrum}

Appendix A derives bounds for the complete fourth source in three
auxiliary coefficient norms: a Casimir-weighted local norm, a total-spin
budget \(m(v_4)\), and a budget \(t(v_4)\) at the actual anchor link.
These bounds are

\[
\begin{aligned}
\|v_4\|_{\rm loc}&\leq\frac{2296826751679}{30073680},\\
m(v_4)&\leq\frac{17270702970768271}{341697152160},\\
t(v_4)&\leq\frac{110695177857394584026401}{18025447358750832000}.
\end{aligned}
\]

These are bounds on coefficient representations used to construct and
estimate the vacuum; the physical state norm remains the original
weighted Haar norm.

The strongest construction in Appendix B uses the full fourth-order
reference \(q_4=\sum_{n=1}^4\xi^n v_n\) and its entire residual. A
convergent Neumann inverse and a convergent nonlinear series construct
\(v-q_4\), including both infinite tails. The resulting positive
function is identified with the actual normalized physical ground state.
Its source estimate is then returned to the complete physical excitation
problem.

Let \(\Delta_L\) be the lowest positive excitation energy above that
ground state. The result in Appendix B, R30 is

\[\Delta_L\geq\kappa d_{[4]}(\xi),
\qquad L\geq2,\quad a>0,\quad0<\xi\leq\alpha_{[4]}.\]

Here \(d_{[4]}\) is the explicit expression in R29, and \(\alpha_{[4]}\)
is the first positive root of the rational polynomial in R25. The stated
enclosure is

\[0.018104972231644127075<\alpha_{[4]}<0.018104972231644127076.\]

The bound has \(d_{[4]}(\xi)>3/2\) on that whole closed positive
interval. In the original coupling, a sufficient rounded condition is
\(g^2\geq3.715960362535435237\). A concrete consequence is

\[\Delta_L>1.6584\,\kappa=3.3168\,g^2/a,
\qquad g^2\geq15/4.\]

This statement is uniform in the finite box size and concerns every
physical excitation. Appendix B also preserves an earlier, weaker
construction based on \(q_3\); its endpoint \(\alpha_4\) must not be
confused with the stronger \(q_4\) endpoint \(\alpha_{[4]}\) used here.

\subsection{Two energy remainders with different
domains}\label{two-energy-remainders-with-different-domains}

The trace packet gives the physical energy remainder after order six as

\[
|R_{\geq8}|\leq\frac{49\kappa|E_L|}{1200}
\frac{(60|\xi|)^8}{1-(60|\xi|)^2},\qquad0<|\xi|<\frac1{60}.
\]

Its analytic input is the inherited uniform logarithmic-source
construction on the complex circle \(|\xi|=1/60\). After division by
\(M\), the geometric factor \(|E_L|/M=(m+1)/m\) is uniformly bounded.
The \(1/60\) domain of this Cauchy remainder is stated separately from
the larger real-coupling interval of the fourth-reference spectral
result.

The quaternion packet supplies a separate proof using bounded
perturbation of the isolated free vacuum. With

\[R=\frac3{8M},\]

its dimensionless energy remainder satisfies

\[|R_8(\xi)|\leq\frac34
\frac{(|\xi|/R)^8}{1-(|\xi|/R)^2},\qquad |\xi|<R.\]

The physical remainder is \(\kappa R_8\). This radius depends on the
volume; for \(L=2\) it is \(1/640\). The second estimate is
independently derived within its packet and does not replace the uniform
analytic argument of the first.

\subsection{The plaquette observable and the continuum
boundary}\label{the-plaquette-observable-and-the-continuum-boundary}

For the trace packet, differentiating the normalized ground-state
eigenvalue yields the mean half-trace

\[P_L(\xi)=\frac1{2M}\sum_p\langle W_p\rangle_{\rho_L},
\qquad \rho_L=\psi_L^2,\]

and a uniform remainder permits the spatial-volume limit at fixed
lattice spacing. Appendix C, E12 gives

\[P_\infty(\xi)=\frac\xi3-\frac{11}{468}\xi^3
+\frac{211396463}{4691494080}\xi^5+R_{P,\infty},\]

with

\[|R_{P,\infty}|\leq\frac{49}{40}
\frac{(60\xi)^7[8-6(60\xi)^2]}{[1-(60\xi)^2]^2},
\qquad0<\xi<1/60.\]

This is a fixed-lattice-spacing observable. Along the separately
specified path \(a_n=a_0 2^{-n}\) and
\(g_n^{-2}=g_0^{-2}+\beta n\log2\), the spectral domain is
\(g_n^{-2}\leq2\sqrt{\alpha_{[4]}}\). For \(\beta>0\) the path
eventually leaves that domain. The finite-regulator estimates therefore
do not supply a continuum field or a finite positive continuum mass.

\section{Literature and evidentiary
scope}\label{literature-and-evidentiary-scope}

The exponential-vacuum and gauge-invariant character/Casimir framework
has established antecedents. The supplied trace packet explicitly
transports the coupling, energy, derivative, and observable conventions
of its comparison sources. The quaternion packet checks the
one-plaquette coefficient against the Mathieu characteristic value in
NIST DLMF. These comparisons support attribution and stated convention
checks; they do not provide an independent published verification of the
full fourth catalogue or sixth open-box coefficient.

The claims in the appendices combine written arguments with exact finite
computation. A finite evaluation at rational link assignments is
distinguished there from a polynomial identity, and a coefficient
calculation is distinguished from a convergence or spectral theorem.
Neither packet claims independent external analytical review or a new
Lean certificate. This reader preserves those distinctions.

\subsection{References}\label{references}

The identities and versions below were checked against primary records
on 17 September 2026. This bibliographic check does not constitute a new
audit of the mathematical implications drawn from them; the scope of
each use is stated in the supplied proofs.

\begin{enumerate}
\def\labelenumi{\arabic{enumi}.}
\tightlist
\item
  D. Schütte, Zheng Weihong, and C. J. Hamer, \emph{The Coupled Cluster
  Method in Hamiltonian Lattice Field Theory}.
  \href{https://arxiv.org/abs/hep-lat/9603026v1}{arXiv:hep-lat/9603026v1},
  30 March 1996; \emph{Physical Review D} \textbf{55} (1997), 2974-2986,
  \href{https://doi.org/10.1103/PhysRevD.55.2974}{doi:10.1103/PhysRevD.55.2974}.
  Exponential-vacuum and character/Casimir framework; convention map in
  Appendix B, R23.
\item
  C. H. Llewellyn Smith and N. J. Watson, \emph{The Shifted Coupled
  Cluster Method: A New Approach to Hamiltonian Lattice Gauge Theories}.
  \href{https://arxiv.org/abs/hep-lat/9212025v1}{arXiv:hep-lat/9212025v1},
  18 December 1992; \emph{Physics Letters B} \textbf{302} (1993),
  463-471,
  \href{https://doi.org/10.1016/0370-2693(93)90428-K}{doi:10.1016/0370-2693(93)90428-K}.
  Shifted coupled-cluster framework; scope discussed in Appendix C, E7.
\item
  Thomas Spriggs, Eliska Greplova, Juan Carrasquilla, and Jannes Nys,
  \emph{Accurate ground states of SU(2) lattice gauge theory in 2+1D and
  3+1D}. \href{https://arxiv.org/abs/2509.12323v1}{arXiv:2509.12323v1},
  15 September 2025. Explicit coupling, energy, observable, and boundary
  comparison in Appendix C, E13-E14; the cited version is v1.
\item
  NIST Digital Library of Mathematical Functions, Chapter 28,
  \emph{Mathieu Functions and Hill's Equation} (chapter author: Gerhard
  Wolf), §28.6, \emph{Expansions for Small q},
  \href{https://dlmf.nist.gov/28.6.E5}{equation (28.6.5)}. Release 1.2.8
  of 15 September 2026; \href{https://dlmf.nist.gov/help/cite}{official
  citation information}. One-plaquette Mathieu coefficient comparison in
  Appendix E, E17.
\item
  Timo Jakobs, Marco Garofalo, Tobias Hartung, Karl Jansen, Johann
  Ostmeyer, Simone Romiti, and Carsten Urbach, \emph{Dynamics in
  hamiltonian lattice gauge theory: approaching the continuum limit with
  partitionings of SU(2)}. \emph{The European Physical Journal C}
  \textbf{85} (2025), article 1418, published 13 December 2025,
  \href{https://link.springer.com/article/10.1140/epjc/s10052-025-15120-x}{doi:10.1140/epjc/s10052-025-15120-x}.
  Single-plaquette Mathieu representation cited in Appendix E, E17.
\end{enumerate}

\clearpage
\appendix

\section{The complete connected fourth source in the original cubic
SU(2)
Hamiltonian}\label{the-complete-connected-fourth-source-in-the-original-cubic-su2-hamiltonian}

17 September 2026. Parent:
{\ttfamily 9\allowbreak{}1\allowbreak{}4\allowbreak{}3\allowbreak{}4\allowbreak{}b\allowbreak{}6\allowbreak{}9\allowbreak{}6\allowbreak{}2\allowbreak{}0\allowbreak{}6\allowbreak{}2\allowbreak{}b\allowbreak{}d\allowbreak{}8\allowbreak{}0\allowbreak{}4\allowbreak{}3\allowbreak{}9\allowbreak{}d\allowbreak{}4\allowbreak{}c\allowbreak{}b\allowbreak{}2\allowbreak{}c\allowbreak{}a\allowbreak{}e\allowbreak{}9\allowbreak{}d\allowbreak{}2\allowbreak{}4\allowbreak{}7\allowbreak{}9\allowbreak{}2\allowbreak{}6\allowbreak{}4\allowbreak{}d\allowbreak{}d\allowbreak{}e\allowbreak{}},
the unmerged PR8 contribution. This file,
{\ttfamily P\allowbreak{}H\allowbreak{}Y\allowbreak{}S\allowbreak{}I\allowbreak{}C\allowbreak{}A\allowbreak{}L\allowbreak{}\_\allowbreak{}R\allowbreak{}E\allowbreak{}T\allowbreak{}U\allowbreak{}R\allowbreak{}N\allowbreak{}.\allowbreak{}m\allowbreak{}d\allowbreak{}},
and
{\ttfamily S\allowbreak{}I\allowbreak{}X\allowbreak{}T\allowbreak{}H\allowbreak{}\_\allowbreak{}E\allowbreak{}N\allowbreak{}E\allowbreak{}R\allowbreak{}G\allowbreak{}Y\allowbreak{}.\allowbreak{}m\allowbreak{}d\allowbreak{}}
form one continuation. The exact original-coordinate tables are
generated by
{\ttfamily v\allowbreak{}e\allowbreak{}r\allowbreak{}i\allowbreak{}f\allowbreak{}y\allowbreak{}.\allowbreak{}p\allowbreak{}y\allowbreak{}};
every entry and every original anchored multiset is retained. This is a
written proof with finite exact computation, not a Lean certificate or
an independent external analytical review. No claim of historical
priority is made.

\subsection{Q1. Original operator and labelled coefficient
source}\label{q1.-original-operator-and-labelled-coefficient-source}

Vertices are
{\ttfamily \{\allowbreak{}-\allowbreak{}L\allowbreak{},\allowbreak{}.\allowbreak{}.\allowbreak{}.\allowbreak{},\allowbreak{}L\allowbreak{}\}\allowbreak{}\textasciicircum{}\allowbreak{}3\allowbreak{}},
L\textgreater=2. Edges are every contained positively oriented
nearest-neighbor link e=(n,i). Faces p=(n;i,j), i\textless j, carry the
original word

\begin{verbatim}
W_p=tr(U_i(n) U_j(n+e_i) U_i(n+e_j)^(-1) U_j(n)^(-1)).
\end{verbatim}

With T\_alpha=-i sigma\_alpha/2 and X\_e,alpha the original left
derivative, retain

\begin{verbatim}
K=-sum_(e,alpha) X_e,alpha^2,
H=kappa K+kappa xi(2M-S),
S=sum_p W_p, M=|P_L|, kappa=2g^2/a, xi=1/(4g^4).       (Q1)
\end{verbatim}

The scalar 2kappa xi M remains. The Hilbert measure is the original
product Haar probability. All vertex gauge transformations, including
boundary vertices, are imposed in the physical space. The physical
H\^{}2 and H\^{}1 domains and positive ground-state identity are as in
the pinned original
{\ttfamily f\allowbreak{}i\allowbreak{}n\allowbreak{}i\allowbreak{}t\allowbreak{}e\allowbreak{}\_\allowbreak{}b\allowbreak{}o\allowbreak{}x\allowbreak{}\_\allowbreak{}w\allowbreak{}e\allowbreak{}a\allowbreak{}k\allowbreak{}\_\allowbreak{}c\allowbreak{}o\allowbreak{}u\allowbreak{}p\allowbreak{}l\allowbreak{}i\allowbreak{}n\allowbreak{}g\allowbreak{}\_\allowbreak{}p\allowbreak{}h\allowbreak{}y\allowbreak{}s\allowbreak{}i\allowbreak{}c\allowbreak{}a\allowbreak{}l\allowbreak{}\_\allowbreak{}g\allowbreak{}a\allowbreak{}p\allowbreak{}.\allowbreak{}m\allowbreak{}d\allowbreak{}},
blob
{\ttfamily d\allowbreak{}f\allowbreak{}e\allowbreak{}e\allowbreak{}7\allowbreak{}7\allowbreak{}c\allowbreak{}0\allowbreak{}c\allowbreak{}a\allowbreak{}0\allowbreak{}d\allowbreak{}6\allowbreak{}e\allowbreak{}d\allowbreak{}8\allowbreak{}9\allowbreak{}7\allowbreak{}c\allowbreak{}4\allowbreak{}1\allowbreak{}7\allowbreak{}2\allowbreak{}f\allowbreak{}c\allowbreak{}b\allowbreak{}f\allowbreak{}e\allowbreak{}2\allowbreak{}3\allowbreak{}c\allowbreak{}d\allowbreak{}9\allowbreak{}d\allowbreak{}1\allowbreak{}0\allowbreak{}d\allowbreak{}3\allowbreak{}1\allowbreak{}3\allowbreak{}}.

Let P\_H be original Haar integration and Q\_H=I-P\_H. On nonconstant
physical Fourier blocks define

\begin{verbatim}
Gamma(f,h)=sum_(e,alpha)(X_e,alpha f)(X_e,alpha h),
B(f,h)=K^(-1) Q_H Gamma(f,h),
v_1=S/3, v_n=sum_(i=1)^(n-1)B(v_i,v_(n-i)).           (Q2)
\end{verbatim}

Every occurrence keeps its multiset of original face labels. The support
is their original union; cancellation of a coefficient does not delete
that support record. At the function level, assembly is the map

\begin{verbatim}
A: direct_sum_multisets V_multiset -> functions,
   (f_multiset) -> sum f_multiset.                   (Q3)
\end{verbatim}

Its kernel is the set of coefficient families with zero assembled
Fourier matrix at every original spin label. Neither a zero matrix nor a
trace identity is assigned an inverse outside the stated quotient.
Scalars removed by Q\_H are returned to the ground energy in R9 of the
companion file.

\subsection{Q2. A finite trace presentation with an explicit evaluation
map}\label{q2.-a-finite-trace-presentation-with-an-explicit-evaluation-map}

An oriented edge word is a tuple (s\_1 e\_1,\ldots,s\_r e\_r), s\_i=+1
or -1. Its evaluation is tr(product\_i U\_ei\^{}si), in that order. A
monomial is a product of these original trace evaluations. The
coefficient algebra is rational. The retained map
{\ttfamily e\allowbreak{}v\allowbreak{}} sends a formal rational sum of
such monomials to its actual smooth cylinder function.

The program chooses representatives only using the following proved
relations in ker(ev):

\begin{itemize}
\tightlist
\item
  adjacent U\_e U\_e\^{}(-1)=I, with tr(I)=2 kept as the scalar 2;
\item
  cyclic trace rotation;
\item
  reversal of the entire word with every sign reversed, since
  tr(V\^{}(-1))=tr(V) for V in SU(2);
\item
  tr(V\^{}r)=P\_r(tr V), where P\_0=2, P\_1=x, P\_r=x P\_(r-1)-P\_(r-2).
\end{itemize}

The last identity follows by multiplying V\^{}2-tr(V)V+I=0 by V\^{}(r-2)
and taking the trace. Thus {\ttfamily e\allowbreak{}v\allowbreak{}}
factors through the displayed trace relations. It is this factorization,
with the original multiset still attached, that relates the coefficient
representatives. No faithfulness of an unrestricted trace-word algebra
is presumed, and no physical parameter, vector norm, or Haar factor is
changed.

Here are the exact derivative formulas used by
{\ttfamily t\allowbreak{}r\allowbreak{}a\allowbreak{}c\allowbreak{}e\allowbreak{}\_\allowbreak{}a\allowbreak{}l\allowbreak{}g\allowbreak{}e\allowbreak{}b\allowbreak{}r\allowbreak{}a\allowbreak{}.\allowbreak{}p\allowbreak{}y\allowbreak{}}.
An occurrence of U\_e has sign +1 and a cut immediately before that
occurrence. An occurrence of U\_e\^{}(-1) has sign -1 and a cut
immediately after it. If w\_i is the cyclic word starting at that cut,
then

\begin{verbatim}
X_e,alpha tr(w)=sum_occ s_occ tr(T_alpha w_occ).       (Q4)
\end{verbatim}

For two trace factors A,B, the original Pauli matrices give

\begin{verbatim}
sum_alpha tr(T_alpha A)tr(T_alpha B)
   =-(1/2)tr(AB)+(1/4)tr(A)tr(B),
sum_alpha tr(T_alpha A T_alpha B)
   =-(1/2)tr(A)tr(B)+(1/4)tr(AB).                    (Q5)
\end{verbatim}

These follow by substituting
{\ttfamily s\allowbreak{}u\allowbreak{}m\allowbreak{}\_\allowbreak{}a\allowbreak{}l\allowbreak{}p\allowbreak{}h\allowbreak{}a\allowbreak{} \allowbreak{}(\allowbreak{}T\allowbreak{}\_\allowbreak{}a\allowbreak{}l\allowbreak{}p\allowbreak{}h\allowbreak{}a\allowbreak{})\allowbreak{}\_\allowbreak{}i\allowbreak{}j\allowbreak{}(\allowbreak{}T\allowbreak{}\_\allowbreak{}a\allowbreak{}l\allowbreak{}p\allowbreak{}h\allowbreak{}a\allowbreak{})\allowbreak{}\_\allowbreak{}k\allowbreak{}l\allowbreak{}=\allowbreak{}-\allowbreak{}(\allowbreak{}1\allowbreak{}/\allowbreak{}2\allowbreak{})\allowbreak{}d\allowbreak{}e\allowbreak{}l\allowbreak{}t\allowbreak{}a\allowbreak{}\_\allowbreak{}i\allowbreak{}l\allowbreak{} \allowbreak{}d\allowbreak{}e\allowbreak{}l\allowbreak{}t\allowbreak{}a\allowbreak{}\_\allowbreak{}j\allowbreak{}k\allowbreak{}+\allowbreak{}(\allowbreak{}1\allowbreak{}/\allowbreak{}4\allowbreak{})\allowbreak{}d\allowbreak{}e\allowbreak{}l\allowbreak{}t\allowbreak{}a\allowbreak{}\_\allowbreak{}i\allowbreak{}j\allowbreak{} \allowbreak{}d\allowbreak{}e\allowbreak{}l\allowbreak{}t\allowbreak{}a\allowbreak{}\_\allowbreak{}k\allowbreak{}l\allowbreak{}}.

For one word w and edge e, each occurrence contributes (3/4)tr(w) to
E\_e tr(w). Two different occurrences, with cut segments A,B and signs
s,t, contribute

\begin{verbatim}
st tr(A)tr(B) - (st/2)tr(w).                         (Q6)
\end{verbatim}

For a product of trace factors, retain the entire product rule

\begin{verbatim}
K product f_i = sum_i(Kf_i) product_(j!=i)f_j
             -2 sum_(i<j) Gamma(f_i,f_j) product_(k!=i,j)f_k.   (Q7)
\end{verbatim}

Equations Q4-Q7 prove that the rational word operators intertwine
evaluation with the original K and Gamma, on every original link
assignment. They also retain the nonzero cross terms. The separately
implemented
{\ttfamily d\allowbreak{}i\allowbreak{}f\allowbreak{}f\allowbreak{}e\allowbreak{}r\allowbreak{}e\allowbreak{}n\allowbreak{}t\allowbreak{}i\allowbreak{}a\allowbreak{}l\allowbreak{}\_\allowbreak{}a\allowbreak{}u\allowbreak{}d\allowbreak{}i\allowbreak{}t\allowbreak{}.\allowbreak{}p\allowbreak{}y\allowbreak{}}
differentiates quaternion matrix products by explicit generator
insertions and never calls these Fierz contractions.

\subsection{Q3. Original finite Casimir inverse, with its scalar
fibre}\label{q3.-original-finite-casimir-inverse-with-its-scalar-fibre}

For a multiset C of n faces, an edge occurring r\_e times has the full
spin list

\begin{verbatim}
j_e=r_e/2,r_e/2-1,...,r_e mod 2 /2.                 (Q8)
\end{verbatim}

At each original vertex, a contributing invariant tensor has integral
summed spin and max j\_e \textless= sum of the other j\_e. The necessity
follows from the maximum weight of their tensor product. The sufficiency
follows by successively adjoining a spin: the attainable spins fill the
interval from the excess of the largest spin over the others (or the
parity minimum) to their total, in unit steps. This inductively gives
spin zero exactly for the two stated tests. For the computation only
their necessary direction is needed; retaining extra channels would
still give a valid inverse interpolation.

Enumerate all assignments in Q8 satisfying those vertex tests and retain
their original Casimirs c=sum\_e j\_e(j\_e+1). Call the resulting finite
set C\_C. Every actual representation of the coefficient source lies in
that list. Define

\begin{verbatim}
q_C(t)=product_(c in C_C, c>0)(1-t/c),
f_C(t)=(1-q_C(t))/t-q_C(t) sum_(c in C_C,c>0)1/c.    (Q9)
\end{verbatim}

The quotient by t in Q9 is a polynomial because its numerator has zero
constant term. Direct substitution gives f\_C(0)=0 and f\_C(c)=1/c for
every positive listed c.~The original product-group spectral
decomposition therefore proves

\begin{verbatim}
K f_C(K) f=Q_H f                                   (Q10)
\end{verbatim}

on this whole finite coefficient space, with every repeated spin
multiplicity retained. This is a polynomial in the original derivatives,
not an inverse on an independently chosen matrix. The constant removed
by Q10 is recorded explicitly by original Haar integration.

For a face multiset C, let C=A+B range over ordered nonempty proper
submultisets, with their original multiplicities. The complete
coefficient is

\begin{verbatim}
v_C=f_C(K) sum_(A+B=C) Gamma(v_A,v_B),
v_{ {p} }=W_p/3.                                  (Q11)
\end{verbatim}

For repeated labels each distinct submultiset is counted once. For
distinct labels both ordered complementary choices occur. This agrees
exactly with Q2 on expansion of the original sums. It proves each
computed entry without introducing an eigenfunction cutoff for H: n
bounds only the degree of the one Taylor coefficient being calculated.

\subsection{Q4. Exhaustive fourth coefficient and inverse coordinate
transports}\label{q4.-exhaustive-fourth-coefficient-and-inverse-coordinate-transports}

Start at the four original faces incident on the anchor edge (0,0,0;1).
Grow a multiset by adjoining a face already present or a face sharing an
original edge with a present face. Retain all results as multisets. The
cardinalities at orders one through four are

\begin{verbatim}
4, 46, 612, 8621.                                  (Q12)
\end{verbatim}

This exhausts connected anchored multisets: choose a face containing the
anchor, take a spanning tree of the adjacency graph of the distinct
faces, grow along that tree, and insert each remaining repeated
occurrence. This gives a growth order of every such multiset. A
disconnected term is zero in Q11 because either an input source is
already zero or the two differentiated active supports are disjoint.

The fourth multiset counts and representative counts are respectively

{\def\LTcaptype{none} % do not increment counter
\begin{longtable}[]{@{}lrr@{}}
\toprule\noalign{}
Multiplicity & Original anchored multisets & Cubic representatives \\
\midrule\noalign{}
\endhead
\bottomrule\noalign{}
\endlastfoot
4 & 4 & 1 \\
3+1 & 84 & 2 \\
2+2 & 42 & 2 \\
2+1+1 & 1572 & 19 \\
1+1+1+1 & 6919 & 54 \\
Total & 8621 & 78 \\
\end{longtable}
}

The representative map retains an actual permutation pi of the three
axes, signs epsilon\_i, and an integer translation d:

\begin{verbatim}
y_i=epsilon_i x_(pi(i))-d_i,
x_(pi(i))=epsilon_i(y_i+d_i).                       (Q13)
\end{verbatim}

A positive original edge maps to the corresponding signed edge under
Q13. The link-variable map uses that same U or U\^{}(-1). A face maps to
the original face word or its whole inverse, which have equal
fundamental traces. Haar integration and each link Casimir are preserved
by this map. Applying Q4-Q7 proves that Q13 transports every coefficient
in Q11. The full inverse is retained for every one of the 8621 original
multisets in
{\ttfamily q\allowbreak{}u\allowbreak{}a\allowbreak{}r\allowbreak{}t\allowbreak{}i\allowbreak{}c\allowbreak{}\_\allowbreak{}t\allowbreak{}r\allowbreak{}a\allowbreak{}n\allowbreak{}s\allowbreak{}p\allowbreak{}o\allowbreak{}r\allowbreak{}t\allowbreak{}s\allowbreak{}.\allowbreak{}j\allowbreak{}s\allowbreak{}o\allowbreak{}n\allowbreak{}}.

No invariance of the auxiliary trace-coefficient norm under separate
link inversion is asserted. The estimates in Q6-Q8 below use bounds
valid in every original orientation.

{\ttfamily q\allowbreak{}u\allowbreak{}a\allowbreak{}r\allowbreak{}t\allowbreak{}i\allowbreak{}c\allowbreak{}\_\allowbreak{}c\allowbreak{}o\allowbreak{}e\allowbreak{}f\allowbreak{}f\allowbreak{}i\allowbreak{}c\allowbreak{}i\allowbreak{}e\allowbreak{}n\allowbreak{}t\allowbreak{}s\allowbreak{}.\allowbreak{}j\allowbreak{}s\allowbreak{}o\allowbreak{}n\allowbreak{}}
contains all 78 polynomials: 743 nonzero rational trace-monomial terms,
their original edge dictionaries, face words, multiplicities and
complete allowed Casimir lists. The exact self contribution is

\begin{verbatim}
v_{ {p,p,p,p} } = (17/10368) chi_1(Omega_p)
                    -(7/51840) chi_2(Omega_p).      (Q14)
\end{verbatim}

Here chi\_1=W\_p\^{}2-1 and chi\_2=W\_p\^{}4-3W\_p\^{}2+1. All other
contributions are exposed in the table in the same original words. The
producer regenerates the entire table; it does not merely test a few
displayed examples. Its differential audit checks the original equation
for every representative at two declared rational quaternion
assignments. Those finite evaluations accompany, rather than replace,
the general intertwining proof Q4-Q11.

\subsection{Q5. The auxiliary norm and its original physical gauge
inequality}\label{q5.-the-auxiliary-norm-and-its-original-physical-gauge-inequality}

For an original nonempty union label S and spin tuple j, write the
coefficient as Tr(A\_(S,j) pi\_j(U)). Define

\begin{verbatim}
||f||loc=max_a sum_(S contains a,j!=0)c(j)||A_(S,j)||_1,
m(f)=max_a sum_(S contains a,j!=0)(sum_e j_e)||A_(S,j)||_1,
t(f)=max_e sum_(S contains e,j!=0)j_e||A_(S,j)||_1.    (Q15)
\end{verbatim}

All these sums concern the explicitly retained coefficient source. The
physical norm stays int rho \textbar f\textbar\^{}2 and physical energy
stays kappa int rho sum\textbar Xf\textbar\^{}2.

For any active edge e, the endpoint invariant-tensor inequalities force
at least 2j\_e total spin on the other edges at its endpoints. Their
remote endpoints are distinct on the triangle-free original cubic graph.
Applying the same inequality there and counting every further edge at
most twice forces at least j\_e further total spin. Thus sum\_f
j\_f\textgreater=4j\_e. Since j(j+1)\textgreater=3j/2 on nonzero
half-integers,

\begin{verbatim}
c(j)>=6j_e, m(f)<=2||f||loc/3, t(f)<=||f||loc/6.     (Q16)
\end{verbatim}

A nonconstant physical block has c\textgreater=3. Products of
coefficient matrices have trace-norm bound \textbar\textbar A tensor
B\textbar\textbar\_1=\textbar\textbar A\textbar\textbar\_1\textbar\textbar B\textbar\textbar\_1.
Decomposing the actual tensor representation, pinching to its isotypic
blocks and tracing the multiplicity factors gives the complete actual
product coefficients with total trace norm at most that product. The
last assertion follows by trace duality against block-diagonal
contractions; all multiplicities remain. Each original spin generator
has operator norm j\_e. Summing its three indices and retaining both
possible anchored inputs proves

\begin{verbatim}
||B(f,h)||loc<=3[m(f)t(h)+m(h)t(f)]
             <=(2/3)||f||loc||h||loc.              (Q17)
\end{verbatim}

The first source has
(\textbar\textbar v1\textbar\textbar loc,m(v1),t(v1))\textless=(32,64/3,16/3).
The full second source retains (236,5834/39,137/6), from the parent Q2
expansion.

\subsection{Q6. Exact cubic channel capacities, including all mixed
channels}\label{q6.-exact-cubic-channel-capacities-including-all-mixed-channels}

This stage further evaluates the original cubic table; it changes no
coefficient. The per-cluster nuclear masses x are constrained as
follows. Every coordinate is nonnegative.

\begin{itemize}
\tightlist
\item
  Pair: x\_0+x\_1\textless=64, x\_0\textless=16.
\item
  Repeated cubic: x\_(1/2)+x\_(3/2)\textless=216,
  x\_(1/2)\textless=520/3.
\item
  Path cubic: x\_00+x\_01+x\_10+x\_11\textless=512; x\_00\textless=32;
  x\_00+x\_01\textless=128; x\_00+x\_10\textless=128.
\item
  Common-edge cubic: x\_(1/2)+x\_(3/2)\textless=512;
  x\_(1/2)\textless=256.
\item
  Corner cubic: retain 000,011,101,110,111, with total at most512,
  x\_000\textless=8, and x\_000+x\_b\textless=128 for each b in
  \{011,101,110\}. Each one-spin channel is exactly zero by the original
  central-vertex invariant condition.
\end{itemize}

The pair and path bounds are original single-link Haar contractions. For
the repeated cubic, H=(4/3)W\_p P\_0(W\_p W\_q)-(1/3)W\_q gives
\textbar\textbar H\textbar\textbar\_1\textless=(4/3)\emph{8}16+8/3=520/3,
and the full product chi\_1(Omega\_p)W\_q has norm at most27\emph{8=216.
For the common-edge cubic, its entire half-isotypic projection is
(2/3)sum W\_r P\_0(W\_pW\_q), with norm at most(2/3)}3\emph{8}16=256;
this keeps both half-spin copies. The corner single-zero projections
have norm at most8*16=128; its triple-zero projection is one quarter of
the original six-edge boundary trace, giving8. These bounds follow from
literal original Haar contractions and Q17, not from independent
assignment of channel masses.

A self cubic has the two fixed bounds 8,64 on its half- and three-half
character terms. Maximizing each actual weighted coefficient over these
finite polytopes, then summing over the 612 anchored multisets with
their actual anchor spins, gives

\begin{verbatim}
||v3||loc <= 918680/351,
m(v3) <= 336572872/208845,
t(v3) <= 225985217/1253070.                         (Q18)
\end{verbatim}

The weights are respectively c(j)\textbar a\_j\textbar, (sum
j\_e)\textbar a\_j\textbar, and j\_anchor\textbar a\_j\textbar{} for the
complete parent cubic coefficients a\_j.
{\ttfamily c\allowbreak{}h\allowbreak{}a\allowbreak{}n\allowbreak{}n\allowbreak{}e\allowbreak{}l\allowbreak{}\_\allowbreak{}b\allowbreak{}o\allowbreak{}u\allowbreak{}n\allowbreak{}d\allowbreak{}s\allowbreak{}.\allowbreak{}p\allowbreak{}y\allowbreak{}}
records every rational facet and vertex. The exact vertex calculation
solves every candidate set of active facets, retains feasible vertices
and evaluates those displayed weights. Since each polytope is bounded, a
linear functional has a maximum on that list. This is a finite complete
certificate of all three sums. The original v3 coefficients, signs and
physical state norm remain unchanged.

\subsection{Q7. Three independently valid fourth-source
estimates}\label{q7.-three-independently-valid-fourth-source-estimates}

For each original fourth multiset C, the table bounds its full
c-weighted norm by the smallest of the following separately proved
scalar upper bounds.

\textbf{Input-channel bound.} For each ordered split A+B=C in Q11, use
the actual eigencomponents (a\_mu,F\_mu) and (a\_nu,F\_nu) of the lower
sources. The c-weighted inverse cancels the output c.~Thus the bound is

\begin{verbatim}
3 sum_(mu,nu)|a_mu a_nu| (sum_e j_mu,e j_nu,e) x_mu x_nu.    (Q19)
\end{verbatim}

Maximize it over the two actual mass polytopes of Q6. A bilinear
functional on two compact polytopes reaches a maximum at a pair of
vertices: fix one variable, maximize the other linearly, then maximize
the retained first coordinate linearly. The producer evaluates every
such pair, including equal lower multisets without presuming
independence of their actual masses. This is a bound; treating their
equal masses as independent merely enlarges the admitted domain. Sum
every ordered split.

\textbf{Complete trace expression.} Compute Kv\_C by Q6-Q7 from the
exact table. A nonempty reduced closed word of length l has a tensor
trace coefficient of nuclear norm at most2\^{}(l-1). To verify it before
identifying repeated edge variables, use independent fundamental factors
for its occurrences. At a maximum and minimum of the original coordinate
height along the closed walk, the alternating two-index contractions
each have Euclidean norm sqrt(2); the other contractions are
two-dimensional identities. Separate row/column permutations give
norm2\^{}(l-t), with at least one maximum t\textgreater=1. Identifying
repeated original links is the actual tensor-representation restriction,
followed by the trace-norm-contractive isotypic decomposition used in
Q17. Thus the bound2\^{}(l-1) is valid in every orientation, also for
repeated original edges. Products of words multiply these bounds.
Applying this to every rational term of Kv\_C gives a whole-source bound
on \textbar\textbar v\_C\textbar\textbar loc; constants cause no
difficulty because the nonconstant Fourier sum is at most the full sum.
No rotational invariance of this auxiliary norm is used.

\textbf{Joint output-channel bound.} For a four-distinct-face multiset
with no edge in more than two faces, choose s\_e=0,1 on each shared
edge. Single-occurrence edges have spin1/2. Keep every assignment
admitted by the original vertex conditions. Let c\_A(s) be the original
Casimir of a submultiset A in that assignment. For singleton A set
a\_A=1/3; otherwise define

\begin{verbatim}
a_A(s)=sum_(empty!=B proper subset A)
   [(c_B(s)+c_(A\B)(s)-c_A(s))/(2c_A(s))]
    a_B(s)a_(A\B)(s).                              (Q20)
\end{verbatim}

A zero c\_A is assigned zero by the original Q\_H. At every shared edge
a subset has either zero factors, one fundamental factor with scalar
Casimir3/4, or the full pair with the total edge Casimir. Hence these
lifted subset Casimirs commute with the final edge projectors. Q10 and
the product rule prove Q20 on the full product of the original four
traces. A forbidden intermediate component can have a formally displayed
coefficient, but its projected function is literally zero; it is not
assigned a new vector.

For each set Z of shared edges, integrate precisely those original edge
variables. The projection P\_(0,Z) is product\_(e in Z)(I-E\_e/2) on
this four-face product. Its trace expression, produced by Q6-Q7,
supplies the orientation-independent norm bound b\_Z just proved.
Therefore its final channel masses satisfy

\begin{verbatim}
sum_(s with s_e=0 for all e in Z) x_s <= b_Z.         (Q21)
\end{verbatim}

Maximize sum\_s c(s)\textbar a\_C(s)\textbar x\_s over
x\_s\textgreater=0 and every constraint Q21. The file
{\ttfamily q\allowbreak{}u\allowbreak{}a\allowbreak{}r\allowbreak{}t\allowbreak{}i\allowbreak{}c\allowbreak{}\_\allowbreak{}b\allowbreak{}o\allowbreak{}u\allowbreak{}n\allowbreak{}d\allowbreak{}s\allowbreak{}.\allowbreak{}j\allowbreak{}s\allowbreak{}o\allowbreak{}n\allowbreak{}}
retains the entire constraint list, all coefficient weights, a rational
primal vector and a rational dual vector y. The checker verifies
y\textgreater=0, A\^{}T y\textgreater=the weights and equality of primal
and dual objectives. Consequently b\^{}T y is a proved upper bound
independent of any numerical optimization.
{\ttfamily e\allowbreak{}x\allowbreak{}a\allowbreak{}c\allowbreak{}t\allowbreak{}\_\allowbreak{}l\allowbreak{}p\allowbreak{}.\allowbreak{}p\allowbreak{}y\allowbreak{}}
obtains it by exact rational elimination with Bland's rule; all
certification inequalities are checked again. Forty-four of the 78
representatives admit this particular commuting-channel calculation; the
others retain Q19 and the complete trace bound.

\subsection{Q8. Complete summed bound and physical
scope}\label{q8.-complete-summed-bound-and-physical-scope}

Sum the selected per-cluster bounds over every original anchored
multiset using Q13 and the orientation-independent inequalities above.
The complete rational sum is

\begin{verbatim}
||v4||loc <= A4 := 2296826751679/30073680
                   < 76374.                        (Q22)
\end{verbatim}

The complete table includes the 78 exact source polynomials and their
8621 transports. The bound itself includes every ordered input split,
every admitted joint output channel and every original scalar return. It
is stronger than the preceding unevaluated estimate
{\ttfamily (\allowbreak{}1\allowbreak{}2\allowbreak{}8\allowbreak{}/\allowbreak{}3\allowbreak{})\allowbreak{}*\allowbreak{}(\allowbreak{}9\allowbreak{}4\allowbreak{}4\allowbreak{}9\allowbreak{}8\allowbreak{}4\allowbreak{}/\allowbreak{}3\allowbreak{}5\allowbreak{}1\allowbreak{})\allowbreak{}+\allowbreak{}7\allowbreak{}9\allowbreak{}9\allowbreak{}2\allowbreak{}5\allowbreak{}8\allowbreak{}/\allowbreak{}3\allowbreak{}9\allowbreak{}}.

The three quantities in Q18 and the full fourth coefficient Q22 are the
inputs used in
{\ttfamily P\allowbreak{}H\allowbreak{}Y\allowbreak{}S\allowbreak{}I\allowbreak{}C\allowbreak{}A\allowbreak{}L\allowbreak{}\_\allowbreak{}R\allowbreak{}E\allowbreak{}T\allowbreak{}U\allowbreak{}R\allowbreak{}N\allowbreak{}.\allowbreak{}m\allowbreak{}d\allowbreak{}}.
That argument controls every remaining order. The finite
fourth-coefficient calculation is never substituted for the full
interacting Hamiltonian.

\subsection{Q9. The fourth source's total-spin and actual-anchor-spin
bounds}\label{q9.-the-fourth-sources-total-spin-and-actual-anchor-spin-bounds}

The complete fourth coefficient also gives more information than its
c-weighted norm. Retain each original coefficient family and the actual
list of allowed spins from Q8. For a representative C with a proved
whole c-weighted bound b\_C, two immediate estimates are

\begin{verbatim}
m_C <= b_C max_(j in C_C,j!=0) (sum_e j_e)/c(j),
t_C,e <= b_C max_(j in C_C,j!=0) j_e/c(j).           (Q23)
\end{verbatim}

These follow term by term on the original Fourier coefficients. The
maximum is over the complete displayed assignments, not just their
distinct Casimir values. It is finite and rational. Every extra admitted
spin enlarges these upper bounds.

There is also an estimate directly from the complete v\_C trace
polynomial. For a monomial m let r\_e(m) be its full occurrence count on
edge e. Tensor decomposition places every actual output spin at most
r\_e(m)/2. With b\_m=2\^{}(sum\_words(length(word)-1)), the
orientation-independent bound from Q7 gives

\begin{verbatim}
m_C <= sum_m |a_m| b_m (sum_e r_e(m))/2,
t_C,e <= sum_m |a_m| b_m r_e(m)/2.                  (Q24)
\end{verbatim}

Every term and its original occurrences remain. These are bounds on the
same canonical Fourier coefficient quantities as Q23; taking the smaller
of the two scalar upper bounds does not change a coefficient or physical
metric.

For the 44 commuting-channel cases in Q20-Q21, retain the identical mass
constraints A x\textless=b and x\textgreater=0. Use instead the two
complete objective rows

\begin{verbatim}
weight_m(s)=(sum_e j_e(s)) |a_C(s)|,
weight_e(s)=j_e(s) |a_C(s)|.                        (Q25)
\end{verbatim}

Exact rational dual vectors y\textgreater=0 with A\^{}T
y\textgreater=weight bound these quantities by b\^{}T y. The producer
computes and verifies the primal and dual equations as for Q21. The
table
{\ttfamily q\allowbreak{}u\allowbreak{}a\allowbreak{}r\allowbreak{}t\allowbreak{}i\allowbreak{}c\allowbreak{}\_\allowbreak{}s\allowbreak{}p\allowbreak{}i\allowbreak{}n\allowbreak{}\_\allowbreak{}b\allowbreak{}u\allowbreak{}d\allowbreak{}g\allowbreak{}e\allowbreak{}t\allowbreak{}s\allowbreak{}.\allowbreak{}j\allowbreak{}s\allowbreak{}o\allowbreak{}n\allowbreak{}}
keeps every dual, the selected scalar bounds and the original edge
multiplicities.

For the point-spin sum the anchor must be transported as well. In each
actual map Q13 its endpoints 0 and e\_1 go to -d and (epsilon\_i
delta\_(pi(i),1)-d\_i)\_i. The code locates that exact original edge in
the representative dictionary and adds its own t\_C,e bound. The
total-spin bound is added once for every original anchored multiset.
Thus every original orientation and anchor are retained. In particular
this computation uses universal orientation bounds, not a claimed
nuclear-norm invariance under a partial transpose.

The resulting complete rational sums are

\begin{verbatim}
m(v4) <= m4 := 17270702970768271/341697152160,
t(v4) <= t4 := 110695177857394584026401/18025447358750832000. (Q26)
\end{verbatim}

Their approximate sizes are 50543.88912 and 6141.050242. The exact
fractions in Q26, rather than those renderings, are used in every
subsequent bound. The point-spin result is strictly smaller than A4/6.
These are proved bounds for the entire original fourth source, not
selected channels or a sampled cluster.
{\ttfamily P\allowbreak{}H\allowbreak{}Y\allowbreak{}S\allowbreak{}I\allowbreak{}C\allowbreak{}A\allowbreak{}L\allowbreak{}\_\allowbreak{}R\allowbreak{}E\allowbreak{}T\allowbreak{}U\allowbreak{}R\allowbreak{}N\allowbreak{}.\allowbreak{}m\allowbreak{}d\allowbreak{}}
R24-R34 uses them to complete the full correction at q4; R1-R23 retain
the earlier cubic-reference calculation as an intermediate result.

\subsection{Sources and attribution}\label{sources-and-attribution}

The pinned parent provides the complete second and third coordinate
calculations, original support reconstruction, and auxiliary coefficient
norms; its cubic producer was replayed unchanged. Classical
exponential-vacuum and gauge-invariant character/Casimir approaches are
described by Schuette, Zheng Weihong and Hamer, arXiv:hep-lat/9603026v1,
especially equations17,21-26,30,38; the actual physical parameter/energy
dictionary is retained in the companion return file. No new general
coupled-cluster method or worldwide-priority claim is asserted. Source
scopes and the latest peer README discovery are in SOURCE\_INTAKE.json.

\section{Full corrections at the cubic and fourth sources in the
original physical
energy}\label{full-corrections-at-the-cubic-and-fourth-sources-in-the-original-physical-energy}

17 September 2026. All operators, parameters, coefficient spaces and
source maps are those of
{\ttfamily Q\allowbreak{}U\allowbreak{}A\allowbreak{}R\allowbreak{}T\allowbreak{}I\allowbreak{}C\allowbreak{}\_\allowbreak{}S\allowbreak{}O\allowbreak{}U\allowbreak{}R\allowbreak{}C\allowbreak{}E\allowbreak{}.\allowbreak{}m\allowbreak{}d\allowbreak{}}
(Q1-Q26). R24-R34 give the strongest result, using the complete fourth
source; R1-R23 preserve the cubic-reference calculation and its explicit
maps. This argument constructs the full correction rather than
truncating the Hamiltonian. The final estimate is for the original
finite-regulator physical spectrum, uniformly in L. The smooth
four-dimensional continuum mass is not determined by this contribution.

\subsection{R1. The actual cubic reference and its whole
residual}\label{r1.-the-actual-cubic-reference-and-its-whole-residual}

Write x=xi without changing its physical value, and set

\begin{verbatim}
q3=x v1+x^2 v2+x^3 v3,
R3=x v1+B(q3,q3)-q3
  =x^4 v4+2x^5 B(v2,v3)+x^6 B(v3,v3),
J3 h=2B(q3,h), L3=I-J3.                              (R1)
\end{verbatim}

Every term is an original function. Q11 evaluates all of v4. The fifth
and sixth terms remain as the displayed complete bilinear functions,
with the independent bounds below. The Frechet factor two follows by
expansion of B(q3+h,q3+h) using symmetry of the original Gamma.

Put

\begin{verbatim}
A4=2296826751679/30073680,
m3=336572872/208845, t3=225985217/1253070,
m2=5834/39, t2=137/6,
l3=m3+4t3=292337810/125307,
b23=3(m2*t3+m3*t2)=222621900791/1163565,
b33=6m3*t3=76060493515233224/43616234025.              (R2)
\end{verbatim}

The exact Q17 source estimate and the three evaluated lower coefficients
give

\begin{verbatim}
||J3|| <= ell(x)=(128/3)x+(3132/13)x^2+l3*x^3,
||R3||loc <= delta(x)=A4*x^4+2b23*x^5+b33*x^6.        (R3)
\end{verbatim}

For the cubic part of J3, Q17 and Q16 give
{\ttfamily 2\allowbreak{}*\allowbreak{}3\allowbreak{}[\allowbreak{}m\allowbreak{}3\allowbreak{}/\allowbreak{}6\allowbreak{}+\allowbreak{}(\allowbreak{}2\allowbreak{}/\allowbreak{}3\allowbreak{})\allowbreak{}t\allowbreak{}3\allowbreak{}]\allowbreak{}=\allowbreak{}m\allowbreak{}3\allowbreak{}+\allowbreak{}4\allowbreak{}t\allowbreak{}3\allowbreak{}};
this retains the two different spin budgets. The fifth residual
coefficient is twice b23, not b23. All sixth-order residual terms are
included in b33.

Define the explicitly specified rational polynomial

\begin{verbatim}
D4(x)=(1-ell(x))^2-(8/3)delta(x).                    (R4)
\end{verbatim}

Let alpha4 be its unique root between 178/10000 and 179/10000. Exact
rational evaluation proves opposite signs at these endpoints and
ell(179/10000)\textless1. On {[}0,179/10000{]}, ell and delta are
increasing, and

\begin{verbatim}
D4'(x)=-2(1-ell(x))*ell'(x)-(8/3)delta'(x)<0.
\end{verbatim}

Thus alpha4 is also the first positive root; D4\textgreater=0 and
ell\textless1 on {[}0,alpha4{]}. The full constants in R2-R4 define the
polynomial without rounded coefficients.

\subsection{R2. Both infinite constructions and their actual
residuals}\label{r2.-both-infinite-constructions-and-their-actual-residuals}

The actual source operator L3 has the norm-convergent inverse

\begin{verbatim}
L3^(-1)=sum_(j>=0)J3^j,
||L3^(-1)||<=1/(1-ell(x)).                           (R5)
\end{verbatim}

Let z=L3\^{}(-1)R3 and C=L3\^{}(-1)B. Define w\_1=z and
{\ttfamily w\allowbreak{}\_\allowbreak{}n\allowbreak{}=\allowbreak{}s\allowbreak{}u\allowbreak{}m\allowbreak{}\_\allowbreak{}(\allowbreak{}i\allowbreak{}=\allowbreak{}1\allowbreak{})\allowbreak{}\textasciicircum{}\allowbreak{}(\allowbreak{}n\allowbreak{}-\allowbreak{}1\allowbreak{})\allowbreak{} \allowbreak{}C\allowbreak{}(\allowbreak{}w\allowbreak{}\_\allowbreak{}i\allowbreak{},\allowbreak{}w\allowbreak{}\_\allowbreak{}(\allowbreak{}n\allowbreak{}-\allowbreak{}i\allowbreak{})\allowbreak{})\allowbreak{}}.
From R3 and Q17,

\begin{verbatim}
||w_n||loc <= Cat_(n-1) c_*^(n-1) z_*^n,
z_*=delta/(1-ell), c_*=(2/3)/(1-ell),
theta=4 c_* z_*<=1.
\end{verbatim}

The complete series w=sum\_n w\_n is absolutely convergent, including
x=alpha4, and satisfies L3w=R3+B(w,w). Its bound is

\begin{verbatim}
||w||loc<=w_*(x)=(3/4)(1-ell(x)-sqrt(D4(x))).         (R6)
\end{verbatim}

At theta=1 the convergence comes from the Catalan series itself. With
b\_N=binom(2N,N)/4\^{}N, the exact identity
Cat\_N/4\^{}N=2(b\_N-b\_(N+1)) gives the full nonlinear tail

\begin{verbatim}
||sum_(n>N)w_n||loc
   <=(3/4)(1-ell)theta^(N+1)b_N
   <=(3/4)(1-ell)theta^(N+1)/sqrt(N+1).              (R7)
\end{verbatim}

The final inequality follows from b\_N\^{}2\textless=1/(N+1), proved by
induction from the displayed ratio b\_(N+1)/b\_N=(2N+1)/(2N+2).
Independently, the finite linear sum has the exact original residual

\begin{verbatim}
L3^(-1)R3-sum_(j=0)^N J3^j R3
   =J3^(N+1)L3^(-1)R3.                             (R8)
\end{verbatim}

Its norm is at most ell\^{}(N+1)delta/(1-ell). Each J3 operation adjoins
at most three original faces to the union label, and R3 has at most six;
the finite linear sums therefore retain supports of at most6+3N faces.
R7 retains every subsequent nonlinear union.

\subsection{R3. Identification with the original positive vacuum and its
scalar}\label{r3.-identification-with-the-original-positive-vacuum-and-its-scalar}

At each finite L the labelled normed coefficient source is complete, and
its global c-weighted sum is at most \textbar E\_L\textbar{} times its
local norm. The original spin-generator bounds control every first and
second derivative by that global sum. Thus the convergent v=q3+w is
C\^{}2 on the original compact product group. The source equation can be
summed with its full derivatives. Retain

\begin{verbatim}
C(v)=int_H Gamma(v,v),
c_L=-(1/2)log int_H exp(2v),
psi_L=exp(v+c_L),
E0,L=2kappa x M-kappa C(v).                         (R9)
\end{verbatim}

The original product rule gives H psi\_L=E0,L psi\_L. Elliptic
regularity makes psi\_L smooth and strictly positive. Multiplication and
division by psi\_L preserve the original finite-volume H\^{}1 domain.
Integration by parts gives, for every smooth f,

\begin{verbatim}
q_(H-E0,L)(psi_L f)=kappa int psi_L^2 sum_i|X_i f|^2. (R10)
\end{verbatim}

This proves E0,L is the actual ground energy. A second ground state
divided by psi\_L has every original derivative zero and is constant.
The source construction is gauge invariant, so this is the actual
physical vacuum. The argument does not assume an excitation gap.

The approximate exponential has a separate exact identity, which retains
its multiplication defect:

\begin{verbatim}
exp(-q3) H exp(q3)
  =kappa[K-2Gamma(q3,.)]
   +kappa[2Mx-C(q3)]I-kappa M_(K R3).                (R11)
\end{verbatim}

Here M\_(K R3) means multiplication by the full original polynomial K
R3. Q1-Q2 and a direct differentiation prove R11. Its scalar C(q3) and
the multiplication operator remain; they are not replaced by the actual
E0,L before R9 is proved.

The same absolute coefficient arguments give analyticity for complex
\textbar x\textbar\textless alpha4. At every original power the unique
coefficient recurrence is Q2. Consequently this solution agrees, by its
actual coefficients and the ground-state identification, with the
preceding constructions on their common domain. No competing vacuum or
comparison measure is substituted.

\subsection{R4. Return to the complete physical spectral
problem}\label{r4.-return-to-the-complete-physical-spectral-problem}

The original transported operator is

\begin{verbatim}
A=psi_L^(-1)(H-E0,L)psi_L=kappa[K-2Gamma(v,.)].      (R12)
\end{verbatim}

The complete actual point-spin budget is

\begin{verbatim}
t(v)<= (16/3)x+(137/6)x^2+t3*x^3+w_*/6.
\end{verbatim}

For a physical Fourier function h, the three original generator indices
and sum\_e j\_e\textless=2c(j)/3 give

\begin{verbatim}
||2Gamma(v,h)||_X<=chi(x)||Kh||_X,
chi(x)=4[(16/3)x+(137/6)x^2+t3*x^3+w_*/6].           (R13)
\end{verbatim}

The norm X is the sum of the original trace norms of Fourier
coefficients. The scalar output is included in this bound. Substitution
of R6 gives the exact identity

\begin{verbatim}
3(1-chi(x))=d4(x)
  =(3/2)(1+sqrt(D4(x)))+(1136/13)x^2+p3*x^3,
p3=(3/2)m3-6t3=278874091/208845>0.                 (R14)
\end{verbatim}

In particular chi\textless1 throughout the entire closed domain. At
alpha4, d4 exceeds3/2.

Let lambda\textgreater0 be an actual physical excitation eigenvalue, and
h the zero-Haar-mean part of its eigenfunction after division by the
actual vacuum. A constant transported eigenfunction has eigenvalue zero,
so h is nonzero. The original eigen-equation gives

\begin{verbatim}
(K-lambda/kappa)h=Q_H 2Gamma(v,h).                  (R15)
\end{verbatim}

Smoothness on this fixed finite product implies absolute summability of
its c-weighted Fourier trace coefficients: applying arbitrarily high
powers of the elliptic K gives rapid decay, while the number and
dimensions of product-spin representations grow polynomially in c at
this fixed L. Cauchy-Schwarz then gives the required trace-norm
summability. Hence R13 applies to the actual eigenfunction. Q16 gives
c\textgreater=3 on every nonconstant physical block. For
0\textless lambda/kappa\textless3,

\begin{verbatim}
||(K-lambda/kappa)h||_X
   >=(1-lambda/(3kappa))||Kh||_X.
\end{verbatim}

Combining this with R13 proves lambda\textgreater=kappa d4(x).
Excitations with lambda/kappa\textgreater=3 obey the same bound because
d4(x)=3(1-chi(x))\textless=3. The compact physical spectral resolution
returns this to the full original physical form domain:

\begin{verbatim}
Delta_L>=kappa d4(x),
q_A(f)>=kappa d4(x) Var_(rho_L)(f),
L>=2, a>0, 0<x<=alpha4.                             (R16)
\end{verbatim}

Thus the result includes states whose original labels depend on L. It is
a lower bound on all physical excitations, not a lower bound only on a
chosen finite trace family.

\subsection{R5. Exact numerical enclosures in original physical
units}\label{r5.-exact-numerical-enclosures-in-original-physical-units}

Rational bisection and integer-square comparisons give

\begin{verbatim}
0.01781868048538520630 < alpha4 < 0.01781868048538520631,
3.745693472404566975 < 1/(2sqrt(alpha4))
                    < 3.745693472404566976.          (R17)
\end{verbatim}

The exact condition on the original coupling is
g\^{}2\textgreater=1/(2sqrt(alpha4)). A sufficient rounded condition is
g\^{}2\textgreater=3.745693472404566976. The parent threshold was
about3.825973052393386; the comparison retains that stronger saved
value, not an earlier weaker summary.

At g\^{}2=15/4, x=4/225, the bound is

\begin{verbatim}
d4(4/225)>1.5794.
\end{verbatim}

The same rational bound holds for all g\^{}2\textgreater=15/4. To verify
the required monotonicity, for 0\textless=x\textless=4/225 use
D4\textless=1, ell'\textgreater=128/3 and R4 to bound

\begin{verbatim}
d4'(x)<=-(3/2)(1-ell(4/225))(128/3)
         +2(1136/13)(4/225)+3p3(4/225)^2<0.
\end{verbatim}

Every number in this final test is rational. Therefore

\begin{verbatim}
Delta_L>1.5794*kappa=3.1588*g^2/a,
g^2>=15/4.                                        (R18)
\end{verbatim}

At g\^{}2=4, x=1/64,

\begin{verbatim}
d4(1/64)>1.88578,
||v-q3||loc<0.019533267617373.                      (R19)
\end{verbatim}

The unchanged physical scale is kappa=8/a there. These bounds include
the complete linear and nonlinear tails in R7-R8. The generated
verification record contains exact rational brackets, not only these
decimal renderings.

\subsection{R6. Retained source kernels, complete operator defect, and
physical
primitive}\label{r6.-retained-source-kernels-complete-operator-defect-and-physical-primitive}

For the actual finite windows

\begin{verbatim}
V_N=span{R3,J3R3,...,J3^N R3},
V_N --L3--> V --0--> 0,
\end{verbatim}

the inclusion in degree zero and the identity in degree one give the
commuting support transitions. R5 supplies the explicit isomorphism

\begin{verbatim}
V_(N+1)/V_N -> ker[V/L3 V_N -> V/L3 V_(N+1)],
[h] -> [L3h],   inverse [L3h] -> [h].               (R20)
\end{verbatim}

Changing a representative by L3 V\_N changes the inverse by exactly
V\_N, proving both inverse laws. Subtraction of the actual finite
primitive sum leaves J3\^{}(N+1)R3 as the representative. The source
label remains attached at a receiving zero. This is the actual Split
Zero support-indexed complex and its retained primitive, with the same
original coefficient domain throughout.

The map from this source to the physical equation has the exact
identities

\begin{verbatim}
K L3 h=Q_H[K-2Gamma(q3,.)]h,
Q_H A h-kappa K L3h=-2kappa Q_H Gamma(w,h).           (R21)
\end{verbatim}

The right side has X norm at
most(2kappa/3)w\_*\textbar\textbar Kh\textbar\textbar\_X by the original
spin bounds. R11 separately retains the entire multiplication residual
and scalar of the reference exponential. Thus no identity between the
linearized source and the full physical operator is asserted without its
displayed defect.

On centered original physical functions let d\_X f=(X\_i f)\_i, with
squared derivative norm kappa int rho\_L sum\_i \textbar X\_i
f\textbar\^{}2. The inverse on its actual range is p\_X(d\_X f)=f.~Its
uniqueness follows from the derivative-zero functions being constants
and the centering condition. The original ground-state map
f-\textgreater psi\_L f and its inverse h-\textgreater h/psi\_L preserve
the specified pairings. The variational characterization and R16 give

\begin{verbatim}
||p_X||^2=1/Delta_L<=1/(kappa d4(x)).                 (R22)
\end{verbatim}

An actual gauge-invariant conditional-expectation kernel is a centered
subspace. Restriction of R16 to its original closed form consequently
gives the same lower bound for its physical restricted operator. This
returns the stronger inverse estimate to the previously computed
zero-shift response without changing its forcing or Gram matrix.

\subsection{R7. Continuum scope and literature
correspondence}\label{r7.-continuum-scope-and-literature-correspondence}

The standing path is a\_n=a0\emph{2\^{}(-n), g\_n\^{}2=1/c\_n,
c\_n=g0\^{}(-2)+beta}n*log2. Its inclusion in R16 is exactly
c\_n\textless=2sqrt(alpha4). For beta\textgreater0 the path eventually
leaves this proved domain. This finite-regulator, volume-uniform result
does not assign a finite continuum mass or a nontrivial ultraviolet
field to that path.

The exponential vacuum and linear excitation equations have established
antecedents. In Schuette, Zheng Weihong and Hamer,
arXiv:hep-lat/9603026v1, the coefficient x\_C=2/g\_C\^{}4 and physical
operator H\_C=g\_C\^{}2{[}K-x\_C S{]}/(2a) have the exact
original-coordinate return

\begin{verbatim}
g_C=2^(3/4)g,
H_original=sqrt(2)H_C(a,g_C)+2kappa xi M I.           (R23)
\end{verbatim}

Substitution verifies both coefficients and the additive energy. Their
Hermitian-generator commutator derivative corresponds to iX here, giving
the original signs K v-Gamma(v,v)-xi S. The exact physical Haar/energy
pairings and complete source residuals in R1-R22 are kept after this
comparison. The general method and existence of higher-order
calculations are not claimed as new. No global priority or
best-known-threshold statement follows from this scoped literature
reading.

\subsection{R8. Execute the full correction at the calculated fourth
source}\label{r8.-execute-the-full-correction-at-the-calculated-fourth-source}

The fourth table is now used as an actual reference, rather than left as
a remainder estimate. Put

\begin{verbatim}
q4=x v1+x^2 v2+x^3 v3+x^4 v4,
J_[4]h=2B(q4,h), L_[4]=I-J_[4].
\end{verbatim}

Expansion of the unchanged source equation Q2 gives the entire residual,
with all ordered contributions:

\begin{verbatim}
R_[4]=x v1+B(q4,q4)-q4
  =x^5[2B(v1,v4)+2B(v2,v3)]
   +x^6[2B(v2,v4)+B(v3,v3)]
   +2x^7 B(v3,v4)+x^8 B(v4,v4).                   (R24)
\end{verbatim}

Each displayed B is the original function with its full union support.
The fifth source is not claimed to have been separately expanded into
its final channel table. R24 is an exact expression for every part of
the reference residual, including the sixth, seventh and eighth degrees.

Use the exact m4,t4 from Q26, together with the earlier m\_i,t\_i.
Define

\begin{verbatim}
l4=m4+4t4,
b14=3(m1*t4+m4*t1), b24=3(m2*t4+m4*t2),
b34=3(m3*t4+m4*t3), b44=6m4*t4,
ell_[4](x)=ell(x)+l4*x^4,
delta_[4](x)=(2b14+2b23)x^5+(2b24+b33)x^6
                   +2b34*x^7+b44*x^8,
D_[4](x)=(1-ell_[4](x))^2-(8/3)delta_[4](x).        (R25)
\end{verbatim}

All of these constants are rational and positive. The underlying
estimates are exactly Q17:

\begin{verbatim}
||J_[4]||<=ell_[4], ||R_[4]||loc<=delta_[4].
\end{verbatim}

In particular the fourth contribution to the derivative is
{\ttfamily 2\allowbreak{}*\allowbreak{}3\allowbreak{}[\allowbreak{}m\allowbreak{}4\allowbreak{}/\allowbreak{}6\allowbreak{}+\allowbreak{}(\allowbreak{}2\allowbreak{}/\allowbreak{}3\allowbreak{})\allowbreak{}t\allowbreak{}4\allowbreak{}]\allowbreak{}=\allowbreak{}m\allowbreak{}4\allowbreak{}+\allowbreak{}4\allowbreak{}t\allowbreak{}4\allowbreak{}},
and the factor two in each unequal reference pair is retained.

Let alpha\_{[}4{]} be the first positive root of the polynomial
D\_{[}4{]}. On {[}0,182/10000{]}, ell\_{[}4{]}\textless1, ell\_{[}4{]}
and delta\_{[}4{]} are increasing, and

\begin{verbatim}
D_[4]'=-2(1-ell_[4])ell_[4]'-(8/3)delta_[4]'<0.
\end{verbatim}

Exact signs at 181/10000 and 182/10000 therefore prove existence and
uniqueness of this root in that interval, and its first-positive-root
status. Integer-square comparisons and rational bisection give

\begin{verbatim}
0.018104972231644127075 < alpha_[4]
                       < 0.018104972231644127076,
3.715960362535435236 < 1/(2sqrt(alpha_[4]))
                    < 3.715960362535435237.         (R26)
\end{verbatim}

For every 0\textless x\textless=alpha\_{[}4{]}, L\_{[}4{]} has its
actual norm-convergent Neumann inverse and the complete correction
w\footnote{4}=v-q4 is constructed by the same explicit binary-tree
series as R5-R7, with the R25 data. Thus

\begin{verbatim}
||L_[4]^(-1)||<=1/(1-ell_[4]),
||w^[4]||loc <= w_[4]*(x)
   =(3/4)(1-ell_[4](x)-sqrt(D_[4](x))).             (R27)
\end{verbatim}

For clarity, set z=L\_{[}4{]}\^{}(-1)R\_{[}4{]}, C=L\_{[}4{]}\^{}(-1)B,
w\_1=z and
{\ttfamily w\allowbreak{}\_\allowbreak{}n\allowbreak{}=\allowbreak{}s\allowbreak{}u\allowbreak{}m\allowbreak{}\_\allowbreak{}(\allowbreak{}i\allowbreak{}=\allowbreak{}1\allowbreak{})\allowbreak{}\textasciicircum{}\allowbreak{}(\allowbreak{}n\allowbreak{}-\allowbreak{}1\allowbreak{})\allowbreak{} \allowbreak{}C\allowbreak{}(\allowbreak{}w\allowbreak{}\_\allowbreak{}i\allowbreak{},\allowbreak{}w\allowbreak{}\_\allowbreak{}(\allowbreak{}n\allowbreak{}-\allowbreak{}i\allowbreak{})\allowbreak{})\allowbreak{}}.
Its scalar majorants are
{\ttfamily z\allowbreak{}\_\allowbreak{}*\allowbreak{}=\allowbreak{}d\allowbreak{}e\allowbreak{}l\allowbreak{}t\allowbreak{}a\allowbreak{}\_\allowbreak{}[\allowbreak{}4\allowbreak{}]\allowbreak{}/\allowbreak{}(\allowbreak{}1\allowbreak{}-\allowbreak{}e\allowbreak{}l\allowbreak{}l\allowbreak{}\_\allowbreak{}[\allowbreak{}4\allowbreak{}]\allowbreak{})\allowbreak{}}
and
{\ttfamily c\allowbreak{}\_\allowbreak{}*\allowbreak{}=\allowbreak{}(\allowbreak{}2\allowbreak{}/\allowbreak{}3\allowbreak{})\allowbreak{}/\allowbreak{}(\allowbreak{}1\allowbreak{}-\allowbreak{}e\allowbreak{}l\allowbreak{}l\allowbreak{}\_\allowbreak{}[\allowbreak{}4\allowbreak{}]\allowbreak{})\allowbreak{}}.
The parameter theta\_{[}4{]}=4c\_\emph{z\_} lies in {[}0,1{]}. The exact
complete tail is bounded by

\begin{verbatim}
||sum_(n>N)w_n||loc
   <=(3/4)(1-ell_[4]) theta_[4]^(N+1)
                      binom(2N,N)/4^N
   <=(3/4)(1-ell_[4]) theta_[4]^(N+1)/sqrt(N+1).    (R28)
\end{verbatim}

At alpha\_{[}4{]} this converges to zero by the central-binomial bound.
The inverse truncation separately has residual
{\ttfamily J\allowbreak{}\_\allowbreak{}[\allowbreak{}4\allowbreak{}]\allowbreak{}\textasciicircum{}\allowbreak{}(\allowbreak{}N\allowbreak{}+\allowbreak{}1\allowbreak{})\allowbreak{}L\allowbreak{}\_\allowbreak{}[\allowbreak{}4\allowbreak{}]\allowbreak{}\textasciicircum{}\allowbreak{}(\allowbreak{}-\allowbreak{}1\allowbreak{})\allowbreak{}R\allowbreak{}\_\allowbreak{}[\allowbreak{}4\allowbreak{}]\allowbreak{}}
and norm at most
{\ttfamily e\allowbreak{}l\allowbreak{}l\allowbreak{}\_\allowbreak{}[\allowbreak{}4\allowbreak{}]\allowbreak{}\textasciicircum{}\allowbreak{}(\allowbreak{}N\allowbreak{}+\allowbreak{}1\allowbreak{})\allowbreak{}d\allowbreak{}e\allowbreak{}l\allowbreak{}t\allowbreak{}a\allowbreak{}\_\allowbreak{}[\allowbreak{}4\allowbreak{}]\allowbreak{}/\allowbreak{}(\allowbreak{}1\allowbreak{}-\allowbreak{}e\allowbreak{}l\allowbreak{}l\allowbreak{}\_\allowbreak{}[\allowbreak{}4\allowbreak{}]\allowbreak{})\allowbreak{}}.
These equations control both infinities without replacing their actual
coefficient sources. A finite inverse term has original support in at
most8+4N faces, and every nonlinear union remains recorded.

The resulting v=q4+w\footnote{4} is C\^{}2 in the original finite-volume
coordinates by the same global c-summability proof in R3. Its full
equation proves R9-R10 with this v, hence identifies the same actual
positive physical vacuum and its scalar. The exact approximate
exponential formula is R11 with q3,R3 replaced by q4,R\_{[}4{]}; the
full multiplication defect K R\_{[}4{]} and the scalar C(q4) remain.

\subsection{R9. Strongest full physical lower bound and retained
inverse}\label{r9.-strongest-full-physical-lower-bound-and-retained-inverse}

The actual point-spin bound is now

\begin{verbatim}
t(v)<=t1*x+t2*x^2+t3*x^3+t4*x^4+w_[4]*/6.
\end{verbatim}

The complete physical excitation argument R12-R16 applies to this actual
v. Its returned margin is exactly

\begin{verbatim}
d_[4](x)=(3/2)(1+sqrt(D_[4](x)))
         +(1136/13)x^2+p3*x^3+p4*x^4,
p4=(3/2)m4-6t4
   =32092619324045301088/823531037954625>0.          (R29)
\end{verbatim}

Indeed expanding
{\ttfamily 3\allowbreak{}[\allowbreak{}1\allowbreak{}-\allowbreak{}4\allowbreak{}(\allowbreak{}t\allowbreak{}1\allowbreak{}*\allowbreak{}x\allowbreak{}+\allowbreak{}t\allowbreak{}2\allowbreak{}*\allowbreak{}x\allowbreak{}\textasciicircum{}\allowbreak{}2\allowbreak{}+\allowbreak{}t\allowbreak{}3\allowbreak{}*\allowbreak{}x\allowbreak{}\textasciicircum{}\allowbreak{}3\allowbreak{}+\allowbreak{}t\allowbreak{}4\allowbreak{}*\allowbreak{}x\allowbreak{}\textasciicircum{}\allowbreak{}4\allowbreak{}+\allowbreak{}w\allowbreak{}\_\allowbreak{}[\allowbreak{}4\allowbreak{}]\allowbreak{}*\allowbreak{}/\allowbreak{}6\allowbreak{})\allowbreak{}]\allowbreak{}}
using R27 proves R29 term by term. In particular
d\_{[}4{]}\textgreater3/2 on the full closed positive-coupling interval.
The exact original theorem is

\begin{verbatim}
Delta_L>=kappa d_[4](xi),
q_A(f)>=kappa d_[4](xi) Var_(rho_L)(f),
L>=2, a>0, 0<xi<=alpha_[4].                        (R30)
\end{verbatim}

The proof includes every physical form-domain vector through the
complete compact spectral resolution as in R15-R16. No preassigned
missing excitation gap is used in the source construction. A sufficient
rounded original coupling condition is
{\ttfamily g\allowbreak{}\textasciicircum{}\allowbreak{}2\allowbreak{}\textgreater{}\allowbreak{}=\allowbreak{}3\allowbreak{}.\allowbreak{}7\allowbreak{}1\allowbreak{}5\allowbreak{}9\allowbreak{}6\allowbreak{}0\allowbreak{}3\allowbreak{}6\allowbreak{}2\allowbreak{}5\allowbreak{}3\allowbreak{}5\allowbreak{}4\allowbreak{}3\allowbreak{}5\allowbreak{}2\allowbreak{}3\allowbreak{}7\allowbreak{}};
the exact condition is
{\ttfamily g\allowbreak{}\textasciicircum{}\allowbreak{}2\allowbreak{}\textgreater{}\allowbreak{}=\allowbreak{}1\allowbreak{}/\allowbreak{}(\allowbreak{}2\allowbreak{}s\allowbreak{}q\allowbreak{}r\allowbreak{}t\allowbreak{}(\allowbreak{}a\allowbreak{}l\allowbreak{}p\allowbreak{}h\allowbreak{}a\allowbreak{}\_\allowbreak{}[\allowbreak{}4\allowbreak{}]\allowbreak{})\allowbreak{})\allowbreak{}}.

At the concrete original parameters,

\begin{verbatim}
g^2=15/4, xi=4/225: d_[4](xi)>1.6584,
g^2=4, xi=1/64:     d_[4](xi)>1.89811.              (R31)
\end{verbatim}

The first bound holds throughout g\^{}2\textgreater=15/4. For x0=4/225
the exact rational inequality

\begin{verbatim}
-(3/2)(1-ell_[4](x0))(128/3)
   +2(1136/13)x0+3p3*x0^2+4p4*x0^3<0
\end{verbatim}

bounds d\_{[}4{]}' from above for 0\textless x\textless=x0, because
sqrt(D\_{[}4{]})\textless=1 and ell\_{[}4{]}' \textgreater=128/3. Thus
d\_{[}4{]} decreases on that interval. It follows in original physical
units that

\begin{verbatim}
Delta_L>1.6584*kappa=3.3168*g^2/a,
g^2>=15/4.                                        (R32)
\end{verbatim}

At g\^{}2=4 the complete reference error obeys
{\ttfamily |\allowbreak{}|\allowbreak{}v\allowbreak{}-\allowbreak{}q\allowbreak{}4\allowbreak{}|\allowbreak{}|\allowbreak{}l\allowbreak{}o\allowbreak{}c\allowbreak{}\textless{}\allowbreak{}0\allowbreak{}.\allowbreak{}0\allowbreak{}1\allowbreak{}1\allowbreak{}1\allowbreak{}6\allowbreak{}9\allowbreak{}7\allowbreak{}6\allowbreak{}8\allowbreak{}0\allowbreak{}0\allowbreak{}0\allowbreak{}3\allowbreak{}1\allowbreak{}2\allowbreak{}}.
The generated receipt retains full rational brackets for the correction,
discriminant, and original energy factors.

For the actual source windows
{\ttfamily V\allowbreak{}\_\allowbreak{}N\allowbreak{}=\allowbreak{}s\allowbreak{}p\allowbreak{}a\allowbreak{}n\allowbreak{}\{\allowbreak{}R\allowbreak{}\_\allowbreak{}[\allowbreak{}4\allowbreak{}]\allowbreak{},\allowbreak{}J\allowbreak{}\_\allowbreak{}[\allowbreak{}4\allowbreak{}]\allowbreak{}R\allowbreak{}\_\allowbreak{}[\allowbreak{}4\allowbreak{}]\allowbreak{},\allowbreak{}.\allowbreak{}.\allowbreak{}.\allowbreak{},\allowbreak{}J\allowbreak{}\_\allowbreak{}[\allowbreak{}4\allowbreak{}]\allowbreak{}\textasciicircum{}\allowbreak{}N\allowbreak{} \allowbreak{}R\allowbreak{}\_\allowbreak{}[\allowbreak{}4\allowbreak{}]\allowbreak{}\}\allowbreak{}},
the two degree maps and inverse in R20 are now the explicit maps

\begin{verbatim}
[h] -> [L_[4]h], inverse [L_[4]h] -> [h].
\end{verbatim}

They are well-defined on
{\ttfamily V\allowbreak{}\_\allowbreak{}(\allowbreak{}N\allowbreak{}+\allowbreak{}1\allowbreak{})\allowbreak{}/\allowbreak{}V\allowbreak{}\_\allowbreak{}N\allowbreak{}}
and the corresponding transported cohomology kernel because L\_{[}4{]}
is injective and its inverse is R27. Both inverse compositions remain
literal. The receiving scalar, support and full physical defect are

\begin{verbatim}
K L_[4]h=Q_H[K-2Gamma(q4,.)]h,
Q_H A h-kappa K L_[4]h=-2kappa Q_H Gamma(w^[4],h),
||p_X||^2=1/Delta_L<=1/(kappa d_[4](xi)).             (R33)
\end{verbatim}

The defect norm is at most
{\ttfamily (\allowbreak{}2\allowbreak{}k\allowbreak{}a\allowbreak{}p\allowbreak{}p\allowbreak{}a\allowbreak{}/\allowbreak{}3\allowbreak{})\allowbreak{}w\allowbreak{}\_\allowbreak{}[\allowbreak{}4\allowbreak{}]\allowbreak{}*\allowbreak{}|\allowbreak{}|\allowbreak{}K\allowbreak{}h\allowbreak{}|\allowbreak{}|\allowbreak{}\_\allowbreak{}X\allowbreak{}}.
The physical derivative norm and ground-state multiplication/division
maps are exactly those in R22. Thus the stronger inverse also restricts
to the original physical conditional kernel, without changing a forcing
vector or its raw Gram matrix.

The standing simultaneous path has the exact domain

\begin{verbatim}
a_n=a0*2^(-n), g_n^2=1/c_n,
c_n=g0^(-2)+beta*n*log2 <= 2sqrt(alpha_[4]).          (R34)
\end{verbatim}

For beta\textgreater0 it eventually leaves the certified interval. The
root in R26 is the endpoint of this particular proved majorant, not an
assigned singularity of the actual theory. This contribution gives no
ultraviolet continuum field or finite positive continuum mass. The
complete fourth table and its q4 correction are now finished inputs; the
next unevaluated coefficient is the original channel-resolved v5.

\section{The complete sixth-order ground energy and its original cube
contribution}\label{the-complete-sixth-order-ground-energy-and-its-original-cube-contribution}

17 September 2026. This calculation uses the actual fourth source
evaluated in
{\ttfamily Q\allowbreak{}U\allowbreak{}A\allowbreak{}R\allowbreak{}T\allowbreak{}I\allowbreak{}C\allowbreak{}\_\allowbreak{}S\allowbreak{}O\allowbreak{}U\allowbreak{}R\allowbreak{}C\allowbreak{}E\allowbreak{}.\allowbreak{}m\allowbreak{}d\allowbreak{}},
not a finite Hamiltonian approximation. The original scalar energy and
all extensive and boundary factors are retained. The coefficient is
reproduced by an independent eigenvector recurrence, and the cube
contribution is additionally computed by all720 original face orderings.

\subsection{E1. Eliminate the fifth source through its exact
equation}\label{e1.-eliminate-the-fifth-source-through-its-exact-equation}

Retain the original real Haar pairing and source recurrence. Since

\begin{verbatim}
C(v)=int_H Gamma(v,v), E0=2kappa xi M-kappa C(v),
\end{verbatim}

its sixth energy coefficient is

\begin{verbatim}
e6:= [xi^6](E0/kappa)
   =-2 int Gamma(v1,v5)-2 int Gamma(v2,v4)-int Gamma(v3,v3).
\end{verbatim}

The original fifth equation is Kv5=2Q\_H Gamma(v1,v4)+2Q\_H
Gamma(v2,v3). Integrating against the zero-Haar-mean v1, and using int
Gamma(f,h)=int f Kh, gives

\begin{verbatim}
e6=-4 int v1 Gamma(v1,v4)-4 int v1 Gamma(v2,v3)
    -2 int Gamma(v2,v4)-int Gamma(v3,v3).             (E1)
\end{verbatim}

This identity includes all fifth-source contributions; none is estimated
or discarded. For a specified original multiset of six faces, expand
every product in E1 over ordered allocations of that multiset to the
displayed sources. A repeated label is allocated by its actual
multiplicity, as in Q11. The resulting coefficients depend only on
sources through degree four, all evaluated in the adjoining tables.

\subsection{E2. Exact original Haar
integration}\label{e2.-exact-original-haar-integration}

Each word is evaluated in its original SU(2) edge variables. For an edge
of even total occurrence r, its full spin-zero projector on that
polynomial is

\begin{verbatim}
P_e,0=product_(j=1,...,r/2)[I-E_e/(j(j+1))].         (E2)
\end{verbatim}

Odd occurrence has zero Haar integral, proved by U\_e-\textgreater-U\_e.
Formula E2 follows by its values1 on Casimir0 and0 on each original
positive spin Casimir in that edge's complete occurrence range. Its
output is independent of U\_e. Thus replacing U\_e by I after E2 is the
actual original-edge integration, not a substitution for an unevaluated
integral. Repeat over every edge. Q4-Q7 make each step a rational
original trace calculation. The code retains the resulting constants,
including tr(I)=2.

The integration algorithm is independently checked on the fundamental
character moments1,2,5 at powers2,4,6. For every sixth-order coefficient
below, the same final number is obtained from the separate
Rayleigh-Schrodinger recurrence for the original K-xi S. Its overall
vacuum amplitude remains a separate scalar factor; the energy
coefficients are independent of that factor because the original
eigen-equation is homogeneous in it. The numerical polynomial entries in
that recurrence are coefficient multipliers, not an assigned physical
vacuum norm. The actual amplitude and original unit Haar norm are
returned by R9.

Explicitly, with multivariate face couplings, write those coefficient
multipliers as u\_C, u\_empty=1, int u\_C=0 for C nonempty, and keep the
overall scalar amplitude a as a separate factor a u\_C. The energy
coefficients satisfy

\begin{verbatim}
e_C=-sum_(p distinct in C) int W_p u_(C-{p}),
u_C=K^(-1)Q_H[sum_(p distinct in C)W_p u_(C-{p})
               +sum_(empty!=A proper submultiset C)e_A u_(C-A)].    (E3)
\end{verbatim}

Both equations follow by comparing each original face-coupling monomial
in K(a sum u\_C x\^{}C)-sum\_p x\_p W\_p(a sum u\_C x\^{}C)=e(x)(a sum
u\_C x\^{}C). The same scalar a appears on both sides; no energy or
coupling is rescaled. E3 and E1 give the same six local coefficients.
The first calculation uses the logarithmic-vacuum equation, while E3
uses the original linear eigen-equation.

\subsection{E3. Exhaustion of the sixth-order connected
supports}\label{e3.-exhaustion-of-the-sixth-order-connected-supports}

For each face p, changing U\_e to -U\_e on one original edge multiplies
a source coefficient by the parity of that edge's total original
occurrences. Differentiation and all Casimir inverses commute with this
center action. A nonzero Haar scalar therefore requires the set of
odd-multiplicity faces to have even boundary over F2.

A nonempty finite even-boundary set of elementary cubic faces has at
least six faces. Equality is the six-face boundary of one original cube.
Here is an explicit finite filling argument. Let c12,c13,c23 be its
binary face arrays. Set the binary cube array

\begin{verbatim}
b(n1,n2,n3)=sum_(k>n3)c12(n1,n2,k) mod2.
\end{verbatim}

The edge-boundary equations imply that the sum over k of c12 is constant
in n1 and n2; finite support makes that constant zero. Thus b is finite
also toward negative n3. Its upper-minus-lower cube boundary has the
required c12 array. The remaining face difference is a finite2-cycle
with c12=0. The boundary equations on directions1 and2 force its c13,c23
arrays to be constant in direction3; finite support makes them zero.
This proves the finite filling and its exact boundary map.

For a nonempty finite cube array b, each nonempty column along an axis
has at least two boundary faces normal to that axis. Therefore its
boundary has at least twice the sum of the cardinalities of its three
coordinate projections, hence at least six. Equality forces each
projection to be a singleton, so b is exactly one cube. This proves the
stated six-face classification without deleting nonreduced or repeated
face labels.

At total order six, all-even multiplicities are6,4+2,2+2+2. Their
connected distinct-face supports are respectively a single face, an
adjacent pair, or a path/common-edge/corner triple. The only all-odd
surviving support is one complete cube. Disconnected supports have zero
logarithmic energy coefficient by Q11 and E1; equivalently the original
operator on disjoint active link supports is the sum of its commuting
tensor-factor operators, so their scalar eigen-energy has no mixed
coefficient. The original scalar and all variables remain in that tensor
comparison.

\subsection{E4. Complete local coefficient
table}\label{e4.-complete-local-coefficient-table}

The coefficient for each indicated original multiset is

{\def\LTcaptype{none} % do not increment counter
\begin{longtable}[]{@{}lr@{}}
\toprule\noalign{}
Original multiset & e\_C \\
\midrule\noalign{}
\endhead
\bottomrule\noalign{}
\endlastfoot
\{p,p,p,p,p,p\} & -289/77760 \\
\{p,p,p,p,q,q\}, p\textasciitilde q & 22285/47309184 \\
\{p,p,q,q,r,r\}, distinct path & -4909/118272960 \\
Same multiplicities, common-edge triple & 244/4312035 \\
Same multiplicities, cube corner & -212/542997 \\
Six distinct faces bounding one cube & -83/1944 \\
\end{longtable}
}

The first pair row has the separate original role exchange
p\textless-\textgreater q; a sum over unordered adjacent pairs counts
both roles. The complete finite evidence includes both pair embeddings
(coplanar and perpendicular), all nine original cubic-orbit embeddings
of distinct triples, and the exact coefficients computed from both E1
and E3.

For clarity, the four terms in E1 for the cube are

\begin{verbatim}
-43/1458, -7/972, -13/2916, -1/648.
\end{verbatim}

For the single face they are

\begin{verbatim}
-17/5832, -1/1458, 17/46656, -7/14580.
\end{verbatim}

Their sums give the corresponding rows.
{\ttfamily s\allowbreak{}i\allowbreak{}x\allowbreak{}t\allowbreak{}h\allowbreak{}\_\allowbreak{}e\allowbreak{}n\allowbreak{}e\allowbreak{}r\allowbreak{}g\allowbreak{}y\allowbreak{}.\allowbreak{}j\allowbreak{}s\allowbreak{}o\allowbreak{}n\allowbreak{}}
retains all four terms for every row, together with their original face
words and multiplicities.

\subsection{E5. A second, path-resolved cube
proof}\label{e5.-a-second-path-resolved-cube-proof}

Orient the six cube faces outward. Reversing a whole SU(2) loop
preserves its trace, so this is the exact original trace product. Each
of the twelve edges now occurs once as U and once as U\^{}(-1). The
original integral

\begin{verbatim}
int U_ij conjugate(U_kl)=delta_ik delta_jl/2
\end{verbatim}

supplies twelve factors1/2. Its delta identifications leave one free
fundamental color at each of the eight vertices. Hence

\begin{verbatim}
int product_(six cube faces) W_p=2^8/2^12=1/16.     (E4)
\end{verbatim}

In the independent six-distinct-face perturbation coefficient, each
proper subset has zero scalar energy coefficient by E3 and the
even-boundary argument. In any original ordering, edges occurring twice
in an intermediate subset occur in no later face. Integration of those
edges projects them onto spin zero. Every edge occurring once carries
spin1/2. The original intermediate Casimir is consequently exactly
(3/4)\textbar boundary(subset)\textbar. These projections commute with
later independent edge variables and with the original total K.
Therefore the full cube coefficient is

\begin{verbatim}
-(1/16) sum_(pi in S6) product_(j=1)^5
                    4/[3 |boundary(first j faces of pi)|].       (E5)
\end{verbatim}

No intermediate denominator is replaced by the plaquette denominator3.
The complete path count is

{\def\LTcaptype{none} % do not increment counter
\begin{longtable}[]{@{}lr@{}}
\toprule\noalign{}
Five successive boundary lengths & Multiplicity \\
\midrule\noalign{}
\endhead
\bottomrule\noalign{}
\endlastfoot
(4,8,8,8,4) & 48 \\
(4,8,8,6,4) & 96 \\
(4,6,8,8,4) & 96 \\
(4,6,8,6,4) & 192 \\
(4,6,6,6,4) & 288 \\
\end{longtable}
}

There are720 paths. Their rational sum before the1/16 factor is166/243.
Equation E5 therefore gives exactly -83/1944, matching the complete
coefficient recurrence. This retains all original intermediate
boundaries, their denominators, and the terminal Haar contraction. It is
the first closed-surface scalar term on six distinct original faces.

\subsection{E6. All original finite-box factors and the complete
higher-order
remainder}\label{e6.-all-original-finite-box-factors-and-the-complete-higher-order-remainder}

Let m=2L. Denote by M,J,P,T,C,B the numbers of original faces, unordered
adjacent pairs, three-face paths, common-edge triples, corners and
cubes. Their exact values are

\begin{verbatim}
M=3m^2(m+1), J=6m(3m^2-1),
P=138m^3-126m^2-24m+12,
T=12m^2(m-1), C=8m^3, B=m^3.                       (E6)
\end{verbatim}

For T sum binom(r\_e,3) over the original edge degrees r\_e in
\{2,3,4\}. There are3m(m-1)\^{}2 degree-four edges and12m(m-1)
degree-three edges, giving the stated T. Each cube has eight three-face
corners. To obtain P, count centered face-adjacency wedges and subtract
three for each triangular triple. A face degree is s\_i+s\_j+4r\_k-4,
where s\_i,s\_j are3 at one of the two cell-end positions and4
otherwise, and r\_k is1 at either normal boundary and2 internally. The
direct count is

\begin{verbatim}
sum_p binom(deg p,2)=198m^3-162m^2-24m+12.
\end{verbatim}

Subtract3(T+C), proving P. These formulas preserve every original open
boundary. The verifier independently enumerates the actual boxes
L=2,3,4.

Summing all rows in E4 gives

\begin{verbatim}
e6= -289M/77760 +22285J/23654592 -4909P/118272960
                   +244T/4312035 -212C/542997 -83B/1944
   =-(211396463m^3+30959193m^2+21845782m+2336684)/4691494080.   (E7)
\end{verbatim}

Thus the original full ground energy is

\begin{verbatim}
E0,L=2kappa M xi-kappa M xi^2/3
   +kappa(5M/216-2J/1053)xi^4+kappa e6 xi^6+R_ge8.  (E8)
\end{verbatim}

The parent analyticity proof on the complex circle
\textbar xi\textbar=1/60 bounds the complete scalar C(v)
by49\textbar E\_L\textbar/1200 there. The new source agrees
coefficientwise with that proved construction. The exact center
involution makes E0-2kappa Mxi even. Cauchy's estimate and the full
geometric tail therefore give

\begin{verbatim}
|R_ge8| <= (49kappa |E_L|/1200)
             (60|xi|)^8/[1-(60|xi|)^2],
0<|xi|<1/60.                                       (E9)
\end{verbatim}

The finite coefficient polynomial is accompanied by all higher orders
through this bound. Its validity range remains stated; it is not
silently extended past the Cauchy circle.

At L=2, e6=-3528610133/1172873520. Along the original spatial-volume
limit m-\textgreater infinity, the coefficient per face is

\begin{verbatim}
lim e6/M=-211396463/14074482240.                     (E10)
\end{verbatim}

This is a limit of the specified Taylor coefficient with its complete
original counts. It does not by itself exchange a spatial or ultraviolet
limit with a coupling series beyond the proved domain.

\subsection{E7. Literature and scope}\label{e7.-literature-and-scope}

Schuette-Zheng-Hamer's coupled-cluster/character method and Llewellyn
Smith-Watson's shifted linked-cluster formulation are primary
antecedents, arXiv:hep-lat/9603026v1 and hep-lat/9212025v1. The former's
operator dictionary to Q1 is R23. The latter explicitly describes its
truncation and cluster-shift prescriptions; this calculation keeps the
original products and uses the full residual bounds R3-R8 and E9. This
establishes the method correspondence and the actual all-order error
returned here; it is not a claim that the general expansion method or
every displayed coefficient is historically new. The scoped search did
not establish an independent published match or global priority for the
full spatial sixth coefficient E7.

\subsection{E8. Return to an actual measured plaquette
observable}\label{e8.-return-to-an-actual-measured-plaquette-observable}

For real 0\textless xi\textless1/60 at fixed original kappa, the
constructed simple vacuum is differentiable in xi. Its original mass is
one, hence
{\ttfamily 2\allowbreak{} \allowbreak{}R\allowbreak{}e\allowbreak{}\textless{}\allowbreak{}p\allowbreak{}s\allowbreak{}i\allowbreak{},\allowbreak{}p\allowbreak{}a\allowbreak{}r\allowbreak{}t\allowbreak{}i\allowbreak{}a\allowbreak{}l\allowbreak{}\_\allowbreak{}x\allowbreak{}i\allowbreak{} \allowbreak{}p\allowbreak{}s\allowbreak{}i\allowbreak{}\textgreater{}\allowbreak{}=\allowbreak{}0\allowbreak{}}.
Differentiating
{\ttfamily H\allowbreak{} \allowbreak{}p\allowbreak{}s\allowbreak{}i\allowbreak{}=\allowbreak{}E\allowbreak{}0\allowbreak{} \allowbreak{}p\allowbreak{}s\allowbreak{}i\allowbreak{}}
and testing with the same psi therefore proves the exact
Hellmann-Feynman identity

\begin{verbatim}
dE0/dxi = kappa(2M-sum_p <W_p>_rho).
\end{verbatim}

No derivative of the vacuum is omitted: its two terms cancel through the
displayed mass derivative and the actual eigen-equation. Define the
original mean half-trace as an explicit observation map

\begin{verbatim}
P_L(xi)=(1/(2M))sum_p <W_p>_rho.
\end{verbatim}

Keeping M and the complete coefficients e4,e6 from E8 gives

\begin{verbatim}
P_L(xi)=xi/3-(2e4/M)xi^3-(3e6/M)xi^5+R_P,L,
|R_P,L| <= (49|E_L|/(40M))
     (60xi)^7 [8-6(60xi)^2]/[1-(60xi)^2]^2.         (E11)
\end{verbatim}

To prove the tail, the same original Cauchy bound used in E9 gives
{\ttfamily |\allowbreak{}e\allowbreak{}\_\allowbreak{}(\allowbreak{}2\allowbreak{}n\allowbreak{})\allowbreak{}|\allowbreak{}\textless{}\allowbreak{}=\allowbreak{}(\allowbreak{}4\allowbreak{}9\allowbreak{}|\allowbreak{}E\allowbreak{}\_\allowbreak{}L\allowbreak{}|\allowbreak{}/\allowbreak{}1\allowbreak{}2\allowbreak{}0\allowbreak{}0\allowbreak{})\allowbreak{}6\allowbreak{}0\allowbreak{}\textasciicircum{}\allowbreak{}(\allowbreak{}2\allowbreak{}n\allowbreak{})\allowbreak{}}
for every n\textgreater=4. Differentiate its absolutely convergent
series on \textbar xi\textbar\textless1/60. The exact scalar sum

\begin{verbatim}
sum_(n>=4)2n r^(2n-1)=r^7(8-6r^2)/(1-r^2)^2
\end{verbatim}

follows by differentiating r\^{}8/(1-r\^{}2). Its factor60 and the
factor1/(2M) in the actual observation give precisely E11. The ratio
\textbar E\_L\textbar/M=(m+1)/m is retained and is at most5/4 for
L\textgreater=2; consequently the entire error is independent of
exterior volume after this explicitly defined spatial average.

For each fixed coefficient, all connected source supports have a finite
number of original faces. Their interior translates contribute the same
coefficient, and the boundary embeddings have vanishing proportion as
m-\textgreater infinity. The bounds just used are uniform after division
by M. Dominated power-series summation therefore proves the full
mean-observable limit on \textbar xi\textbar\textless1/60, as well as
differentiation of its energy-density series. Its first terms are

\begin{verbatim}
P_infinity(xi)=xi/3-(11/468)xi^3
                   +(211396463/4691494080)xi^5+R_P,infinity,
|R_P,infinity| <= (49/40)
     (60xi)^7 [8-6(60xi)^2]/[1-(60xi)^2]^2.         (E12)
\end{verbatim}

This is a fixed-lattice-spacing observable return. No continuum scaling
limit is used in the coefficientwise or dominated-series argument.

\subsection{E9. Explicit correspondence to a current variational
benchmark
convention}\label{e9.-explicit-correspondence-to-a-current-variational-benchmark-convention}

Spriggs, Greplova, Carrasquilla and Nys, arXiv:2509.12323v1, equation1,
use the physical convention

\begin{verbatim}
H_N(g_N,a)=g_N^2/a [K/2+(4/g_N^4)sum_p(1-W_p/2)].
\end{verbatim}

Specialize this formula to the same original edge and face sets as Q1.
Then the literal coefficient substitution is

\begin{verbatim}
g_N=2g, lambda_N=4/g_N^4=xi,
H_N(2g,a)=H_original(g,a),
(a/g_N^2)H_N=H_original/(2kappa).                   (E13)
\end{verbatim}

The generator commutators and E\^{}2=-Delta\_(S3)/4 in their
equations1-3 agree with the original T\_alpha=-i sigma\_alpha/2
coordinates through their Hermitian derivative iX. Direct substitution
in E13 gives exactly kappa and kappa xi for the two original
coefficients. The trace map in their equation4 is
W\_p-\textgreater W\_p/2, precisely the measured observation in E11. All
energy and trace factors are exhibited.

Their reported finite numerical lattices have their own boundary choice.
To retain the comparison on a periodic extension of our same vertex set,
let J be pullback from the old edge variables, with inverse adjoint the
Haar integral over the additional wrapping links. For
side2L+1\textgreater=5 its actual operator defect is

\begin{verbatim}
H_periodic J-J H_open
   =kappa xi sum_(p in P_periodic\P_open)(2-W_p)J.  (E14)
\end{verbatim}

The electric terms on the added links annihilate the pullback, and all
old terms agree; expansion proves E14. Thus E13 is a convention
correspondence on matched index sets, and E14 keeps the added
interactions needed for the other boundary choice. No numerical data
from that paper are presented as an evaluation or independent
verification of E7 or E12. The earlier 1985 perturbative source cited in
that paper was identified bibliographically but its full text was not
retrieved here. The new coefficients and remainder are offered as
explicit original-convention benchmark data, with their written
derivation and exact replay.

\section{Completed fourth-order source in the original lattice
coordinates}\label{completed-fourth-order-source-in-the-original-lattice-coordinates}

Audit and continuation, 17 September 2026. The requested basis is the
supplied
{\ttfamily P\allowbreak{}a\allowbreak{}s\allowbreak{}t\allowbreak{}e\allowbreak{}d\allowbreak{} \allowbreak{}m\allowbreak{}a\allowbreak{}r\allowbreak{}k\allowbreak{}d\allowbreak{}o\allowbreak{}w\allowbreak{}n\allowbreak{}(\allowbreak{}6\allowbreak{})\allowbreak{}.\allowbreak{}m\allowbreak{}d\allowbreak{}}.
Its original bytes are preserved in
{\ttfamily i\allowbreak{}n\allowbreak{}p\allowbreak{}u\allowbreak{}t\allowbreak{}/\allowbreak{}}.
The published cubic predecessor was read at Git revision
{\ttfamily 9\allowbreak{}1\allowbreak{}4\allowbreak{}3\allowbreak{}4\allowbreak{}b\allowbreak{}6\allowbreak{}9\allowbreak{}6\allowbreak{}2\allowbreak{}0\allowbreak{}6\allowbreak{}2\allowbreak{}b\allowbreak{}d\allowbreak{}8\allowbreak{}0\allowbreak{}4\allowbreak{}3\allowbreak{}9\allowbreak{}d\allowbreak{}4\allowbreak{}c\allowbreak{}b\allowbreak{}2\allowbreak{}c\allowbreak{}a\allowbreak{}e\allowbreak{}9\allowbreak{}d\allowbreak{}2\allowbreak{}4\allowbreak{}7\allowbreak{}9\allowbreak{}2\allowbreak{}6\allowbreak{}4\allowbreak{}d\allowbreak{}d\allowbreak{}e\allowbreak{}};
source recovery and its limits are in
{\ttfamily s\allowbreak{}o\allowbreak{}u\allowbreak{}r\allowbreak{}c\allowbreak{}e\allowbreak{}s\allowbreak{}/\allowbreak{}}.
This note gives a complete, executable fourth-order source catalogue. It
makes no new claim about a uniform gap or the four-dimensional
continuum.

\subsection{1. Operator, source, and the distinction that needs to be
retained}\label{operator-source-and-the-distinction-that-needs-to-be-retained}

Use the original finite open cubic graph and positive edges, generators
{\ttfamily T\allowbreak{}\_\allowbreak{}a\allowbreak{}=\allowbreak{}-\allowbreak{}i\allowbreak{} \allowbreak{}s\allowbreak{}i\allowbreak{}g\allowbreak{}m\allowbreak{}a\allowbreak{}\_\allowbreak{}a\allowbreak{}/\allowbreak{}2\allowbreak{}},
derivatives
{\ttfamily X\allowbreak{}\_\allowbreak{}e\allowbreak{},\allowbreak{}a\allowbreak{}},
and Haar probability. Write

\[
 K=-\sum_{e,a}X_{e,a}^{2},\qquad
 H=\kappa[K+\xi(2M-S)],\quad S=\sum_pW_p,
 \quad\kappa=2g^2/a,\quad\xi=1/(4g^4).
\tag{Q1}
\] The original ordered face is
{\ttfamily W\allowbreak{}\_\allowbreak{}(\allowbreak{}n\allowbreak{};\allowbreak{}i\allowbreak{},\allowbreak{}j\allowbreak{})\allowbreak{}=\allowbreak{}t\allowbreak{}r\allowbreak{}(\allowbreak{}U\allowbreak{}\_\allowbreak{}i\allowbreak{}(\allowbreak{}n\allowbreak{})\allowbreak{}U\allowbreak{}\_\allowbreak{}j\allowbreak{}(\allowbreak{}n\allowbreak{}+\allowbreak{}e\allowbreak{}\_\allowbreak{}i\allowbreak{})\allowbreak{}U\allowbreak{}\_\allowbreak{}i\allowbreak{}(\allowbreak{}n\allowbreak{}+\allowbreak{}e\allowbreak{}\_\allowbreak{}j\allowbreak{})\allowbreak{}\textasciicircum{}\allowbreak{}-\allowbreak{}1\allowbreak{} \allowbreak{}U\allowbreak{}\_\allowbreak{}j\allowbreak{}(\allowbreak{}n\allowbreak{})\allowbreak{}\textasciicircum{}\allowbreak{}-\allowbreak{}1\allowbreak{})\allowbreak{}}.
Let {\ttfamily P\allowbreak{}\_\allowbreak{}H\allowbreak{}} be Haar
expectation and
{\ttfamily Q\allowbreak{}\_\allowbreak{}H\allowbreak{}=\allowbreak{}I\allowbreak{}-\allowbreak{}P\allowbreak{}\_\allowbreak{}H\allowbreak{}}.
The original physical source is the zero-Haar-mean logarithmic vacuum
{\ttfamily v\allowbreak{}=\allowbreak{}Q\allowbreak{}\_\allowbreak{}H\allowbreak{} \allowbreak{}l\allowbreak{}o\allowbreak{}g\allowbreak{} \allowbreak{}p\allowbreak{}s\allowbreak{}i\allowbreak{}},
with scalar return

\[
 \psi=\exp(v+c),\quad c=-\tfrac12\log\int e^{2v}dU,\qquad
 E_0=\kappa[2M\xi-P_H\Gamma(v,v)],
 \quad\Gamma(f,h)=\sum_{e,a}(X_{e,a}f)(X_{e,a}h).
\tag{Q2}
\] The scalar {\ttfamily c\allowbreak{}} and the additive original
energy have not been discarded. The coefficient equation is

\[
 v=\xi v_1+B(v,v),\qquad v_1=S/3,\quad
 B(f,h)=K^{-1}Q_H\Gamma(f,h).
\tag{Q3}
\] The inverse in Q3 is on its actual nonconstant physical polynomial
image. On a fixed finite graph its uniqueness follows from the positive
Casimir and Haar centering.

Consequently, at
{\ttfamily q\allowbreak{}2\allowbreak{}=\allowbreak{}x\allowbreak{}i\allowbreak{}*\allowbreak{}v\allowbreak{}1\allowbreak{}+\allowbreak{}x\allowbreak{}i\allowbreak{}\textasciicircum{}\allowbreak{}2\allowbreak{}*\allowbreak{}v\allowbreak{}2\allowbreak{}},
the exact residual and full fourth coefficient are different specified
expressions:

\[
 R_2=\xi^3v_3+\xi^4b_{22},\quad b_{22}=B(v_2,v_2),\qquad
 \boxed{v_4=2B(v_1,v_3)+b_{22}.}
\tag{Q4}
\] The uploaded L5--L7 give
{\ttfamily b\allowbreak{}2\allowbreak{}2\allowbreak{}}. Equation Q4
retains the additional incoming source and its unique
{\ttfamily K\allowbreak{}\textasciicircum{}\allowbreak{}\{\allowbreak{}-\allowbreak{}1\allowbreak{}\}\allowbreak{}Q\allowbreak{}\_\allowbreak{}H\allowbreak{}}
primitive. This package computes their signed sum before any norm bound.

\subsection{2. The exact tree-coordinate map and its
inverse}\label{the-exact-tree-coordinate-map-and-its-inverse}

For the original edge union E of a connected face cluster, choose the
stored spanning tree T, root o, and original oriented chords
{\ttfamily c\allowbreak{}\_\allowbreak{}1\allowbreak{},\allowbreak{}.\allowbreak{}.\allowbreak{}.\allowbreak{},\allowbreak{}c\allowbreak{}\_\allowbreak{}r\allowbreak{}},
where
{\ttfamily r\allowbreak{}=\allowbreak{}|\allowbreak{}E\allowbreak{}|\allowbreak{}-\allowbreak{}|\allowbreak{}V\allowbreak{}|\allowbreak{}+\allowbreak{}1\allowbreak{}}.
For a vertex v let
{\ttfamily h\allowbreak{}\_\allowbreak{}v\allowbreak{}(\allowbreak{}U\allowbreak{})\allowbreak{}}
be the original ordered product along the tree path from o to v, using
inverse links on reverse traversals. Define

\[
 Z_j=h_{s(c_j)}(U)\,U_{c_j}\,h_{t(c_j)}(U)^{-1},
 \qquad
 \Phi:SU(2)^E\longrightarrow SU(2)^T\times SU(2)^r,
 \quad U\longmapsto(U_T,Z).
\tag{Q5}
\] Its inverse keeps
{\ttfamily U\allowbreak{}\_\allowbreak{}T\allowbreak{}} and sets
{\ttfamily U\allowbreak{}\_\allowbreak{}c\allowbreak{}=\allowbreak{}h\allowbreak{}\_\allowbreak{}s\allowbreak{}(\allowbreak{}U\allowbreak{}\_\allowbreak{}T\allowbreak{})\allowbreak{}\textasciicircum{}\allowbreak{}-\allowbreak{}1\allowbreak{} \allowbreak{}Z\allowbreak{}\_\allowbreak{}c\allowbreak{} \allowbreak{}h\allowbreak{}\_\allowbreak{}t\allowbreak{}(\allowbreak{}U\allowbreak{}\_\allowbreak{}T\allowbreak{})\allowbreak{}}.
Both compositions are literal identities. Successive left and right Haar
translations in the chord variables prove
{\ttfamily d\allowbreak{}U\allowbreak{}=\allowbreak{}d\allowbreak{}U\allowbreak{}\_\allowbreak{}T\allowbreak{} \allowbreak{}d\allowbreak{}Z\allowbreak{}}.
A gauge-invariant function pulls back from a function of Z invariant
under simultaneous conjugation; its original Haar norm is exactly its
product-Haar norm in these variables. The remaining root gauge action is
retained, rather than treated as an extra independent observable.

At the slice
{\ttfamily U\allowbreak{}\_\allowbreak{}T\allowbreak{}=\allowbreak{}I\allowbreak{}},
an original chord derivative is
{\ttfamily L\allowbreak{}\_\allowbreak{}(\allowbreak{}j\allowbreak{},\allowbreak{}a\allowbreak{})\allowbreak{}},
the original left generator of Z\_j. For a tree edge e, let
{\ttfamily s\allowbreak{}\_\allowbreak{}e\allowbreak{}(\allowbreak{}c\allowbreak{})\allowbreak{},\allowbreak{}t\allowbreak{}\_\allowbreak{}e\allowbreak{}(\allowbreak{}c\allowbreak{})\allowbreak{}}
be the zero-one indicators that the original paths to the two chord
endpoints cross e. Up to the common path-orientation sign for that one
e, its field is

\[
 V_{e,a}=\sum_c[s_e(c)L_{c,a}-t_e(c)R_{c,a}].
\tag{Q6}
\] The common sign cancels in the original squared generator and in
Gamma. The reduced kinetic operator is the exact original operator

\[
 K_T=-\sum_{e\in E,a}V_{e,a}^{2}.
\tag{Q7}
\] To prove Q6--Q7, vary one original tree link as
{\ttfamily e\allowbreak{}x\allowbreak{}p\allowbreak{}(\allowbreak{}t\allowbreak{}T\allowbreak{}\_\allowbreak{}a\allowbreak{})\allowbreak{}U\allowbreak{}\_\allowbreak{}e\allowbreak{}},
and then apply Q5. Each affected chord becomes
{\ttfamily e\allowbreak{}x\allowbreak{}p\allowbreak{}(\allowbreak{}e\allowbreak{}p\allowbreak{}s\allowbreak{}i\allowbreak{}l\allowbreak{}o\allowbreak{}n\allowbreak{}*\allowbreak{}s\allowbreak{}\_\allowbreak{}e\allowbreak{}*\allowbreak{}t\allowbreak{}T\allowbreak{}\_\allowbreak{}a\allowbreak{})\allowbreak{} \allowbreak{}Z\allowbreak{}\_\allowbreak{}c\allowbreak{} \allowbreak{}e\allowbreak{}x\allowbreak{}p\allowbreak{}(\allowbreak{}-\allowbreak{}e\allowbreak{}p\allowbreak{}s\allowbreak{}i\allowbreak{}l\allowbreak{}o\allowbreak{}n\allowbreak{}*\allowbreak{}t\allowbreak{}\_\allowbreak{}e\allowbreak{}*\allowbreak{}t\allowbreak{}T\allowbreak{}\_\allowbreak{}a\allowbreak{})\allowbreak{}}.
Its first and second derivatives are precisely the stated fields. Gauge
covariance returns this identity from the tree slice to every original
configuration. The checker additionally verifies it by independent
second-order quaternion jets through every original edge of all 78
representatives. Those are supplementary coordinate tests; Q5--Q7 are
the analytic identification of the operator.

\subsection{3. Complete polynomial coordinates with exact
relations}\label{complete-polynomial-coordinates-with-exact-relations}

Each stored chord has quaternion coordinates
{\ttfamily Z\allowbreak{}=\allowbreak{}q\allowbreak{}0\allowbreak{} \allowbreak{}I\allowbreak{}-\allowbreak{}i\allowbreak{}(\allowbreak{}q\allowbreak{}1\allowbreak{} \allowbreak{}s\allowbreak{}i\allowbreak{}g\allowbreak{}m\allowbreak{}a\allowbreak{}1\allowbreak{}+\allowbreak{}q\allowbreak{}2\allowbreak{} \allowbreak{}s\allowbreak{}i\allowbreak{}g\allowbreak{}m\allowbreak{}a\allowbreak{}2\allowbreak{}+\allowbreak{}q\allowbreak{}3\allowbreak{} \allowbreak{}s\allowbreak{}i\allowbreak{}g\allowbreak{}m\allowbreak{}a\allowbreak{}3\allowbreak{})\allowbreak{}},
with
{\ttfamily q\allowbreak{}0\allowbreak{}\textasciicircum{}\allowbreak{}2\allowbreak{}+\allowbreak{}q\allowbreak{}1\allowbreak{}\textasciicircum{}\allowbreak{}2\allowbreak{}+\allowbreak{}q\allowbreak{}2\allowbreak{}\textasciicircum{}\allowbreak{}2\allowbreak{}+\allowbreak{}q\allowbreak{}3\allowbreak{}\textasciicircum{}\allowbreak{}2\allowbreak{}=\allowbreak{}1\allowbreak{}}.
The algebra is the rational polynomial quotient by these original sphere
relations, one per chord. In the stored representative monomials each q0
exponent is zero or one. The relation maps are explicit:
{\ttfamily q\allowbreak{}0\allowbreak{}\textasciicircum{}\allowbreak{}2\allowbreak{} \allowbreak{}-\allowbreak{}\textgreater{}\allowbreak{} \allowbreak{}1\allowbreak{}-\allowbreak{}q\allowbreak{}1\allowbreak{}\textasciicircum{}\allowbreak{}2\allowbreak{}-\allowbreak{}q\allowbreak{}2\allowbreak{}\textasciicircum{}\allowbreak{}2\allowbreak{}-\allowbreak{}q\allowbreak{}3\allowbreak{}\textasciicircum{}\allowbreak{}2\allowbreak{}};
their inverse interpretation is evaluation on the original sphere. The
ideal generators have relatively prime leading monomials in different
chord groups, so this representation is unique. No state or energy
metric is replaced by a Euclidean coefficient norm.

The generator matrices are the original quaternion multiplications by
{\ttfamily (\allowbreak{}0\allowbreak{},\allowbreak{}e\allowbreak{}\_\allowbreak{}a\allowbreak{}/\allowbreak{}2\allowbreak{})\allowbreak{}}
on the left or right. Thus all their entries are rational halves. The
original moments are

\[
 \int q_0^{d_0}q_1^{d_1}q_2^{d_2}q_3^{d_3}\,dq=0
 \quad\hbox{if any }d_i\hbox{ is odd},
\] \[
 \int\prod_iq_i^{d_i}dq
 =\frac{\prod_i(d_i-1)!!}{4\cdot6\cdots(2+\sum_i d_i)}
 \quad\hbox{when all }d_i\hbox{ are even}.
\tag{Q8}
\] The empty denominator and
{\ttfamily (\allowbreak{}-\allowbreak{}1\allowbreak{})\allowbreak{}!\allowbreak{}!\allowbreak{}}
are one. Products over distinct chords use independent original Haar
factors. Q8 supplies every mean and full polynomial Gram used in the
package.

\subsection{4. The exact finite inverse and coefficient
recurrence}\label{the-exact-finite-inverse-and-coefficient-recurrence}

Introduce formal face source coordinates
{\ttfamily l\allowbreak{}a\allowbreak{}m\allowbreak{}b\allowbreak{}d\allowbreak{}a\allowbreak{}\_\allowbreak{}p\allowbreak{}},
with the original physical diagonal map
{\ttfamily l\allowbreak{}a\allowbreak{}m\allowbreak{}b\allowbreak{}d\allowbreak{}a\allowbreak{}\_\allowbreak{}p\allowbreak{}=\allowbreak{}x\allowbreak{}i\allowbreak{}}
after all coefficients have been computed. The exact physical operator
on that diagonal is Q1. For a multiindex nu, its original edge
multiplicity is
{\ttfamily m\allowbreak{}\_\allowbreak{}e\allowbreak{}(\allowbreak{}n\allowbreak{}u\allowbreak{})\allowbreak{}=\allowbreak{}s\allowbreak{}u\allowbreak{}m\allowbreak{}\_\allowbreak{}(\allowbreak{}p\allowbreak{} \allowbreak{}c\allowbreak{}o\allowbreak{}n\allowbreak{}t\allowbreak{}a\allowbreak{}i\allowbreak{}n\allowbreak{}i\allowbreak{}n\allowbreak{}g\allowbreak{} \allowbreak{}e\allowbreak{})\allowbreak{} \allowbreak{}n\allowbreak{}u\allowbreak{}\_\allowbreak{}p\allowbreak{}}.
Its allowed edge spins are
{\ttfamily j\allowbreak{}\_\allowbreak{}e\allowbreak{}=\allowbreak{}m\allowbreak{}\_\allowbreak{}e\allowbreak{}/\allowbreak{}2\allowbreak{},\allowbreak{}m\allowbreak{}\_\allowbreak{}e\allowbreak{}/\allowbreak{}2\allowbreak{}-\allowbreak{}1\allowbreak{},\allowbreak{}.\allowbreak{}.\allowbreak{}.\allowbreak{}}
down to zero or one half. Therefore the finite list of possible positive
Casimirs on that source is

\[
 \mathcal C_\nu=
 \left\{\sum_ej_e(j_e+1)>0:
 j_e\in\{m_e/2,m_e/2-1,\ldots\}\right\}.
\tag{Q9}
\] Extra allowed values are harmless. The implementation proves that its
returned inverse solves the original equation, not just a matrix fit.
For a centered forcing f, start
{\ttfamily r\allowbreak{}0\allowbreak{}=\allowbreak{}f\allowbreak{}},
{\ttfamily y\allowbreak{}0\allowbreak{}=\allowbreak{}0\allowbreak{}},
and for each
{\ttfamily c\allowbreak{} \allowbreak{}i\allowbreak{}n\allowbreak{} \allowbreak{}C\allowbreak{}\_\allowbreak{}n\allowbreak{}u\allowbreak{}}
set

\[
 y_{j+1}=y_j+r_j/c,\qquad r_{j+1}=r_j-K_Tr_j/c.
\tag{Q10}
\] The exact identity is
{\ttfamily K\allowbreak{}\_\allowbreak{}T\allowbreak{} \allowbreak{}y\allowbreak{}\_\allowbreak{}j\allowbreak{}=\allowbreak{}f\allowbreak{}-\allowbreak{}r\allowbreak{}\_\allowbreak{}j\allowbreak{}}
at every stage. The final polynomial
{\ttfamily r\allowbreak{}\_\allowbreak{}j\allowbreak{}} is checked to be
zero in the original sphere quotient. This gives a unique centered
inverse because K\_T is a sum of nonnegative squares and its chord
fields alone force a zero-derivative function to be constant. All these
operations have finite degree because this is a coefficient calculation
at a stated order; no finite-spin approximation to the full Hamiltonian
or its spectrum is used.

For explicit arithmetic it is convenient to use a formal eigenline
representative
{\ttfamily u\allowbreak{}(\allowbreak{}l\allowbreak{}a\allowbreak{}m\allowbreak{}b\allowbreak{}d\allowbreak{}a\allowbreak{})\allowbreak{}=\allowbreak{}1\allowbreak{}+\allowbreak{}s\allowbreak{}u\allowbreak{}m\allowbreak{}\_\allowbreak{}(\allowbreak{}n\allowbreak{}u\allowbreak{}!\allowbreak{}=\allowbreak{}0\allowbreak{})\allowbreak{}u\allowbreak{}\_\allowbreak{}n\allowbreak{}u\allowbreak{} \allowbreak{}l\allowbreak{}a\allowbreak{}m\allowbreak{}b\allowbreak{}d\allowbreak{}a\allowbreak{}\textasciicircum{}\allowbreak{}n\allowbreak{}u\allowbreak{}},
{\ttfamily P\allowbreak{}\_\allowbreak{}H\allowbreak{} \allowbreak{}u\allowbreak{}\_\allowbreak{}n\allowbreak{}u\allowbreak{}=\allowbreak{}0\allowbreak{}},
of
{\ttfamily K\allowbreak{}-\allowbreak{}s\allowbreak{}u\allowbreak{}m\allowbreak{}\_\allowbreak{}p\allowbreak{} \allowbreak{}l\allowbreak{}a\allowbreak{}m\allowbreak{}b\allowbreak{}d\allowbreak{}a\allowbreak{}\_\allowbreak{}p\allowbreak{} \allowbreak{}W\allowbreak{}\_\allowbreak{}p\allowbreak{}},
with eigenvalue
{\ttfamily e\allowbreak{}(\allowbreak{}l\allowbreak{}a\allowbreak{}m\allowbreak{}b\allowbreak{}d\allowbreak{}a\allowbreak{})\allowbreak{}}.
This is an explicitly recorded choice of the scalar coordinate of the
eigenline. Its exact return is
{\ttfamily v\allowbreak{}=\allowbreak{}Q\allowbreak{}\_\allowbreak{}H\allowbreak{} \allowbreak{}l\allowbreak{}o\allowbreak{}g\allowbreak{} \allowbreak{}u\allowbreak{}},
followed by Q2 for the physical unit vacuum. The inverse coordinate
relation is
{\ttfamily u\allowbreak{}=\allowbreak{}e\allowbreak{}\textasciicircum{}\allowbreak{}v\allowbreak{}/\allowbreak{}(\allowbreak{}P\allowbreak{}\_\allowbreak{}H\allowbreak{} \allowbreak{}e\allowbreak{}\textasciicircum{}\allowbreak{}v\allowbreak{})\allowbreak{}},
whose denominator has constant coefficient one. The original physical
normalization is not dropped. The recurrences used are

\[
 F_\nu=\sum_{p:\nu_p>0}W_pu_{\nu-e_p},\qquad
 e_\nu=-P_HF_\nu,
\] \[
 K_Tu_\nu=Q_HF_\nu+
   \sum_{0<\mu<\nu}e_\mu u_{\nu-\mu},\quad P_Hu_\nu=0.
\tag{Q11}
\] Terms are interpreted coefficientwise, with
{\ttfamily m\allowbreak{}u\allowbreak{}\textless{}\allowbreak{}=\allowbreak{}n\allowbreak{}u\allowbreak{}}.
Writing
{\ttfamily l\allowbreak{}=\allowbreak{}l\allowbreak{}o\allowbreak{}g\allowbreak{} \allowbreak{}u\allowbreak{}},
differentiation of the formal series gives the complete coefficient
relation

\[
 l_\nu=u_\nu-\frac1{|\nu|}
     \sum_{0<\mu<\nu}|\mu|\,l_\mu u_{\nu-\mu},
 \qquad v_\nu=Q_Hl_\nu.
\tag{Q12}
\] As an independent polynomial check for each coefficient the program
verifies

\[
 K_Tv_\nu=
 Q_H\!\sum_{0<\mu<\nu}
 \Gamma_T(v_\mu,v_{\nu-\mu}),\qquad |\nu|>1,
\tag{Q13}
\] with
{\ttfamily K\allowbreak{}\_\allowbreak{}T\allowbreak{}v\allowbreak{}\_\allowbreak{}(\allowbreak{}e\allowbreak{}\_\allowbreak{}p\allowbreak{})\allowbreak{}=\allowbreak{}W\allowbreak{}\_\allowbreak{}p\allowbreak{}}
and
{\ttfamily 2\allowbreak{} \allowbreak{}G\allowbreak{}a\allowbreak{}m\allowbreak{}m\allowbreak{}a\allowbreak{}\_\allowbreak{}T\allowbreak{}(\allowbreak{}f\allowbreak{},\allowbreak{}h\allowbreak{})\allowbreak{}=\allowbreak{}(\allowbreak{}K\allowbreak{}\_\allowbreak{}T\allowbreak{} \allowbreak{}f\allowbreak{})\allowbreak{}h\allowbreak{}+\allowbreak{}f\allowbreak{}(\allowbreak{}K\allowbreak{}\_\allowbreak{}T\allowbreak{} \allowbreak{}h\allowbreak{})\allowbreak{}-\allowbreak{}K\allowbreak{}\_\allowbreak{}T\allowbreak{}(\allowbreak{}f\allowbreak{}h\allowbreak{})\allowbreak{}}.
These are identities of complete rational polynomials, not agreements
only at finitely many sampled configurations.

\subsection{5. All connected fourth-order supports are
included}\label{all-connected-fourth-order-supports-are-included}

The source operation Gamma differentiates a common original link. Its
inverse Casimir creates no new link. Consequently every nonzero
logarithmic coefficient has an edge-connected face multiset. Starting
from the four faces incident on the original anchor edge
{\ttfamily (\allowbreak{}0\allowbreak{},\allowbreak{}0\allowbreak{},\allowbreak{}0\allowbreak{};\allowbreak{}d\allowbreak{}i\allowbreak{}r\allowbreak{}e\allowbreak{}c\allowbreak{}t\allowbreak{}i\allowbreak{}o\allowbreak{}n\allowbreak{}0\allowbreak{})\allowbreak{}},
the following exhaustive recursion suffices: add a repeated present face
or a face sharing an original edge with a present face, and retain the
sorted multiset. Any connected multiset containing that anchor can be
recovered in this way by removing a leaf of a rooted spanning tree of
its adjacency graph, or one repeated occurrence. Counts by order are

\[
 4,\quad46,\quad612,\quad8621.
\tag{Q14}
\] The 48 signed coordinate permutations and translations act on the
original faces and links. A reversed edge maps to its inverse; this is a
Haar- and Casimir-preserving coordinate map. A reversed complete SU(2)
face word has the same trace. Each orbit is represented by its
explicitly stored minimum corner coordinates and ordered axes. The
fourth-order orbit counts are

{\def\LTcaptype{none} % do not increment counter
\begin{longtable}[]{@{}lr@{}}
\toprule\noalign{}
Face-multiplicity partition & Orbits \\
\midrule\noalign{}
\endhead
\bottomrule\noalign{}
\endlastfoot
4 & 1 \\
3+1 & 2 \\
2+2 & 2 \\
2+1+1 & 19 \\
1+1+1+1 & 54 \\
Total & 78 \\
\end{longtable}
}

{\ttfamily r\allowbreak{}e\allowbreak{}s\allowbreak{}u\allowbreak{}l\allowbreak{}t\allowbreak{}s\allowbreak{}/\allowbreak{}q\allowbreak{}u\allowbreak{}a\allowbreak{}r\allowbreak{}t\allowbreak{}i\allowbreak{}c\allowbreak{}\_\allowbreak{}c\allowbreak{}o\allowbreak{}e\allowbreak{}f\allowbreak{}f\allowbreak{}i\allowbreak{}c\allowbreak{}i\allowbreak{}e\allowbreak{}n\allowbreak{}t\allowbreak{}s\allowbreak{}/\allowbreak{}c\allowbreak{}l\allowbreak{}a\allowbreak{}s\allowbreak{}s\allowbreak{}\_\allowbreak{}0\allowbreak{}0\allowbreak{}.\allowbreak{}j\allowbreak{}s\allowbreak{}o\allowbreak{}n\allowbreak{}}
through
{\ttfamily c\allowbreak{}l\allowbreak{}a\allowbreak{}s\allowbreak{}s\allowbreak{}\_\allowbreak{}7\allowbreak{}7\allowbreak{}.\allowbreak{}j\allowbreak{}s\allowbreak{}o\allowbreak{}n\allowbreak{}}
contain all 4,044 nonzero rational monomials of these 78 fully
calculated source coefficients. They also give original faces, tree
links, chords, root, original derivative fields, coupling multiindex,
Haar mean and full Haar state/energy norms. These auxiliary Haar norms
are not asserted to be the interacting physical norms.

One class (index25) has four distinct faces but five chord variables:
the four-sided tube formed by the cube's lateral faces. Its original
graph has an extra independent cycle. The implementation retains it.
Replacing the chord count by the number of faces would change this
original source.

The source for any original finite box is reconstructed by summing each
contained connected multiset once, pulling its representative polynomial
back through the stored lattice symmetry and Q5, and setting every
formal face coordinate to xi. Q11--13 already include the correct
multiplicities; no additional division by 4! is applied.
{\ttfamily r\allowbreak{}e\allowbreak{}s\allowbreak{}u\allowbreak{}l\allowbreak{}t\allowbreak{}s\allowbreak{}/\allowbreak{}a\allowbreak{}n\allowbreak{}c\allowbreak{}h\allowbreak{}o\allowbreak{}r\allowbreak{}e\allowbreak{}d\allowbreak{}\_\allowbreak{}q\allowbreak{}u\allowbreak{}a\allowbreak{}r\allowbreak{}t\allowbreak{}i\allowbreak{}c\allowbreak{}\_\allowbreak{}t\allowbreak{}r\allowbreak{}a\allowbreak{}n\allowbreak{}s\allowbreak{}p\allowbreak{}o\allowbreak{}r\allowbreak{}t\allowbreak{}s\allowbreak{}.\allowbreak{}j\allowbreak{}s\allowbreak{}o\allowbreak{}n\allowbreak{}}
also records all 8,621 original face multisets, their class numbers, and
an explicit signed coordinate permutation plus translation for each. The
original edge orientation and inverse-link rule are implemented by
{\ttfamily t\allowbreak{}r\allowbreak{}a\allowbreak{}n\allowbreak{}s\allowbreak{}f\allowbreak{}o\allowbreak{}r\allowbreak{}m\allowbreak{}\_\allowbreak{}e\allowbreak{}d\allowbreak{}g\allowbreak{}e\allowbreak{}}
in
{\ttfamily c\allowbreak{}a\allowbreak{}l\allowbreak{}c\allowbreak{}u\allowbreak{}l\allowbreak{}a\allowbreak{}t\allowbreak{}i\allowbreak{}o\allowbreak{}n\allowbreak{}s\allowbreak{}/\allowbreak{}r\allowbreak{}e\allowbreak{}c\allowbreak{}o\allowbreak{}n\allowbreak{}s\allowbreak{}t\allowbreak{}r\allowbreak{}u\allowbreak{}c\allowbreak{}t\allowbreak{}.\allowbreak{}p\allowbreak{}y\allowbreak{}}.
This is a complete finite, coordinate-level specification of v4, rather
than an unevaluated B term.

\subsection{6. A readable closed component and the signed
derivative}\label{a-readable-closed-component-and-the-signed-derivative}

For a single face the fully evaluated contribution is

\[
 \boxed{v_4^{(p^4)}=
   \frac{17}{10368}\chi_1(\Omega_p)
   -\frac7{51840}\chi_2(\Omega_p).}
\tag{Q15}
\] The own-face part of the already available b22 is
{\ttfamily c\allowbreak{}h\allowbreak{}i\allowbreak{}1\allowbreak{}/\allowbreak{}1\allowbreak{}0\allowbreak{}3\allowbreak{}6\allowbreak{}8\allowbreak{}-\allowbreak{}c\allowbreak{}h\allowbreak{}i\allowbreak{}2\allowbreak{}/\allowbreak{}3\allowbreak{}1\allowbreak{}1\allowbreak{}0\allowbreak{}4\allowbreak{}}.
The other incoming term in Q4 has coefficients
{\ttfamily c\allowbreak{}h\allowbreak{}i\allowbreak{}1\allowbreak{}/\allowbreak{}6\allowbreak{}4\allowbreak{}8\allowbreak{}-\allowbreak{}c\allowbreak{}h\allowbreak{}i\allowbreak{}2\allowbreak{}/\allowbreak{}9\allowbreak{}7\allowbreak{}2\allowbreak{}0\allowbreak{}}.
Adding them gives Q15 exactly. The signs remain in the catalogue before
any norms are taken.

The catalogue also determines every formal plaquette-direction response
through degree three by the exact derivative

\[
 Z_p^{[3]}=\partial_{\lambda_p}
       \sum_{1\le|\nu|\le4}\lambda^\nu v_\nu.
\tag{Q16}
\] Differentiating Q13 coefficientwise proves that the terms through
total degree three in
{\ttfamily [\allowbreak{}K\allowbreak{}-\allowbreak{}2\allowbreak{}Q\allowbreak{}\_\allowbreak{}H\allowbreak{} \allowbreak{}G\allowbreak{}a\allowbreak{}m\allowbreak{}m\allowbreak{}a\allowbreak{}(\allowbreak{}q\allowbreak{}3\allowbreak{},\allowbreak{}.\allowbreak{})\allowbreak{}]\allowbreak{} \allowbreak{}Z\allowbreak{}\_\allowbreak{}p\allowbreak{}\textasciicircum{}\allowbreak{}[\allowbreak{}3\allowbreak{}]\allowbreak{} \allowbreak{}-\allowbreak{} \allowbreak{}W\allowbreak{}\_\allowbreak{}p\allowbreak{}}
vanish, with
{\ttfamily q\allowbreak{}3\allowbreak{}=\allowbreak{}s\allowbreak{}u\allowbreak{}m\allowbreak{}\_\allowbreak{}(\allowbreak{}1\allowbreak{}\textless{}\allowbreak{}=\allowbreak{}|\allowbreak{}n\allowbreak{}u\allowbreak{}|\allowbreak{}\textless{}\allowbreak{}=\allowbreak{}3\allowbreak{})\allowbreak{} \allowbreak{}l\allowbreak{}a\allowbreak{}m\allowbreak{}b\allowbreak{}d\allowbreak{}a\allowbreak{}\textasciicircum{}\allowbreak{}n\allowbreak{}u\allowbreak{} \allowbreak{}v\allowbreak{}\_\allowbreak{}n\allowbreak{}u\allowbreak{}}.
The remaining degree4--6 polynomial is exactly the corresponding sum of
those unremoved higher products. This records the signed tangent data;
it is not being promoted to a uniform bound on the full linearized
inverse.

The previous numerical coupling threshold requires a coefficient
trace-norm bound and a full infinite-tail argument. This session does
not replace those with the catalogue's Haar norms or assert a new
threshold.

\subsection{7. Evidence scope}\label{evidence-scope}

Every returned coefficient has zero Haar mean, an exact inverse
residual, and an independently expanded exact logarithmic-source
residual. The sum over all 78 coefficient downsets executes 1,053
inverse/source steps. The additional second-order-jet audit
differentiates original unfixed links and verifies the tree return and
original vertex-gauge invariance at retained rational assignments for
every representative. Those assignments are an independent regression;
the polynomial equalities and operator maps above are the identity
argument. No independent external or Lean review is claimed. The full
raw input, recovered code, calculations and receipts are included in
this session's ZIP.

\section{Completed sixth-order ground-energy calculation, including the
cube
surface}\label{completed-sixth-order-ground-energy-calculation-including-the-cube-surface}

Audit and continuation, 17 September 2026. This result is derived in
this session. The uploaded status line about sixth-order coefficients
contained no such coefficient table, proof, or execution record. The
formulas below retain the original Hamiltonian and open-box geometry.
They are finite-order coefficients with a separately proved
finite-volume remainder, not a new uniform-coupling or continuum
mass-gap theorem.

\subsection{1. Original parameters and exact coefficient
convention}\label{original-parameters-and-exact-coefficient-convention}

Use Q1 of
{\ttfamily F\allowbreak{}O\allowbreak{}U\allowbreak{}R\allowbreak{}T\allowbreak{}H\allowbreak{}\_\allowbreak{}O\allowbreak{}R\allowbreak{}D\allowbreak{}E\allowbreak{}R\allowbreak{}\_\allowbreak{}S\allowbreak{}O\allowbreak{}U\allowbreak{}R\allowbreak{}C\allowbreak{}E\allowbreak{}.\allowbreak{}m\allowbreak{}d\allowbreak{}}:

\[
 H=\kappa\{K+\xi(2M-S)\},\qquad S=\sum_p W_p,
 \quad\kappa=2g^2/a,\quad\xi=1/(4g^4).
\tag{E1}
\] Write the ground eigenvalue of
{\ttfamily K\allowbreak{}-\allowbreak{}x\allowbreak{}i\allowbreak{} \allowbreak{}S\allowbreak{}}
as
{\ttfamily e\allowbreak{}(\allowbreak{}x\allowbreak{}i\allowbreak{})\allowbreak{}=\allowbreak{}e\allowbreak{}2\allowbreak{} \allowbreak{}x\allowbreak{}i\allowbreak{}\textasciicircum{}\allowbreak{}2\allowbreak{}+\allowbreak{}e\allowbreak{}4\allowbreak{} \allowbreak{}x\allowbreak{}i\allowbreak{}\textasciicircum{}\allowbreak{}4\allowbreak{}+\allowbreak{}e\allowbreak{}6\allowbreak{} \allowbreak{}x\allowbreak{}i\allowbreak{}\textasciicircum{}\allowbreak{}6\allowbreak{}+\allowbreak{}.\allowbreak{}.\allowbreak{}.\allowbreak{}}.
Then the original physical energy is
{\ttfamily E\allowbreak{}0\allowbreak{}=\allowbreak{}k\allowbreak{}a\allowbreak{}p\allowbreak{}p\allowbreak{}a\allowbreak{}(\allowbreak{}2\allowbreak{}M\allowbreak{} \allowbreak{}x\allowbreak{}i\allowbreak{}+\allowbreak{}e\allowbreak{}(\allowbreak{}x\allowbreak{}i\allowbreak{})\allowbreak{})\allowbreak{}}.
The additive original scalar stays in this equation. A link-center
transformation
{\ttfamily U\allowbreak{}\_\allowbreak{}(\allowbreak{}n\allowbreak{},\allowbreak{}i\allowbreak{})\allowbreak{} \allowbreak{}-\allowbreak{}\textgreater{}\allowbreak{} \allowbreak{}(\allowbreak{}-\allowbreak{}1\allowbreak{})\allowbreak{}\textasciicircum{}\allowbreak{}(\allowbreak{}s\allowbreak{}u\allowbreak{}m\allowbreak{}\_\allowbreak{}(\allowbreak{}j\allowbreak{}\textless{}\allowbreak{}i\allowbreak{})\allowbreak{}n\allowbreak{}\_\allowbreak{}j\allowbreak{})\allowbreak{} \allowbreak{}U\allowbreak{}\_\allowbreak{}(\allowbreak{}n\allowbreak{},\allowbreak{}i\allowbreak{})\allowbreak{}}
preserves Haar and K and changes every W\_p to -W\_p.~It intertwines
{\ttfamily K\allowbreak{}-\allowbreak{}x\allowbreak{}i\allowbreak{} \allowbreak{}S\allowbreak{}}
with
{\ttfamily K\allowbreak{}+\allowbreak{}x\allowbreak{}i\allowbreak{} \allowbreak{}S\allowbreak{}};
the ground branch near zero is simple, so e is even.

Use the formal eigenline representative
{\ttfamily u\allowbreak{}=\allowbreak{}1\allowbreak{}+\allowbreak{}x\allowbreak{}i\allowbreak{} \allowbreak{}u\allowbreak{}1\allowbreak{}+\allowbreak{}x\allowbreak{}i\allowbreak{}\textasciicircum{}\allowbreak{}2\allowbreak{} \allowbreak{}u\allowbreak{}2\allowbreak{}+\allowbreak{}.\allowbreak{}.\allowbreak{}.\allowbreak{}},
{\ttfamily P\allowbreak{}\_\allowbreak{}H\allowbreak{} \allowbreak{}u\allowbreak{}n\allowbreak{}=\allowbreak{}0\allowbreak{}}.
The exact scalar return to the original unit vacuum is Q2 and Q12 of the
companion. The low-order equations are

\[
 u_1=K^{-1}S=S/3,\qquad
 u_2=K^{-1}Q_H(Su_1),\qquad
 u_3=K^{-1}Q_H(Su_2+e_2u_1),
\] \[
 e_2=-\langle S,u_1\rangle_H=-M/3,\qquad
 e_4=-\langle S,u_3\rangle_H=5M/216-2J/1053.
\tag{E2}
\] Here J counts original unordered pairs of plaquettes sharing an edge.
All pairings are original Haar pairings.

The complete sixth coefficient depends only on these three actual
wavefunction coefficients:

\[
 \boxed{e_6=-\langle u_3,Ku_3\rangle_H
             -e_4\|u_1\|_H^2-e_2\|u_2\|_H^2.}
\tag{E3}
\] For an explicit proof, use
{\ttfamily e\allowbreak{}6\allowbreak{}=\allowbreak{}-\allowbreak{}\textless{}\allowbreak{}S\allowbreak{},\allowbreak{}u\allowbreak{}5\allowbreak{}\textgreater{}\allowbreak{}=\allowbreak{}-\allowbreak{}\textless{}\allowbreak{}u\allowbreak{}1\allowbreak{},\allowbreak{}K\allowbreak{}u\allowbreak{}5\allowbreak{}\textgreater{}\allowbreak{}}.
The order-five equation is
{\ttfamily K\allowbreak{}u\allowbreak{}5\allowbreak{}=\allowbreak{}Q\allowbreak{}\_\allowbreak{}H\allowbreak{} \allowbreak{}S\allowbreak{} \allowbreak{}u\allowbreak{}4\allowbreak{}+\allowbreak{}e\allowbreak{}2\allowbreak{} \allowbreak{}u\allowbreak{}3\allowbreak{}+\allowbreak{}e\allowbreak{}4\allowbreak{} \allowbreak{}u\allowbreak{}1\allowbreak{}}.
Move K through
{\ttfamily \textless{}\allowbreak{}S\allowbreak{} \allowbreak{}u\allowbreak{}1\allowbreak{},\allowbreak{}u\allowbreak{}4\allowbreak{}\textgreater{}\allowbreak{}}
by
{\ttfamily K\allowbreak{}u\allowbreak{}2\allowbreak{}=\allowbreak{}S\allowbreak{}u\allowbreak{}1\allowbreak{}+\allowbreak{}e\allowbreak{}2\allowbreak{}},
and then use
{\ttfamily K\allowbreak{}u\allowbreak{}4\allowbreak{}=\allowbreak{}Q\allowbreak{}\_\allowbreak{}H\allowbreak{} \allowbreak{}S\allowbreak{} \allowbreak{}u\allowbreak{}3\allowbreak{}+\allowbreak{}e\allowbreak{}2\allowbreak{} \allowbreak{}u\allowbreak{}2\allowbreak{}+\allowbreak{}e\allowbreak{}4\allowbreak{}}.
Finally substitute
{\ttfamily S\allowbreak{}u\allowbreak{}2\allowbreak{}=\allowbreak{}K\allowbreak{}u\allowbreak{}3\allowbreak{}-\allowbreak{}e\allowbreak{}2\allowbreak{} \allowbreak{}u\allowbreak{}1\allowbreak{}}.
The two cross terms
{\ttfamily e\allowbreak{}2\allowbreak{}\textless{}\allowbreak{}u\allowbreak{}1\allowbreak{},\allowbreak{}u\allowbreak{}3\allowbreak{}\textgreater{}\allowbreak{}}
cancel; the remaining terms are precisely E3. This also shows why the
two state-denominator corrections in E3 must be retained.

\subsection{2. Complete third-vector spectral
data}\label{complete-third-vector-spectral-data}

For one face the spin-3/2 character component of u3 is
{\ttfamily c\allowbreak{}h\allowbreak{}i\allowbreak{}\_\allowbreak{}(\allowbreak{}3\allowbreak{}/\allowbreak{}2\allowbreak{})\allowbreak{}(\allowbreak{}O\allowbreak{}m\allowbreak{}e\allowbreak{}g\allowbreak{}a\allowbreak{}\_\allowbreak{}p\allowbreak{})\allowbreak{}/\allowbreak{}3\allowbreak{}6\allowbreak{}0\allowbreak{}}.
Its kinetic norm contribution is 1/8640. The coefficient of the original
W\_p in the assembled u3 is

\[
 a_p=-5/216+d_p/1053,
\tag{E4}
\] where d\_p is its original plaquette adjacency degree. This is
obtained by multiplying the self and adjacent components of u2 by each
possible outer face in
{\ttfamily S\allowbreak{} \allowbreak{}u\allowbreak{}2\allowbreak{}+\allowbreak{}e\allowbreak{}2\allowbreak{} \allowbreak{}u\allowbreak{}1\allowbreak{}}.
It is also the coefficient required by
{\ttfamily e\allowbreak{}4\allowbreak{}=\allowbreak{}-\allowbreak{}s\allowbreak{}u\allowbreak{}m\allowbreak{}\_\allowbreak{}p\allowbreak{} \allowbreak{}a\allowbreak{}\_\allowbreak{}p\allowbreak{}}
and
{\ttfamily s\allowbreak{}u\allowbreak{}m\allowbreak{}\_\allowbreak{}p\allowbreak{} \allowbreak{}d\allowbreak{}\_\allowbreak{}p\allowbreak{}=\allowbreak{}2\allowbreak{}J\allowbreak{}}.

For the repeated multiset p,p,q on adjacent faces, put
{\ttfamily F\allowbreak{}=\allowbreak{}(\allowbreak{}W\allowbreak{}\_\allowbreak{}p\allowbreak{}\textasciicircum{}\allowbreak{}2\allowbreak{}-\allowbreak{}1\allowbreak{})\allowbreak{}W\allowbreak{}\_\allowbreak{}q\allowbreak{}}.
The two actual shared-edge spins are 1/2 and 3/2, with Casimirs 9 and
12. The coefficients in u3 are

\[
 r_{1/2}=\frac{1/24+1/9+1/39}{9}=\frac{167}{8424},
 \qquad
 r_{3/2}=\frac{1/24+4/39}{12}=\frac5{416}.
\tag{E5}
\] Their Haar squared masses are 1/3 and 2/3. To verify the lower mass,
write the original shared unit quaternion u and complementary unit
quaternions v,w. The harmonic projection is
{\ttfamily H\allowbreak{}=\allowbreak{}(\allowbreak{}8\allowbreak{}/\allowbreak{}3\allowbreak{})\allowbreak{}(\allowbreak{}u\allowbreak{} \allowbreak{}d\allowbreak{}o\allowbreak{}t\allowbreak{} \allowbreak{}v\allowbreak{})\allowbreak{}(\allowbreak{}v\allowbreak{} \allowbreak{}d\allowbreak{}o\allowbreak{}t\allowbreak{} \allowbreak{}w\allowbreak{})\allowbreak{}-\allowbreak{}(\allowbreak{}2\allowbreak{}/\allowbreak{}3\allowbreak{})\allowbreak{}(\allowbreak{}u\allowbreak{} \allowbreak{}d\allowbreak{}o\allowbreak{}t\allowbreak{} \allowbreak{}w\allowbreak{})\allowbreak{}}.
Using
{\ttfamily i\allowbreak{}n\allowbreak{}t\allowbreak{} \allowbreak{}u\allowbreak{}\_\allowbreak{}i\allowbreak{} \allowbreak{}u\allowbreak{}\_\allowbreak{}j\allowbreak{}=\allowbreak{}d\allowbreak{}e\allowbreak{}l\allowbreak{}t\allowbreak{}a\allowbreak{}\_\allowbreak{}i\allowbreak{}j\allowbreak{}/\allowbreak{}4\allowbreak{}}
and independent Haar in v,w gives
{\ttfamily \textless{}\allowbreak{}H\allowbreak{}\textasciicircum{}\allowbreak{}2\allowbreak{}\textgreater{}\allowbreak{}=\allowbreak{}4\allowbreak{}/\allowbreak{}9\allowbreak{}+\allowbreak{}1\allowbreak{}/\allowbreak{}9\allowbreak{}-\allowbreak{}2\allowbreak{}/\allowbreak{}9\allowbreak{}=\allowbreak{}1\allowbreak{}/\allowbreak{}3\allowbreak{}}.
The full
{\ttfamily \textless{}\allowbreak{}F\allowbreak{}\textasciicircum{}\allowbreak{}2\allowbreak{}\textgreater{}\allowbreak{}=\allowbreak{}1\allowbreak{}},
giving the other mass. Thus each oriented repeated-adjacent contribution
to
{\ttfamily \textless{}\allowbreak{}u\allowbreak{}3\allowbreak{},\allowbreak{}K\allowbreak{}u\allowbreak{}3\allowbreak{}\textgreater{}\allowbreak{}}
is

\[
 R_{\mathrm{adj}}=3r_{1/2}^{2}+8r_{3/2}^{2}
                 =110453/47309184.
\] A repeated nonadjacent pair has coefficient
{\ttfamily 1\allowbreak{}/\allowbreak{}7\allowbreak{}2\allowbreak{}},
Casimir 11, Haar squared norm one, hence contribution
{\ttfamily R\allowbreak{}\_\allowbreak{}d\allowbreak{}i\allowbreak{}s\allowbreak{}=\allowbreak{}1\allowbreak{}1\allowbreak{}/\allowbreak{}5\allowbreak{}1\allowbreak{}8\allowbreak{}4\allowbreak{}}.

For a distinct pair in u2 the adjacent coefficients are
{\ttfamily u\allowbreak{}\_\allowbreak{}0\allowbreak{}=\allowbreak{}4\allowbreak{}/\allowbreak{}2\allowbreak{}7\allowbreak{}},
{\ttfamily u\allowbreak{}\_\allowbreak{}1\allowbreak{}=\allowbreak{}4\allowbreak{}/\allowbreak{}3\allowbreak{}9\allowbreak{}},
with shared-edge masses 1/4 and 3/4. The adjacent-pair squared norm in
u2 is therefore

\[
 P_2=\tfrac14(4/27)^2+\tfrac34(4/39)^2=1648/123201.
\tag{E6}
\] The nonadjacent-pair value is 1/81, and the self value is 1/576.

For three distinct faces the original projection masses and Casimirs are
as follows. Each coefficient is the sum of its three possible outer-face
sources from E2, divided by that actual Casimir.

\begin{itemize}
\tightlist
\item
  No adjacency: coefficient 1/27, Casimir 9, squared mass one.
\item
  One adjacent pair:
  {\ttfamily u\allowbreak{}\_\allowbreak{}s\allowbreak{}/\allowbreak{}3\allowbreak{}},
  Casimir
  {\ttfamily 1\allowbreak{}5\allowbreak{}/\allowbreak{}2\allowbreak{}+\allowbreak{}2\allowbreak{}s\allowbreak{}},
  masses 1/4,3/4.
\item
  Two-edge path: coefficient
  {\ttfamily (\allowbreak{}1\allowbreak{}/\allowbreak{}9\allowbreak{}+\allowbreak{}u\allowbreak{}\_\allowbreak{}s\allowbreak{}+\allowbreak{}u\allowbreak{}\_\allowbreak{}t\allowbreak{})\allowbreak{}/\allowbreak{}(\allowbreak{}6\allowbreak{}+\allowbreak{}2\allowbreak{}s\allowbreak{}+\allowbreak{}2\allowbreak{}t\allowbreak{})\allowbreak{}},
  with masses
  {\ttfamily (\allowbreak{}1\allowbreak{},\allowbreak{}3\allowbreak{},\allowbreak{}3\allowbreak{},\allowbreak{}9\allowbreak{})\allowbreak{}/\allowbreak{}1\allowbreak{}6\allowbreak{}}
  for shared-edge spins
  {\ttfamily (\allowbreak{}0\allowbreak{},\allowbreak{}0\allowbreak{})\allowbreak{},\allowbreak{}(\allowbreak{}0\allowbreak{},\allowbreak{}1\allowbreak{})\allowbreak{},\allowbreak{}(\allowbreak{}1\allowbreak{},\allowbreak{}0\allowbreak{})\allowbreak{},\allowbreak{}(\allowbreak{}1\allowbreak{},\allowbreak{}1\allowbreak{})\allowbreak{}}.
\item
  Common-edge triple: coefficients
  {\ttfamily (\allowbreak{}3\allowbreak{}/\allowbreak{}2\allowbreak{})\allowbreak{}(\allowbreak{}u\allowbreak{}\_\allowbreak{}0\allowbreak{}+\allowbreak{}u\allowbreak{}\_\allowbreak{}1\allowbreak{})\allowbreak{}/\allowbreak{}(\allowbreak{}1\allowbreak{}5\allowbreak{}/\allowbreak{}2\allowbreak{})\allowbreak{}}
  and
  {\ttfamily 3\allowbreak{}u\allowbreak{}\_\allowbreak{}1\allowbreak{}/\allowbreak{}(\allowbreak{}2\allowbreak{}1\allowbreak{}/\allowbreak{}2\allowbreak{})\allowbreak{}};
  squared masses 1/2,1/2.
\item
  Cube corner: coefficient
  {\ttfamily [\allowbreak{}(\allowbreak{}3\allowbreak{}-\allowbreak{}n\allowbreak{})\allowbreak{}u\allowbreak{}\_\allowbreak{}0\allowbreak{}+\allowbreak{}n\allowbreak{} \allowbreak{}u\allowbreak{}\_\allowbreak{}1\allowbreak{}]\allowbreak{}/\allowbreak{}(\allowbreak{}9\allowbreak{}/\allowbreak{}2\allowbreak{}+\allowbreak{}2\allowbreak{}n\allowbreak{})\allowbreak{}};
  the n=0,2,3 squared masses, with n=2 aggregated, are
  {\ttfamily 1\allowbreak{}/\allowbreak{}1\allowbreak{}6\allowbreak{},\allowbreak{}9\allowbreak{}/\allowbreak{}1\allowbreak{}6\allowbreak{},\allowbreak{}3\allowbreak{}/\allowbreak{}8\allowbreak{}}.
  The n=1 subspace is killed by the original vertex invariant
  projection.
\end{itemize}

Here is a derivation of all nontrivial masses. For a path, one
shared-edge Haar projection is one half of the original two-face
boundary times the third trace, with squared norm 1/4. Both such
projections give one quarter of the eight-link boundary trace, with
squared norm 1/16. Orthogonality of the commuting original edge
projections gives the remaining three masses. For a common edge, expose
three independent complementary unit quaternions and the shared original
unit quaternion. Each pair singlet term has squared norm 1/4, and each
distinct cross pairing is 1/16. The identity
{\ttfamily P\allowbreak{}\_\allowbreak{}(\allowbreak{}1\allowbreak{}/\allowbreak{}2\allowbreak{})\allowbreak{}F\allowbreak{}=\allowbreak{}(\allowbreak{}2\allowbreak{}/\allowbreak{}3\allowbreak{})\allowbreak{}s\allowbreak{}u\allowbreak{}m\allowbreak{} \allowbreak{}W\allowbreak{}\_\allowbreak{}r\allowbreak{} \allowbreak{}P\allowbreak{}\_\allowbreak{}0\allowbreak{}(\allowbreak{}W\allowbreak{}\_\allowbreak{}p\allowbreak{} \allowbreak{}W\allowbreak{}\_\allowbreak{}q\allowbreak{})\allowbreak{}}
gives mass
{\ttfamily (\allowbreak{}4\allowbreak{}/\allowbreak{}9\allowbreak{})\allowbreak{}(\allowbreak{}3\allowbreak{}/\allowbreak{}4\allowbreak{}+\allowbreak{}6\allowbreak{}/\allowbreak{}1\allowbreak{}6\allowbreak{})\allowbreak{}=\allowbreak{}1\allowbreak{}/\allowbreak{}2\allowbreak{}};
its complement has the other half. For a corner, the triple-singlet
projection is one quarter of the six-link boundary, of squared mass
1/16. Each single-edge-singlet projection has squared mass 1/4. Its only
additional component is one of the three n=2 channels, of mass 3/16. The
full mass is one, so the n=3 remainder is 3/8. All simple-loop Haar
squared norms here equal one by integrating one of its original links.
Unique exterior links justify the stated product masses.

These give the following exact kinetic contributions for each distinct
triple in u3:

{\def\LTcaptype{none} % do not increment counter
\begin{longtable}[]{@{}lr@{}}
\toprule\noalign{}
Adjacency type & Contribution to
{\ttfamily \textless{}\allowbreak{}u\allowbreak{}3\allowbreak{},\allowbreak{}K\allowbreak{} \allowbreak{}u\allowbreak{}3\allowbreak{}\textgreater{}\allowbreak{}} \\
\midrule\noalign{}
\endhead
\bottomrule\noalign{}
\endlastfoot
none & 1/81 \\
one pair & 4768/369603 \\
path & 1595629/118272960 \\
common edge & 20032/1437345 \\
cube corner & 210880/14660919 \\
\end{longtable}
}

The direct original-generator calculation in
{\ttfamily g\allowbreak{}a\allowbreak{}u\allowbreak{}g\allowbreak{}e\allowbreak{}\_\allowbreak{}p\allowbreak{}o\allowbreak{}l\allowbreak{}y\allowbreak{}n\allowbreak{}o\allowbreak{}m\allowbreak{}i\allowbreak{}a\allowbreak{}l\allowbreak{}.\allowbreak{}p\allowbreak{}y\allowbreak{}}
constructs u1,u2,u3 without this table, using the exact K fields and
inverse residuals. It agrees with E3 and the table for every geometric
orbit of a connected one-, two-, or three-face set. That calculation
also verifies the complete Rayleigh numerator and denominator identity,
preserving all cross terms.

\subsection{3. Linked contributions and why six distinct cube faces
occur}\label{linked-contributions-and-why-six-distinct-cube-faces-occur}

With independent formal face couplings, an energy coefficient with
support split into edge-disjoint components is additive in those
components. The full Haar space, electric operator and potential factor
over the disjoint edge sets; their positive vacuum product is already
gauge invariant at any shared vertices. Thus mixed connected-energy
coefficients vanish on an edge-disconnected set. Subtracting proper
subcluster energies gives the following linked sixth-order weights:

{\def\LTcaptype{none} % do not increment counter
\begin{longtable}[]{@{}
  >{\raggedright\arraybackslash}p{(\linewidth - 2\tabcolsep) * \real{0.4286}}
  >{\raggedleft\arraybackslash}p{(\linewidth - 2\tabcolsep) * \real{0.5714}}@{}}
\toprule\noalign{}
\begin{minipage}[b]{\linewidth}\raggedright
Support
\end{minipage} & \begin{minipage}[b]{\linewidth}\raggedleft
Weight
\end{minipage} \\
\midrule\noalign{}
\endhead
\bottomrule\noalign{}
\endlastfoot
One plaquette &
{\ttfamily -\allowbreak{}2\allowbreak{}8\allowbreak{}9\allowbreak{}/\allowbreak{}7\allowbreak{}7\allowbreak{}7\allowbreak{}6\allowbreak{}0\allowbreak{}} \\
Adjacent pair, with all sixth-order powers on that support &
{\ttfamily 2\allowbreak{}2\allowbreak{}2\allowbreak{}8\allowbreak{}5\allowbreak{}/\allowbreak{}2\allowbreak{}3\allowbreak{}6\allowbreak{}5\allowbreak{}4\allowbreak{}5\allowbreak{}9\allowbreak{}2\allowbreak{}} \\
Distinct path of three plaquettes &
{\ttfamily -\allowbreak{}4\allowbreak{}9\allowbreak{}0\allowbreak{}9\allowbreak{}/\allowbreak{}1\allowbreak{}1\allowbreak{}8\allowbreak{}2\allowbreak{}7\allowbreak{}2\allowbreak{}9\allowbreak{}6\allowbreak{}0\allowbreak{}} \\
Three plaquettes sharing an edge &
{\ttfamily 2\allowbreak{}4\allowbreak{}4\allowbreak{}/\allowbreak{}4\allowbreak{}3\allowbreak{}1\allowbreak{}2\allowbreak{}0\allowbreak{}3\allowbreak{}5\allowbreak{}} \\
Three faces at a cube corner &
{\ttfamily -\allowbreak{}2\allowbreak{}1\allowbreak{}2\allowbreak{}/\allowbreak{}5\allowbreak{}4\allowbreak{}2\allowbreak{}9\allowbreak{}9\allowbreak{}7\allowbreak{}} \\
All six faces of an elementary cube, each once &
{\ttfamily -\allowbreak{}8\allowbreak{}3\allowbreak{}/\allowbreak{}1\allowbreak{}9\allowbreak{}4\allowbreak{}4\allowbreak{}} \\
\end{longtable}
}

For example, the original two-adjacent-face energy has
{\ttfamily e\allowbreak{}6\allowbreak{}=\allowbreak{}-\allowbreak{}7\allowbreak{}6\allowbreak{}7\allowbreak{}7\allowbreak{}1\allowbreak{}3\allowbreak{}/\allowbreak{}1\allowbreak{}1\allowbreak{}8\allowbreak{}2\allowbreak{}7\allowbreak{}2\allowbreak{}9\allowbreak{}6\allowbreak{}0\allowbreak{}}.
Subtracting twice the one-face value gives the pair weight. The original
path/common/corner three-face e6 values are
{\ttfamily -\allowbreak{}1\allowbreak{}8\allowbreak{}3\allowbreak{}4\allowbreak{}6\allowbreak{}1\allowbreak{}/\allowbreak{}1\allowbreak{}9\allowbreak{}7\allowbreak{}1\allowbreak{}2\allowbreak{}1\allowbreak{}6\allowbreak{}0\allowbreak{}},
{\ttfamily -\allowbreak{}5\allowbreak{}2\allowbreak{}8\allowbreak{}1\allowbreak{}/\allowbreak{}6\allowbreak{}3\allowbreak{}8\allowbreak{}8\allowbreak{}2\allowbreak{}0\allowbreak{}},
and
{\ttfamily -\allowbreak{}2\allowbreak{}5\allowbreak{}5\allowbreak{}5\allowbreak{}0\allowbreak{}5\allowbreak{}1\allowbreak{}/\allowbreak{}2\allowbreak{}9\allowbreak{}3\allowbreak{}2\allowbreak{}1\allowbreak{}8\allowbreak{}3\allowbreak{}8\allowbreak{}0\allowbreak{}}
respectively. Subtracting the three one-face values and the actual
number of adjacent pair weights gives the table. These are rational
identities, executed both from E3 and from the independently constructed
polynomial states.

The last row needs an explicit additional calculation. Let Z\_e multiply
one original link by -I. It commutes with K and preserves Haar. The
operator

\[
 \Pi_{\mathrm{even}}=\prod_e(I+Z_e)/2
\tag{E7}
\] projects the original polynomial source onto the edge-even subspace.
A monomial in face couplings can contribute to the scalar energy only if
its set of odd-multiplicity faces is a mod-two closed cubical surface.
For total order six, either every face has even multiplicity (supports
of at most three faces), or the six distinct faces form an elementary
cube surface.

For completeness, every finite cubical two-cycle in this open cubic box
bounds a finite mod-two set of cubes. This can be constructed by
assigning cube membership from the exterior along one coordinate: the
cycle condition ensures the assignments obtained from the other faces
agree. For a nonempty cube set, each occupied coordinate line has two
boundary faces. Two distinct cubes give at least five occupied
two-dimensional projection cells in the three directions, hence at least
ten boundary faces. A boundary with at most six faces therefore comes
from exactly one cube. This proves the stated exhaustion at order six.

\subsubsection{Exact cube calculation}\label{exact-cube-calculation}

Orient the six original face words outward. Complete face reversal
leaves an SU(2) fundamental trace unchanged. Each of the twelve original
edges is then traversed once in each direction. The identity
{\ttfamily i\allowbreak{}n\allowbreak{}t\allowbreak{} \allowbreak{}U\allowbreak{}\_\allowbreak{}a\allowbreak{}b\allowbreak{} \allowbreak{}c\allowbreak{}o\allowbreak{}n\allowbreak{}j\allowbreak{}u\allowbreak{}g\allowbreak{}a\allowbreak{}t\allowbreak{}e\allowbreak{}(\allowbreak{}U\allowbreak{}\_\allowbreak{}c\allowbreak{}d\allowbreak{})\allowbreak{}d\allowbreak{}U\allowbreak{}=\allowbreak{}d\allowbreak{}e\allowbreak{}l\allowbreak{}t\allowbreak{}a\allowbreak{}\_\allowbreak{}a\allowbreak{}c\allowbreak{} \allowbreak{}d\allowbreak{}e\allowbreak{}l\allowbreak{}t\allowbreak{}a\allowbreak{}\_\allowbreak{}b\allowbreak{}d\allowbreak{}/\allowbreak{}2\allowbreak{}}
supplies one factor 1/2 per edge. The remaining trace-index
identifications have eight independent vertex-color loops, each of
dimension two. Consequently

\[
 \left\langle\prod_{p\in\partial c}W_p\right\rangle_H
 =2^8/2^{12}=1/16.
\tag{E8}
\] The coordinate union-find audit retains all twelve edge gluings and
verifies the eight color classes.

For an order of the six distinct face insertions, let A\_k be its first
k faces. Every edge internal to A\_k has both its insertions already
present. Any nontrivial spin on that edge is killed by final Haar
expectation, because no later face contains it. Its singlet projector
commutes with the remaining insertions and the Casimir. Each boundary
edge carries its original spin 1/2, so the retained intermediate energy
is
{\ttfamily (\allowbreak{}3\allowbreak{}/\allowbreak{}4\allowbreak{})\allowbreak{}|\allowbreak{}p\allowbreak{}a\allowbreak{}r\allowbreak{}t\allowbreak{}i\allowbreak{}a\allowbreak{}l\allowbreak{} \allowbreak{}A\allowbreak{}\_\allowbreak{}k\allowbreak{}|\allowbreak{}}.
No proper A\_k is a closed surface, so no intermediate vacuum
subtraction occurs. Therefore the full linked coefficient is

\[
 -\frac1{16}\sum_{\pi\in S_6}
   \prod_{k=1}^{5}\frac4{3|\partial A_k(\pi)|}
 =-\frac1{16}\frac{166}{243}
 =\boxed{-83/1944}.
\tag{E9}
\] All 720 original orders and their five boundary counts are in
{\ttfamily r\allowbreak{}e\allowbreak{}s\allowbreak{}u\allowbreak{}l\allowbreak{}t\allowbreak{}s\allowbreak{}/\allowbreak{}c\allowbreak{}u\allowbreak{}b\allowbreak{}e\allowbreak{}\_\allowbreak{}s\allowbreak{}i\allowbreak{}x\allowbreak{}t\allowbreak{}h\allowbreak{}\_\allowbreak{}p\allowbreak{}a\allowbreak{}t\allowbreak{}h\allowbreak{}s\allowbreak{}.\allowbreak{}j\allowbreak{}s\allowbreak{}o\allowbreak{}n\allowbreak{}}.

There is also an independent three-plus-three calculation. Of the twenty
three-face cube subsets, twelve are paths and eight are corners. Their
all-internal-singlet u3 coefficients are 11/162 and 8/81, with boundary
energies 6 and 9/2. The cross pairing with the complementary triple is
E8. Thus the contribution to E3 is

\[
 -\frac1{16}\left[12\cdot6(11/162)^2
                   +8\cdot(9/2)(8/81)^2\right]=-83/1944.
\tag{E10}
\] The cube term is a required original three-dimensional contribution.
It cannot be assigned a planar value by deleting the extra faces.

\subsection{4. The complete finite-box
coefficient}\label{the-complete-finite-box-coefficient}

Let M count original plaquettes, J adjacent pairs, P3 three-face paths,
T\_e common-edge triples, T\_v cube-corner triples, and C cubes. The
result is

\[
 \boxed{\begin{aligned}
 e_6={}&-\frac{289}{77760}M
 +\frac{22285}{23654592}J
 -\frac{4909}{118272960}P_3\\
 &+\frac{244}{4312035}T_e
 -\frac{212}{542997}T_v
 -\frac{83}{1944}C.
 \end{aligned}}
\tag{E11}
\] This formula can also be verified before linked subtraction. The full
{\ttfamily |\allowbreak{}|\allowbreak{}u\allowbreak{}2\allowbreak{}|\allowbreak{}|\allowbreak{}\textasciicircum{}\allowbreak{}2\allowbreak{}}
is
{\ttfamily M\allowbreak{}/\allowbreak{}5\allowbreak{}7\allowbreak{}6\allowbreak{}+\allowbreak{}(\allowbreak{}c\allowbreak{}h\allowbreak{}o\allowbreak{}o\allowbreak{}s\allowbreak{}e\allowbreak{}(\allowbreak{}M\allowbreak{},\allowbreak{}2\allowbreak{})\allowbreak{}-\allowbreak{}J\allowbreak{})\allowbreak{}/\allowbreak{}8\allowbreak{}1\allowbreak{}+\allowbreak{}J\allowbreak{}*\allowbreak{}(\allowbreak{}1\allowbreak{}6\allowbreak{}4\allowbreak{}8\allowbreak{}/\allowbreak{}1\allowbreak{}2\allowbreak{}3\allowbreak{}2\allowbreak{}0\allowbreak{}1\allowbreak{})\allowbreak{}}.
The full
{\ttfamily \textless{}\allowbreak{}u\allowbreak{}3\allowbreak{},\allowbreak{}K\allowbreak{}u\allowbreak{}3\allowbreak{}\textgreater{}\allowbreak{}}
has:
{\ttfamily 3\allowbreak{} \allowbreak{}s\allowbreak{}u\allowbreak{}m\allowbreak{}\_\allowbreak{}p\allowbreak{} \allowbreak{}a\allowbreak{}\_\allowbreak{}p\allowbreak{}\textasciicircum{}\allowbreak{}2\allowbreak{}+\allowbreak{}M\allowbreak{}/\allowbreak{}8\allowbreak{}6\allowbreak{}4\allowbreak{}0\allowbreak{}};
the repeated contributions
{\ttfamily [\allowbreak{}M\allowbreak{}(\allowbreak{}M\allowbreak{}-\allowbreak{}1\allowbreak{})\allowbreak{}-\allowbreak{}2\allowbreak{}J\allowbreak{}]\allowbreak{}*\allowbreak{}(\allowbreak{}1\allowbreak{}1\allowbreak{}/\allowbreak{}5\allowbreak{}1\allowbreak{}8\allowbreak{}4\allowbreak{})\allowbreak{}+\allowbreak{}2\allowbreak{}J\allowbreak{}*\allowbreak{}(\allowbreak{}1\allowbreak{}1\allowbreak{}0\allowbreak{}4\allowbreak{}5\allowbreak{}3\allowbreak{}/\allowbreak{}4\allowbreak{}7\allowbreak{}3\allowbreak{}0\allowbreak{}9\allowbreak{}1\allowbreak{}8\allowbreak{}4\allowbreak{})\allowbreak{}};
the distinct-triple table in section2; and the positive norm cross term
{\ttfamily (\allowbreak{}8\allowbreak{}3\allowbreak{}/\allowbreak{}1\allowbreak{}9\allowbreak{}4\allowbreak{}4\allowbreak{})\allowbreak{}C\allowbreak{}}.
Inserting these into E3 cancels every disconnected term and gives E11.

For
{\ttfamily m\allowbreak{}=\allowbreak{}2\allowbreak{}L\allowbreak{}\textgreater{}\allowbreak{}=\allowbreak{}4\allowbreak{}},
the original open-box counts are

\[
 \begin{gathered}
 M=3m^2(m+1),\quad J=6m(3m^2-1),\quad C=m^3,\\
 T_e=12m^2(m-1),\quad T_v=8m^3,\\
 P_3=138m^3-126m^2-24m+12.
 \end{gathered}
\tag{E12}
\] Here
{\ttfamily T\allowbreak{}\_\allowbreak{}e\allowbreak{}=\allowbreak{}s\allowbreak{}u\allowbreak{}m\allowbreak{}\_\allowbreak{}e\allowbreak{} \allowbreak{}c\allowbreak{}h\allowbreak{}o\allowbreak{}o\allowbreak{}s\allowbreak{}e\allowbreak{}(\allowbreak{}r\allowbreak{}\_\allowbreak{}e\allowbreak{},\allowbreak{}3\allowbreak{})\allowbreak{}}
for the original incidence r\_e; each cube has eight different corner
triples. For paths, count the wedges centered at an original face, then
subtract three per adjacency triangle:
{\ttfamily P\allowbreak{}3\allowbreak{}=\allowbreak{}s\allowbreak{}u\allowbreak{}m\allowbreak{}\_\allowbreak{}p\allowbreak{} \allowbreak{}c\allowbreak{}h\allowbreak{}o\allowbreak{}o\allowbreak{}s\allowbreak{}e\allowbreak{}(\allowbreak{}d\allowbreak{}\_\allowbreak{}p\allowbreak{},\allowbreak{}2\allowbreak{})\allowbreak{}-\allowbreak{}3\allowbreak{}(\allowbreak{}T\allowbreak{}\_\allowbreak{}e\allowbreak{}+\allowbreak{}T\allowbreak{}\_\allowbreak{}v\allowbreak{})\allowbreak{}}.
In one orientation, faces in each boundary coordinate plane have degrees
{\ttfamily 8\allowbreak{}-\allowbreak{}b\allowbreak{}\_\allowbreak{}x\allowbreak{}-\allowbreak{}b\allowbreak{}\_\allowbreak{}y\allowbreak{}},
and those in an interior plane have
{\ttfamily 1\allowbreak{}2\allowbreak{}-\allowbreak{}b\allowbreak{}\_\allowbreak{}x\allowbreak{}-\allowbreak{}b\allowbreak{}\_\allowbreak{}y\allowbreak{}},
where
{\ttfamily b\allowbreak{}\_\allowbreak{}x\allowbreak{},\allowbreak{}b\allowbreak{}\_\allowbreak{}y\allowbreak{}}
indicate whether the face interval touches its corresponding boundary.
The numbers for
{\ttfamily b\allowbreak{}\_\allowbreak{}x\allowbreak{}+\allowbreak{}b\allowbreak{}\_\allowbreak{}y\allowbreak{}=\allowbreak{}0\allowbreak{},\allowbreak{}1\allowbreak{},\allowbreak{}2\allowbreak{}}
are
{\ttfamily (\allowbreak{}m\allowbreak{}-\allowbreak{}2\allowbreak{})\allowbreak{}\textasciicircum{}\allowbreak{}2\allowbreak{},\allowbreak{}4\allowbreak{}(\allowbreak{}m\allowbreak{}-\allowbreak{}2\allowbreak{})\allowbreak{},\allowbreak{}4\allowbreak{}}.
There are two boundary and m-1 interior planes and three orientations.
Their exact sum is
{\ttfamily 1\allowbreak{}9\allowbreak{}8\allowbreak{}m\allowbreak{}\textasciicircum{}\allowbreak{}3\allowbreak{}-\allowbreak{}1\allowbreak{}6\allowbreak{}2\allowbreak{}m\allowbreak{}\textasciicircum{}\allowbreak{}2\allowbreak{}-\allowbreak{}2\allowbreak{}4\allowbreak{}m\allowbreak{}+\allowbreak{}1\allowbreak{}2\allowbreak{}},
giving E12. The executable also enumerates all original faces and
triples directly for L=2,3,4.

Substitution yields the closed finite-box coefficient

\[
 \boxed{e_6=
 -\frac{211396463m^3+30959193m^2+21845782m+2336684}
        {4691494080}.}
\tag{E13}
\] In particular, at L=2,

\[
 \boxed{\frac{E_0}{\kappa}
 =480\xi-80\xi^2+\frac{1198}{351}\xi^4
 -\frac{3528610133}{1172873520}\xi^6+R_8(\xi).}
\tag{E14}
\] The limit of this exact coefficient divided by M is
{\ttfamily -\allowbreak{}2\allowbreak{}1\allowbreak{}1\allowbreak{}3\allowbreak{}9\allowbreak{}6\allowbreak{}4\allowbreak{}6\allowbreak{}3\allowbreak{}/\allowbreak{}1\allowbreak{}4\allowbreak{}0\allowbreak{}7\allowbreak{}4\allowbreak{}4\allowbreak{}8\allowbreak{}2\allowbreak{}2\allowbreak{}4\allowbreak{}0\allowbreak{}}.
This is a limit of a finite-order coefficient; no interchange with an
infinite-volume perturbative sum is asserted.

\subsection{5. An independently proved finite-volume
remainder}\label{an-independently-proved-finite-volume-remainder}

The original physical free Casimir has lowest nonconstant eigenvalue 3.
A physical nonconstant spin support has no degree-one active vertex; a
finite graph of minimum active degree two contains a cycle, and the
original cubic graph has no cycle shorter than four. Each active spin
has Casimir at least 3/4. Four spin-one-half edges around one original
face attain 3. This proves the free value used here.

On the physical space,
{\ttfamily |\allowbreak{}|\allowbreak{}S\allowbreak{}|\allowbreak{}|\allowbreak{}\textless{}\allowbreak{}=\allowbreak{}2\allowbreak{}M\allowbreak{}}.
On the contour
{\ttfamily |\allowbreak{}z\allowbreak{}|\allowbreak{}=\allowbreak{}3\allowbreak{}/\allowbreak{}2\allowbreak{}},
the free resolvent
{\ttfamily (\allowbreak{}K\allowbreak{}-\allowbreak{}z\allowbreak{})\allowbreak{}\textasciicircum{}\allowbreak{}-\allowbreak{}1\allowbreak{}}
has norm at most 2/3. Set

\[
 R=\frac3{8M}.
\tag{E15}
\] For
{\ttfamily |\allowbreak{}x\allowbreak{}i\allowbreak{}|\allowbreak{}\textless{}\allowbreak{}=\allowbreak{}R\allowbreak{}},
the bounded perturbation
{\ttfamily -\allowbreak{}x\allowbreak{}i\allowbreak{} \allowbreak{}S\allowbreak{}}
has norm at most 3/4. The contour resolvent Neumann series is thus
bounded by a ratio at most 1/2, and its isolated spectral projection has
the same rank one as at zero. This gives an analytic ground branch
e(xi). The elementary resolvent inclusion
{\ttfamily d\allowbreak{}i\allowbreak{}s\allowbreak{}t\allowbreak{}(\allowbreak{}z\allowbreak{},\allowbreak{}s\allowbreak{}p\allowbreak{}e\allowbreak{}c\allowbreak{} \allowbreak{}K\allowbreak{})\allowbreak{}\textgreater{}\allowbreak{}|\allowbreak{}|\allowbreak{}x\allowbreak{}i\allowbreak{} \allowbreak{}S\allowbreak{}|\allowbreak{}|\allowbreak{} \allowbreak{}=\allowbreak{}\textgreater{}\allowbreak{} \allowbreak{}z\allowbreak{} \allowbreak{}i\allowbreak{}n\allowbreak{} \allowbreak{}r\allowbreak{}e\allowbreak{}s\allowbreak{}o\allowbreak{}l\allowbreak{}v\allowbreak{}e\allowbreak{}n\allowbreak{}t\allowbreak{}(\allowbreak{}K\allowbreak{}-\allowbreak{}x\allowbreak{}i\allowbreak{} \allowbreak{}S\allowbreak{})\allowbreak{}}
puts its single enclosed eigenvalue in
{\ttfamily |\allowbreak{}e\allowbreak{}(\allowbreak{}x\allowbreak{}i\allowbreak{})\allowbreak{}|\allowbreak{}\textless{}\allowbreak{}=\allowbreak{}3\allowbreak{}/\allowbreak{}4\allowbreak{}};
the other free eigenvalues are at least 3 away from zero. On the real
interval this branch is the original ground energy after subtracting the
explicit scalar 2Mxi.

Cauchy's estimate on the displayed radius and the exact parity in E1
give

\[
 \boxed{|R_8(\xi)|\le\frac34
 \frac{(|\xi|/R)^8}{1-(|\xi|/R)^2},\qquad |\xi|<R.}
\tag{E16}
\] The physical remainder is kappa times E16. For L=2,
{\ttfamily R\allowbreak{}=\allowbreak{}1\allowbreak{}/\allowbreak{}6\allowbreak{}4\allowbreak{}0\allowbreak{}}.
This complete remainder is volume dependent. It uses no unverified
uniform logarithmic-source bound and does not extend an earlier coupling
threshold or prove a continuum mass gap.

\subsection{6. Primary-literature check with explicit
conventions}\label{primary-literature-check-with-explicit-conventions}

For one original plaquette, K acts on class functions of
{\ttfamily O\allowbreak{}m\allowbreak{}e\allowbreak{}g\allowbreak{}a\allowbreak{}=\allowbreak{}c\allowbreak{}o\allowbreak{}s\allowbreak{}(\allowbreak{}t\allowbreak{}h\allowbreak{}e\allowbreak{}t\allowbreak{}a\allowbreak{})\allowbreak{}I\allowbreak{}-\allowbreak{}i\allowbreak{} \allowbreak{}s\allowbreak{}i\allowbreak{}n\allowbreak{}(\allowbreak{}t\allowbreak{}h\allowbreak{}e\allowbreak{}t\allowbreak{}a\allowbreak{})\allowbreak{}n\allowbreak{}.\allowbreak{}s\allowbreak{}i\allowbreak{}g\allowbreak{}m\allowbreak{}a\allowbreak{}}
as
{\ttfamily -\allowbreak{}(\allowbreak{}p\allowbreak{}a\allowbreak{}r\allowbreak{}t\allowbreak{}i\allowbreak{}a\allowbreak{}l\allowbreak{}\_\allowbreak{}t\allowbreak{}h\allowbreak{}e\allowbreak{}t\allowbreak{}a\allowbreak{}\textasciicircum{}\allowbreak{}2\allowbreak{}+\allowbreak{}2\allowbreak{}c\allowbreak{}o\allowbreak{}t\allowbreak{}(\allowbreak{}t\allowbreak{}h\allowbreak{}e\allowbreak{}t\allowbreak{}a\allowbreak{})\allowbreak{}p\allowbreak{}a\allowbreak{}r\allowbreak{}t\allowbreak{}i\allowbreak{}a\allowbreak{}l\allowbreak{}\_\allowbreak{}t\allowbreak{}h\allowbreak{}e\allowbreak{}t\allowbreak{}a\allowbreak{})\allowbreak{}}.
Put
{\ttfamily u\allowbreak{}(\allowbreak{}t\allowbreak{}h\allowbreak{}e\allowbreak{}t\allowbreak{}a\allowbreak{})\allowbreak{}=\allowbreak{}s\allowbreak{}i\allowbreak{}n\allowbreak{}(\allowbreak{}t\allowbreak{}h\allowbreak{}e\allowbreak{}t\allowbreak{}a\allowbreak{})\allowbreak{}f\allowbreak{}(\allowbreak{}t\allowbreak{}h\allowbreak{}e\allowbreak{}t\allowbreak{}a\allowbreak{})\allowbreak{}}
and
{\ttfamily t\allowbreak{}h\allowbreak{}e\allowbreak{}t\allowbreak{}a\allowbreak{}=\allowbreak{}2\allowbreak{}z\allowbreak{}}.
The equation for
{\ttfamily K\allowbreak{}-\allowbreak{}x\allowbreak{}i\allowbreak{} \allowbreak{}W\allowbreak{}}
becomes the Mathieu equation with
{\ttfamily a\allowbreak{}=\allowbreak{}4\allowbreak{}(\allowbreak{}e\allowbreak{}+\allowbreak{}1\allowbreak{})\allowbreak{}}
and
{\ttfamily q\allowbreak{}=\allowbreak{}-\allowbreak{}4\allowbreak{}x\allowbreak{}i\allowbreak{}},
with zero boundary values at z=0,pi/2. Thus the original ground branch
is

\[
 e_{\mathrm{one}}(\xi)=\tfrac14b_2(-4\xi)-1.
\tag{E17}
\] NIST DLMF 28.6.5 gives
{\ttfamily b\allowbreak{}2\allowbreak{}(\allowbreak{}q\allowbreak{})\allowbreak{}=\allowbreak{}4\allowbreak{}-\allowbreak{}q\allowbreak{}\textasciicircum{}\allowbreak{}2\allowbreak{}/\allowbreak{}1\allowbreak{}2\allowbreak{}+\allowbreak{}5\allowbreak{}q\allowbreak{}\textasciicircum{}\allowbreak{}4\allowbreak{}/\allowbreak{}1\allowbreak{}3\allowbreak{}8\allowbreak{}2\allowbreak{}4\allowbreak{}-\allowbreak{}2\allowbreak{}8\allowbreak{}9\allowbreak{}q\allowbreak{}\textasciicircum{}\allowbreak{}6\allowbreak{}/\allowbreak{}7\allowbreak{}9\allowbreak{}6\allowbreak{}2\allowbreak{}6\allowbreak{}2\allowbreak{}4\allowbreak{}0\allowbreak{}+\allowbreak{}.\allowbreak{}.\allowbreak{}.\allowbreak{}}.
The exact return in E17 is

\[
 e_{\mathrm{one}}=-\xi^2/3+5\xi^4/216
                         -289\xi^6/77760+\cdots,
\] matching the independent original character recurrence and the
single-face row above. A modern original-source treatment of the
single-plaquette Mathieu spectrum is Jakobs et al., EPJC85 (2025)1418,
equation39; its Hamiltonian convention must be transported before
numerical comparison. This literature check verifies the one-plaquette
convention and coefficient. It does not claim that the full cubic-box
E13 or the catalogue is new to the literature, nor supply an independent
review of those results.

Sources: https://dlmf.nist.gov/28.6.E5 ;
https://link.springer.com/article/10.1140/epjc/s10052-025-15120-x . For
the established exponential-vacuum and linked-source framework see
Schütte, Zheng Weihong and Hamer, arXiv:hep-lat/9603026, section3.

\section{Explicit signed tangent
certificate}\label{explicit-signed-tangent-certificate}

This accompanies Q16 in
{\ttfamily F\allowbreak{}O\allowbreak{}U\allowbreak{}R\allowbreak{}T\allowbreak{}H\allowbreak{}\_\allowbreak{}O\allowbreak{}R\allowbreak{}D\allowbreak{}E\allowbreak{}R\allowbreak{}\_\allowbreak{}S\allowbreak{}O\allowbreak{}U\allowbreak{}R\allowbreak{}C\allowbreak{}E\allowbreak{}.\allowbreak{}m\allowbreak{}d\allowbreak{}}.
The source convention is the same original family of independent
plaquette couplings, with
{\ttfamily l\allowbreak{}a\allowbreak{}m\allowbreak{}b\allowbreak{}d\allowbreak{}a\allowbreak{}\_\allowbreak{}p\allowbreak{}=\allowbreak{}x\allowbreak{}i\allowbreak{}}
only after the exact coefficients and their directional responses are
retained.

The main coefficient producer solves the original eigenvector
coefficient equation and then computes
{\ttfamily Q\allowbreak{}\_\allowbreak{}H\allowbreak{} \allowbreak{}l\allowbreak{}o\allowbreak{}g\allowbreak{} \allowbreak{}u\allowbreak{}}.
The independent program
{\ttfamily c\allowbreak{}h\allowbreak{}e\allowbreak{}c\allowbreak{}k\allowbreak{}s\allowbreak{}/\allowbreak{}v\allowbreak{}e\allowbreak{}r\allowbreak{}i\allowbreak{}f\allowbreak{}y\allowbreak{}\_\allowbreak{}s\allowbreak{}i\allowbreak{}g\allowbreak{}n\allowbreak{}e\allowbreak{}d\allowbreak{}\_\allowbreak{}t\allowbreak{}a\allowbreak{}n\allowbreak{}g\allowbreak{}e\allowbreak{}n\allowbreak{}t\allowbreak{}.\allowbreak{}p\allowbreak{}y\allowbreak{}}
instead solves the logarithmic-source equation directly on every
coefficient downset:

\begin{verbatim}
K v_nu = Q_H sum_(0<mu<nu) Gamma(v_mu,v_(nu-mu)),
K v_(e_p)=W_p.
\end{verbatim}

Its returned order-four coefficient equals the entire stored catalogue
polynomial for each of the 78 classes. The same polynomial ring,
original link vector fields, and exact inverse residual are specified in
Q5--Q13. The two coefficient constructions are distinct recurrences on
those fixed objects.

For the original face p and each retained source multiindex rho of total
degree at most three, define

\begin{verbatim}
Z_(p,rho)=(rho_p+1) v_(rho+e_p).
\end{verbatim}

The program evaluates the complete signed residual

\begin{verbatim}
K Z_(p,rho) - 2 Q_H sum_(0<mu<=rho)
              Gamma(v_mu,Z_(p,rho-mu))
\end{verbatim}

and proves it is the zero quotient polynomial for positive degree rho.
At rho=0 its value is exactly the original W\_p.~This is a finite
formal-parameter coefficient identity, with every sign and multiplicity
retained. It follows analytically by differentiating the
logarithmic-source equation; the checker also executes all coefficient
products explicitly.

The full receipt
{\ttfamily r\allowbreak{}e\allowbreak{}s\allowbreak{}u\allowbreak{}l\allowbreak{}t\allowbreak{}s\allowbreak{}/\allowbreak{}s\allowbreak{}i\allowbreak{}g\allowbreak{}n\allowbreak{}e\allowbreak{}d\allowbreak{}\_\allowbreak{}t\allowbreak{}a\allowbreak{}n\allowbreak{}g\allowbreak{}e\allowbreak{}n\allowbreak{}t\allowbreak{}.\allowbreak{}j\allowbreak{}s\allowbreak{}o\allowbreak{}n\allowbreak{}}
contains:

\begin{itemize}
\tightlist
\item
  78 independent reconstructed logarithmic sources;
\item
  2,156 exactly zero polynomial tangent residuals;
\item
  282 explicitly expanded degree-three response polynomials, with the
  original differentiated face, remaining coupling multiindex, support
  multiset and chord coordinates for every one.
\end{itemize}

The original symmetry transports in
{\ttfamily a\allowbreak{}n\allowbreak{}c\allowbreak{}h\allowbreak{}o\allowbreak{}r\allowbreak{}e\allowbreak{}d\allowbreak{}\_\allowbreak{}q\allowbreak{}u\allowbreak{}a\allowbreak{}r\allowbreak{}t\allowbreak{}i\allowbreak{}c\allowbreak{}\_\allowbreak{}t\allowbreak{}r\allowbreak{}a\allowbreak{}n\allowbreak{}s\allowbreak{}p\allowbreak{}o\allowbreak{}r\allowbreak{}t\allowbreak{}s\allowbreak{}.\allowbreak{}j\allowbreak{}s\allowbreak{}o\allowbreak{}n\allowbreak{}}
return these representative polynomials to any original four-face
cluster. The source-prefactor is the actual integer occurrence of that
differentiated face; there is no extra factorial or loss of its original
union label.

These are the signed finite coefficients through degree three in the
derivative of the vacuum source. No new bound on the norm of an infinite
linearized inverse, uniform gap or continuum limit is inferred from this
finite coefficient data.
\end{document}
